% Copyright (c) 2026 Scott Armstrong and Vlad Vicol.
% Released under Apache 2.0 license; see ../LICENSE.
\documentclass[11pt,reqno]{amsart}

\usepackage{amsmath,amssymb,amsthm}
\usepackage[T1]{fontenc}
\usepackage[utf8]{inputenc}
\usepackage{mathtools}
\usepackage{xcolor}
\usepackage[margin=1.15in]{geometry}
\usepackage[colorlinks=true,linkcolor=blue,citecolor=blue]{hyperref}

% Copyright (c) 2026 Scott Armstrong and Vlad Vicol.
% Released under Apache 2.0 license; see ../LICENSE.
%% macros.tex --- notation for ckn.tex.
%% Every macro here is used at least twice in ckn.tex.

%% --- Theorem environments -------------------------------------------------
\theoremstyle{plain}
\newtheorem{theorem}{Theorem}[section]
\newtheorem{proposition}[theorem]{Proposition}
\newtheorem{lemma}[theorem]{Lemma}
\newtheorem{corollary}[theorem]{Corollary}
\newtheorem{externalinput}[theorem]{External input}

\theoremstyle{definition}
\newtheorem{definition}[theorem]{Definition}
\newtheorem{remark}[theorem]{Remark}
\newtheorem{convention}[theorem]{Convention}

%% --- Spaces and operators -------------------------------------------------
\newcommand{\R}{\mathbb{R}}
\newcommand{\N}{\mathbb{N}}
\newcommand{\Z}{\mathbb{Z}}
\newcommand{\Q}{\mathbb{Q}}
\newcommand{\dv}{\operatorname{div}}
\newcommand{\curl}{\operatorname{curl}}
\newcommand{\supp}{\operatorname{supp}}
\newcommand{\dist}{\operatorname{dist}}
\newcommand{\diam}{\operatorname{diam}}

%% --- Cylinders, balls, averages -------------------------------------------
\newcommand{\Ball}[2]{B_{#1}(#2)}
\newcommand{\Cyl}[2]{Q_{#1}(#2)}
\newcommand{\spavg}[2]{\bigl[#1\bigr]_{#2}}          % spatial average, a function of t
\newcommand{\stavg}[2]{\bigl(#1\bigr)_{#2}}          % space-time average, a number

%% --- Norms ----------------------------------------------------------------
\newcommand{\abs}[1]{\left\lvert #1 \right\rvert}
\newcommand{\norm}[1]{\left\lVert #1 \right\rVert}
\newcommand{\Lp}[2]{L^{#1}(#2)}
\newcommand{\LtLx}[3]{L^{#1}\bigl(#3;L^{#2}\bigr)}

%% --- Fractions in exponents ----------------------------------------------
\newcommand{\nf}[2]{\tfrac{#1}{#2}}

%% --- Parabolic Hausdorff measure -----------------------------------------
\newcommand{\Phaus}{\mathcal{P}}
\newcommand{\dpar}{d_{\mathrm{par}}}

%% --- Suitable weak solution class ----------------------------------------
\newcommand{\SWS}{\mathrm{SWS}}

%% --- Editorial ------------------------------------------------------------


\title[The Caffarelli-Kohn-Nirenberg theorem]
      {The Caffarelli-Kohn-Nirenberg partial regularity theorem:\\
       a self-contained proof written for formalization}

\author{Scott Armstrong}
\address{CNRS and Laboratoire Jacques-Louis Lions, Sorbonne Université;
Courant Institute School of Mathematics, Computing, and Data Science, New York University}
\thanks{Scott Armstrong was supported by the European Research Council (ERC)
under the European Union's Horizon Europe research and innovation programme,
grant agreement No. 101200828.}
\author{Vlad Vicol}
\address{Courant Institute School of Mathematics, Computing, and Data Science, New York University,
New York, NY 10012}
\thanks{Vlad Vicol was partially supported by the Collaborative NSF grant
DMS-2307681 and by a Simons Investigator Award.}
\date{\today}
\thanks{Copyright \textcopyright\ 2026 Scott Armstrong and Vlad Vicol.
This manuscript is distributed under the Apache License, Version 2.0.}

\begin{document}

\begin{abstract}
We give a complete and self-contained proof of the Caffarelli-Kohn-Nirenberg
partial regularity theorem for suitable weak solutions of the three-dimensional
incompressible Navier-Stokes equations with a force $f\in L^q$, $q>5/2$: the
singular set has vanishing one-dimensional parabolic Hausdorff measure. The
exposition is written so that each statement can be transcribed into a single
formal declaration. Every constant is named and its dependence is displayed,
every cut-off function is written out, every passage to a limit is justified,
and the external analytic inputs are isolated in an appendix with exact
statements.
\end{abstract}

\maketitle

\setcounter{tocdepth}{1}
\tableofcontents

\section*{Introduction}
\addcontentsline{toc}{section}{Introduction}

Leray's weak solutions of the incompressible Navier-Stokes equations in three
space dimensions are global in time but are not known to be regular or unique.
Partial regularity theory replaces the (open) question of full regularity by a
quantitative bound on the size of the hypothetical singular set. Scheffer
\cite{Scheffer1977} initiated this programme, and Caffarelli, Kohn and Nirenberg
\cite{CKN1982} obtained the sharpest result known to date: for a
\emph{suitable} weak solution (one that additionally satisfies a localized
form of the energy inequality) the set $S$ of space-time points at which the
velocity fails to be locally bounded satisfies $\Phaus^{1}(S)=0$, where
$\Phaus^{1}$ is the one-dimensional Hausdorff measure of the parabolic metric.
Their hypothesis on the force was $f\in L^{q}$ with $q>5/2$.

The proof in \cite{CKN1982} rests on a two-step scheme that has since become
standard. The first step is an $\varepsilon$-regularity theorem: if a suitable
weak solution has small scale-invariant $L^{3}$ velocity energy and small
scale-invariant $L^{3/2}$ pressure energy, together with a small force norm,
on one parabolic cylinder, then the velocity has a H\"older representative
on the half-cylinder with the same top centre. The second step
is a covering argument: at every singular point, Theorem~\ref{thm:B} gives
$\limsup_{r\downarrow0}r^{-1}\iint_{\Cyl{r}{z}}\abs{\nabla u}^{2}\ge\varepsilon_{1}^{2}$,
and the finiteness of $\iint\abs{\nabla u}^{2}$ then
forces the singular set to be $\Phaus^{1}$-null. Two later treatments simplify
the first step in different directions. Lin \cite{Lin1998} replaced the explicit
iteration of \cite{CKN1982} by a blow-up (compactness plus linear decay)
argument; Ladyzhenskaya and Seregin \cite{LadyzhenskayaSeregin1999} and
Colombo, De~Lellis and Massaccesi \cite{ColomboDeLellisMassaccesi2018} completed
and corrected that argument. Kukavica \cite{Kukavica2008} instead kept the
direct iteration but organized it around a single monotone quantity. His
regularity criterion uses local $L^{5}$ integrability and applies to
divergence-free forces in $L^{2}$, with an extension to $L^{q}$, $q>5/3$,
stated in his Theorem 2. This is a weaker notion of regularity than the
H\"older conclusion considered here.
A direct proof of the theorem, with a smallness criterion involving the
vorticity component perpendicular to the velocity, is given by Wolf~\cite{Wolf2008}.

The present notes are written with a single purpose: to serve as the
mathematical source of a machine-checked formalization. Accordingly,
every constant is given a name ($C_{1},C_{2},\dots$) and is never
reused; when a constant is absolute (that is, independent of the solution, of
the cylinder, and of every parameter in the statement) we say so explicitly.
The estimates used at later steps are stated explicitly and cross-referenced.
%
%
%
%
%

For background on suitable weak solutions, partial regularity and the
harmonic-analysis tools used below we refer to
\cite{Adams1975,CKN1982,ColomboDeLellisMassaccesi2018,DiestelUhl1977,Giaquinta1983,GilbargTrudinger2001,Kukavica2008,LadyzhenskayaSeregin1999,LadyzhenskayaSolonnikovUraltseva1968,LemarieRieusset2002,LemarieRieusset2016,Lin1998,Mattila1995,OLeary2003,RobinsonRodrigoSadowski2016,Scheffer1977,Stein1970,Stein1993,Temam1977,Tsai2018,Vasseur2007}.

\subsection*{Organization of the manuscript}
Section~\ref{sec:setting} fixes the notation, defines suitable weak solutions
with a force, and proves the two facts about the time variable that the
definition does not literally contain: that $\partial_{t}u$ is a
$L^{3/2}$-in-time distribution of a definite order (Lemma~\ref{lem:dtu-dual}),
and that consequently $u$ has a representative that is weakly continuous in time
with values in $L^{2}$ (Lemma~\ref{lem:weak-time-continuity}), as an optional strengthening. The energy estimates below use only the
almost-everywhere-in-time local energy inequality
(Lemma~\ref{lem:lei-pointwise}), obtained directly by time differentiation. Section~\ref{sec:theorems} states the three
theorems. Section~\ref{sec:scaling} records the exact scaling and translation
invariance and proves the interpolation inequalities.
Section~\ref{sec:pressure} decomposes the pressure and proves the pressure decay
lemma. Section~\ref{sec:core} contains the
$\varepsilon$-regularity core, in the direct-iteration form of
\cite{Kukavica2008} and \cite[\S13.8--13.9]{LemarieRieusset2016}, and proves
Theorems~\ref{thm:A} and \ref{thm:B}. Section~\ref{sec:covering} derives
Theorem~\ref{thm:C} from Theorem~\ref{thm:B} by a Vitali covering argument.
Appendix~\ref{sec:external} lists the external analytic inputs.

\section{Setting and notation}\label{sec:setting}

\subsection{Ambient conventions}\label{ss:ambient}

Throughout, space is $\R^{3}$ with the Euclidean norm $\abs{\cdot}$, time is
$\R$, and a space-time point is written $z=(x,t)\in\R^{3}\times\R$. The
viscosity is $1$. Latin indices $i,j,k$ run over $\{1,2,3\}$ and repeated
indices are summed. We write $\partial_{i}=\partial/\partial x_{i}$,
$\nabla=(\partial_{1},\partial_{2},\partial_{3})$,
$\Delta=\partial_{i}\partial_{i}$ (the spatial Laplacian), and
$\partial_{t}=\partial/\partial t$. For a vector field $u$ we write
$\dv u=\partial_{i}u_{i}$ and $(u\otimes u)_{ij}=u_{i}u_{j}$. Lebesgue measure
on $\R^{3}$ and on $\R^{3}\times\R$ is denoted $dx$ and $dz=dx\,dt$; the measure
of a measurable set $E$ is $\abs{E}$.

\begin{convention}\label{conv:constants}
Constants are numbered $C_{1},C_{2},\dots$ in order of first appearance and are
never reused with a different value. A constant is called \emph{absolute} if it
depends on nothing at all; since the space dimension is fixed at $3$ and the
viscosity at $1$, dimensional constants are absolute here. Any other dependence
is displayed in the notation, as in $C_{17}(q)$.
\end{convention}

\begin{convention}[Representatives and finite local quantities]
Spatial balls and vector norms in statements are Euclidean. A displayed
gradient of a Sobolev function is its weak gradient, determined almost
everywhere. Smooth compactly supported tests mean $C^\infty$ tests, not
real-analytic tests. Integrals of nonnegative functions may first be read in
$[0,\infty]$; real-valued scale quantities are used only on cylinders whose
closures are compactly contained in the solution domain. Local finiteness
then follows from (S1) and the local force assumption. No infinite integral
is assigned a finite numerical value.

Distributional identities are tested against every compactly supported
smooth test. A slice identity initially obtained for each test outside its
own null set is promoted, when needed, to one common null set by a countable
dense test family and continuity of the pairings. Convolutions of nonsmooth
data, and identities between their measurable representatives, are understood
almost everywhere. Pointwise regularity estimates concern the continuous
representative subsequently constructed, not arbitrary values on a null set.
\end{convention}

\subsection{Balls, cylinders, and the parabolic metric}\label{ss:cylinders}

For $x\in\R^{3}$ and $r>0$ put $\Ball{r}{x}=\{y\in\R^{3}:\abs{y-x}<r\}$, so that
$\abs{\Ball{r}{x}}=\tfrac{4}{3}\pi r^{3}$. For $z=(x,t)$ the
\emph{parabolic cylinder of radius $r$ with top centre $z$} is
\begin{equation}\label{eq:cylinder}
  \Cyl{r}{z} \;=\; \Ball{r}{x}\times(t-r^{2},\,t],
  \qquad\text{so}\qquad
  \abs{\Cyl{r}{z}}=\tfrac{4}{3}\pi r^{5}.
\end{equation}
We abbreviate $B_{r}=\Ball{r}{0}$ and $Q_{r}=\Cyl{r}{(0,0)}$. The interval
$(t-r^{2},t]$ is open at the bottom and closed at the top so that $z\in\Cyl{r}{z}$; since the top
face is a null set, every integral below is unaffected by this choice.

The \emph{parabolic metric} on $\R^{3}\times\R$ is
\begin{equation}\label{eq:parabolic-metric}
  \dpar(z,z') \;=\; \max\bigl\{\abs{x-x'},\,\abs{t-t'}^{1/2}\bigr\},
  \qquad z=(x,t),\ z'=(x',t').
\end{equation}

\begin{lemma}\label{lem:parabolic-metric}
$\dpar$ is a metric on $\R^{3}\times\R$, and the open $\dpar$-ball of radius $r$
about $z=(x,t)$ is
\begin{equation}\label{eq:parabolic-ball}
  \mathfrak{B}_{r}(z)
  \;=\;\{z':\dpar(z',z)<r\}
  \;=\;\Ball{r}{x}\times(t-r^{2},t+r^{2}).
\end{equation}
Moreover $\Cyl{r}{z}\subset\mathfrak{B}_{r}(z)$ and
$\mathfrak{B}_{r}(z)\subset \Cyl{2r}{(x,t+r^{2})}$ for every $z$ and every
$r>0$.
\end{lemma}

\begin{proof}
Symmetry and the vanishing of $\dpar(z,z')$ exactly when $z=z'$ are immediate
from the definition. For the triangle inequality, let $z,z',z''$ be given. Since
$\abs{\cdot}$ is a metric on $\R^{3}$,
$\abs{x-x''}\le\abs{x-x'}+\abs{x'-x''}\le\dpar(z,z')+\dpar(z',z'')$. For the
time entry, the elementary inequality $(a+b)^{1/2}\le a^{1/2}+b^{1/2}$, valid
for $a,b\ge0$ because $a+b\le a+2a^{1/2}b^{1/2}+b$, gives
\[
  \abs{t-t''}^{1/2}\le\bigl(\abs{t-t'}+\abs{t'-t''}\bigr)^{1/2}
  \le\abs{t-t'}^{1/2}+\abs{t'-t''}^{1/2}
  \le\dpar(z,z')+\dpar(z',z'').
\]
Taking the maximum of the two bounds yields
$\dpar(z,z'')\le\dpar(z,z')+\dpar(z',z'')$.

Identity \eqref{eq:parabolic-ball} is a restatement of the definition:
$\dpar(z',z)<r$ holds if and only if $\abs{x'-x}<r$ and $\abs{t'-t}<r^{2}$. The
inclusion $\Cyl{r}{z}\subset\mathfrak{B}_{r}(z)$ holds because
$(t-r^{2},t]\subset(t-r^{2},t+r^{2})$. Finally
$\Cyl{2r}{(x,t+r^{2})}=\Ball{2r}{x}\times(t-3r^{2},t+r^{2}]$ contains
$\Ball{r}{x}\times(t-r^{2},t+r^{2})$.
\end{proof}

\subsection{Averages}\label{ss:averages}

Two distinct averaging operations occur, and they must not be confused. For a
locally integrable $g$ on a space-time set, the \emph{spatial average} over
$\Ball{r}{x}$ is the function of time
\begin{equation}\label{eq:spatial-average}
  \spavg{g}{\Ball{r}{x}}(t)
  \;=\;\frac{1}{\abs{\Ball{r}{x}}}\int_{\Ball{r}{x}}g(y,t)\,dy,
\end{equation}
defined for almost every $t$; and the \emph{space-time average} over
$\Cyl{r}{z}$ is the number
\begin{equation}\label{eq:spacetime-average}
  \stavg{g}{\Cyl{r}{z}}
  \;=\;\frac{1}{\abs{\Cyl{r}{z}}}\iint_{\Cyl{r}{z}}g(y,s)\,dy\,ds .
\end{equation}
For a vector field these are applied componentwise. We also write
$\spavg{g}{\Ball{r}{x}}$ for the function $t\mapsto\spavg{g}{\Ball{r}{x}}(t)$
when no confusion is possible.

\subsection{Function spaces}\label{ss:spaces}

Let $U\subset\R^{3}$ be open and bounded. We write $L^{s}(U)$, $1\le s\le\infty$,
for the usual Lebesgue spaces, $H^{1}(U)=W^{1,2}(U)$, and $H^{2}_{0}(U)$ for the
closure of $C_{c}^{\infty}(U)$ in the norm
$\norm{\varphi}_{H^{2}(U)}=\bigl(\sum_{\abs{a}\le2}\norm{\partial^{a}\varphi}_{L^{2}(U)}^{2}\bigr)^{1/2}$.
For vector fields we write $H^{2}_{0}(U;\R^{3})$ and we set
\begin{equation}\label{eq:Z-space}
  Z_{U}\;=\;\bigl(H^{2}_{0}(U;\R^{3})\bigr)^{*},
\end{equation}
the topological dual, with the dual norm. For a time interval $J$ and a Banach
space $X$, $L^{s}(J;X)$ is the usual Bochner space. We abbreviate
\[
  \norm{g}_{L^{s}(E)}=\Bigl(\iint_{E}\abs{g}^{s}\Bigr)^{1/s}
\]
for a space-time set $E$ and $1\le s<\infty$.

We record the two Sobolev embeddings that will be used on $H^{2}_{0}$. Because a
function in $H^{2}_{0}(U)$ extends by zero to a function in $H^{2}(\R^{3})$ of
the same norm, the constants are absolute and do not depend on $U$; see
External Input~\ref{ext:sobolev-ball}.

\begin{lemma}\label{lem:H2-embeddings}
There is an absolute constant $C_{1}$ such that for every bounded open
$U\subset\R^{3}$ and every $\varphi\in H^{2}_{0}(U;\R^{3})$,
\begin{equation}\label{eq:H2-emb}
  \norm{\varphi}_{L^{\infty}(U)}\le C_{1}\norm{\varphi}_{H^{2}(U)},
  \qquad
  \norm{\nabla\varphi}_{L^{6}(U)}\le C_{1}\norm{\varphi}_{H^{2}(U)} .
\end{equation}
In particular, for every $s\in[1,\infty]$,
$\norm{\varphi}_{L^{s}(U)}\le C_{1}\abs{U}^{1/s}\norm{\varphi}_{H^{2}(U)}$,
and for every $s\in[1,6]$,
$\norm{\nabla\varphi}_{L^{s}(U)}\le C_{1}\abs{U}^{1/s-1/6}\norm{\varphi}_{H^{2}(U)}$.
\end{lemma}

\begin{proof}
Let $\tilde\varphi$ denote the extension of $\varphi$ to $\R^{3}$ by zero. Since
$\varphi\in H^{2}_{0}(U;\R^{3})$ is a limit in $H^{2}$ of elements of
$C_{c}^{\infty}(U;\R^{3})$, and the zero extension of such an element belongs to
$C_{c}^{\infty}(\R^{3};\R^{3})$ with unchanged $H^{2}$ norm, the map
$\varphi\mapsto\tilde\varphi$ is an isometry of $H^{2}_{0}(U;\R^{3})$ into
$H^{2}(\R^{3};\R^{3})$. By External Input~\ref{ext:sobolev-ball}(iii) there is
an absolute $C$ with
$\norm{\tilde\varphi}_{L^{\infty}(\R^{3})}\le C\norm{\tilde\varphi}_{H^{2}(\R^{3})}$,
and by External Input~\ref{ext:sobolev-ball}(i) applied to each
$\partial_{i}\tilde\varphi\in H^{1}(\R^{3})$ there is an absolute $C'$ with
$\norm{\nabla\tilde\varphi}_{L^{6}(\R^{3})}\le C'\norm{\nabla\tilde\varphi}_{H^{1}(\R^{3})}
 \le C'\norm{\tilde\varphi}_{H^{2}(\R^{3})}$. Take $C_{1}=\max\{C,C'\}$ and
restrict to $U$. The two consequences follow from H\"older's inequality on $U$:
$\norm{\varphi}_{L^{s}(U)}\le\abs{U}^{1/s}\norm{\varphi}_{L^{\infty}(U)}$ and,
for $s\le6$,
$\norm{\nabla\varphi}_{L^{s}(U)}\le\abs{U}^{1/s-1/6}\norm{\nabla\varphi}_{L^{6}(U)}$.
\end{proof}

\subsection{Suitable weak solutions with a force}\label{ss:sws}

Throughout the paper $\Omega\subset\R^{3}$ is open, $I\subset\R$ is an open
interval, and we write
\[
  \mathcal{O}=\Omega\times I\subset\R^{3}\times\R .
\]
Only local statements are made: no initial condition, no boundary condition, and
no global energy inequality is imposed, and the existence of solutions in the
class below is not discussed.

\begin{definition}[Suitable weak solution]\label{def:sws}
Let $q>5/2$ and let $f\in L^{q}_{\mathrm{loc}}(\mathcal{O};\R^{3})$. A pair
$(u,p)$ is a \emph{suitable weak solution of the Navier-Stokes equations on
$\mathcal{O}$ with force $f$} if the following four conditions hold.
\begin{enumerate}
\item[(S1)] \emph{(Regularity class.)} For every open $\Omega'$ with
  $\overline{\Omega'}\subset\Omega$ compact and every interval $J$ with
  $\overline{J}\subset I$ compact,
  \[
    u\in L^{\infty}\bigl(J;L^{2}(\Omega';\R^{3})\bigr)
        \cap L^{2}\bigl(J;H^{1}(\Omega';\R^{3})\bigr),
    \qquad
    p\in L^{3/2}(\Omega'\times J).
  \]
\item[(S2)] \emph{(Incompressibility.)} $\dv u=0$ in
  $\mathcal{D}'(\mathcal{O})$, that is,
  $\iint_{\mathcal{O}}u\cdot\nabla\psi\,dz=0$ for every
  $\psi\in C_{c}^{\infty}(\mathcal{O})$. \emph{No condition is imposed on
  $\dv f$.}
\item[(S3)] \emph{(Momentum equation.)} For every
  $\varphi\in C^{\infty}_{c}(\mathcal{O};\R^{3})$,
  \begin{equation}\label{eq:weak-momentum}
    \iint_{\mathcal{O}}
      \Bigl[\,-\,u\cdot\partial_{t}\varphi
        \;-\;u_{i}u_{j}\,\partial_{j}\varphi_{i}
        \;+\;\nabla u:\nabla\varphi
        \;-\;p\,\dv\varphi
        \;-\;f\cdot\varphi\,\Bigr]dz \;=\;0 .
  \end{equation}
  Here $\nabla u:\nabla\varphi=\partial_{j}u_{i}\,\partial_{j}\varphi_{i}$.
  Equation \eqref{eq:weak-momentum} is the distributional form of
  \begin{equation}\label{eq:NS}
    \partial_{t}u+\dv(u\otimes u)-\Delta u+\nabla p=f,
    \qquad \dv u=0 .
  \end{equation}
\item[(S4)] \emph{(Local energy inequality.)} For every
  $\phi\in C_{c}^{\infty}(\mathcal{O})$ with $\phi\ge0$,
  \begin{equation}\label{eq:lei}
    2\iint_{\mathcal{O}}\abs{\nabla u}^{2}\phi\,dz
    \;\le\;
    \iint_{\mathcal{O}}\Bigl[\,\abs{u}^{2}\bigl(\partial_{t}\phi+\Delta\phi\bigr)
      +\bigl(\abs{u}^{2}+2p\bigr)\,u\cdot\nabla\phi
      +2(f\cdot u)\,\phi\,\Bigr]dz .
  \end{equation}
\end{enumerate}
\end{definition}

\begin{remark}[The gradient is a datum]\label{rem:gradient-datum}
Condition (S1) asserts $u(\cdot,t)\in H^{1}(\Omega')$ for a.e.\ $t$, so the weak
spatial gradient $\nabla u$ exists and is uniquely determined almost everywhere
on $\mathcal{O}$. In a formalization it is convenient to carry $\nabla u$ as an
explicit datum $\mathbf{G}$ satisfying
$\iint\mathbf{G}_{ij}\chi=-\iint u_{i}\partial_{j}\chi$ for all
$\chi\in C_{c}^{\infty}(\mathcal{O})$, rather than as a derived object; the
a.e.\ uniqueness of $\mathbf{G}$ makes the two formulations equivalent. With
that convention, all the quantities of \S\ref{ss:quantities} (in particular
$\beta$, and hence $\theta$ of \eqref{eq:theta}) are functions of the tuple
$(u,\mathbf{G},p,f)$ and of $(z,r)$, and every statement below that mentions
$\nabla u$ is to be read with $\mathbf{G}$ in its place. No mathematical content
changes.
\end{remark}

The sign conventions in \eqref{eq:lei} are those of \cite{CKN1982},
\cite{Lin1998} and \cite{Kukavica2008}: every term on the right carries the
sign it acquires when $\abs{u}^{2}\phi$ is formally differentiated in time using
\eqref{eq:NS} and integrated by parts, and the dissipation appears on the left
with the factor $2$. Formally, \eqref{eq:lei} is what one gets by multiplying
\eqref{eq:NS} by $2u\phi$ and discarding the (nonnegative) defect measure.

\begin{remark}[The force class]\label{rem:force}
Following \cite{CKN1982} we assume $f\in L^{q}$ with $q>5/2$; unlike
\cite{Kukavica2008} and \cite[Definition 13.4]{LemarieRieusset2016} we do
\emph{not} assume $\dv f=0$. Two comments.

(i) The exponent $q$ enters only through the scale-invariant force functional
$\lambda$ of \eqref{eq:alpha-beta}, for which
$\lambda(z,r)\le\norm{f}_{L^{q}}\,r^{\sigma}$ with
\begin{equation}\label{eq:sigma}
  \sigma\;=\;3-\tfrac5q\;>\;1 ,
\end{equation}
so that $\lambda(z,r)\to0$ as $r\to0$ at a rate strictly better than linear.
The strict inequality $\sigma>1$ is used in Lemma~\ref{lem:pk-bounds}(d)--(e)
and, through $\sigma>\varepsilon$, in Proposition~\ref{prop:iteration}; the
threshold $q>5/2$ reappears in Theorem~\ref{thm:endgame} as the condition
$\theta_{0}>5/2$ for the force slot.
Kukavica \cite[Theorem 2]{Kukavica2008} shows that the decay statements of
Sections~\ref{sec:scaling} and \ref{sec:pressure} already hold for $q>5/3$,
where $\sigma>0$ but possibly $\sigma\le1$. The endgame here uses
$\sigma>1$. Remark~\ref{rem:shear} shows that an $L^{2}$ force hypothesis
does not suffice for this H\"older notion of regularity; the example does
not establish sharpness of the exponent $5/2$.
%
%

(ii) We do \emph{not} assume $\dv f=0$, in contrast with
\cite{Kukavica2008}. The hypothesis would matter at exactly two places: it would
make the sum $p_{7}+p_{8}$ vanish of the pressure decomposition of
Proposition~\ref{prop:pressure-decomposition}, and hence the term
$C_{15}\kappa^{1/2}\lambda(z,\rho)^{1/2}$ in the pressure decay estimate
\eqref{eq:pressure-decay} and the term
$C_{32}(q)(r/\rho)^{3/2}\lambda^{3/2}$ in \eqref{eq:lin35-force}. Both terms are estimated in
Lemma~\ref{lem:pk-bounds}(d)--(e) and are carried through
Section~\ref{sec:core}; in the scale iteration they contribute the term
$C_{28}\kappa^{-3}\lambda(z,\rho)$ of \eqref{eq:theta-decay-1}, which is
harmless because $\lambda(z,\rho)\le\norm{f}_{L^{q}}\rho^{\sigma}$ with
$\sigma>1>\varepsilon$. In the endgame the force contribution to the pressure gradient is
estimated in Lemma~\ref{lem:pressure-gradient-morrey}: the near force
term is a completed second Riesz transform of the localized force, and
the annular term has a fixed-scale $L^{3/2}$ time majorant. Neither
estimate requires $\dv f=0$.
\end{remark}

\begin{remark}[Why an $L^{2}$ force does not suffice: a shear flow]\label{rem:shear}
A planar shear flow gives a counterexample at force exponent $2$. Write
$x=(x',x_{3})$ with $x'=(x_{1},x_{2})$, and let
$w\in W^{2,1}_{2}(\R^{2}\times(-1,1))$ be a solution of a two-dimensional heat
equation that is unbounded near the origin; such $w$ exist because
$W^{2,1}_{2}$ is the borderline Sobolev case in parabolic dimension $4$, for
instance a compactly supported cutoff of $(\log(1/\varrho))^{1/4}$ near
the origin, with $\varrho=(\abs{x'}^{2}+\abs{t})^{1/2}$. Its second
spatial derivatives and first time derivative are bounded by
$C\varrho^{-2}(\log(1/\varrho))^{-3/4}$ near the origin. Their squares are
integrable in parabolic dimension $4$, since
$\int_{0}^{1/2} r^{-1}(\log(1/r))^{-3/2}\,dr<\infty$. Put
\[
  u(x,t)=\bigl(0,0,w(x',t)\bigr),\qquad p\equiv0,\qquad
  f=\bigl(0,0,\partial_{t}w-\Delta_{x'}w\bigr).
\]
Then $\dv u=0$ and $\dv f=0$ (both fields depend only on $(x',t)$ and point in
the $x_{3}$ direction), $u\cdot\nabla u=u_{3}\partial_{3}u=0$, $f\in L^{2}_{loc}$,
$p\in L^{3/2}$, and the local energy inequality holds with equality. Since
$\nabla w\in L^{4}_{loc}$ by the parabolic Sobolev embedding in dimension $4$,
one computes $\iint_{\Cyl{r}{0}}\abs{\nabla u}^{2}=o(r^{3})$, so
$\beta(0,r)^{2}=r^{-1}\iint_{\Cyl{r}{0}}\abs{\nabla u}^{2}\to0$ and the
hypothesis \eqref{eq:thmB-hyp} of Theorem~\ref{thm:B} holds for every threshold.
Yet $u$ is unbounded at the origin, so the origin is not a regular point in the
sense of Definition~\ref{def:regular}. Hence the conclusion of
Theorem~\ref{thm:B} is \emph{false} for $f\in L^{2}$ with the H\"older (indeed
with the $L^{\infty}$) notion of regularity. Kukavica's Theorem~1 for
$f\in L^{2}$ is stated for the weaker notion ``$u\in L^{5}$ on a
neighbourhood'' \cite[p.~719]{Kukavica2008}, which this example does not
contradict.
\end{remark}

\begin{remark}[Comparison with the sources]\label{rem:definition-sources}
Definition~\ref{def:sws} is the class used by Lin \cite[\S2]{Lin1998} and by
Ladyzhenskaya--Seregin \cite[Definition 2.1]{LadyzhenskayaSeregin1999}, with a
force added and with the domain restricted to a product $\Omega\times I$. Two
features of our formulation should be noted. First, $p\in L^{3/2}_{loc}$
is an explicit hypothesis. We use the modern suitable-solution class just
cited, without claiming that its pressure hypothesis is identical to the
original definition in \cite{CKN1982} or deriving it from weaker local data.
Second, \cite{LadyzhenskayaSeregin1999} and \cite{Kukavica2008} state the
local energy inequality directly in the pointwise-in-time form
\eqref{eq:lei-pointwise} below; we take the distributional form \eqref{eq:lei}
as the definition and \emph{prove} the pointwise form in
Lemma~\ref{lem:lei-pointwise}, since the two are not literally the same
statement and the derivation uses a time cutoff and Lebesgue differentiation.
\end{remark}

\subsection{The time derivative and weak continuity in time}\label{ss:time}

We first derive a weakly continuous $L^{2}$ representative from (S3). This
is a useful additional property of suitable solutions. The almost-every-time
local energy inequality proved afterwards uses Lebesgue differentiation
directly and does not require this representative. Fix, for the whole
subsection, a ball $B=\Ball{\rho}{x_{0}}$ and a
bounded open interval $J$ with
\begin{equation}\label{eq:BJ-compact}
  \overline{B}\subset\Omega,\qquad \overline{J}\subset I .
\end{equation}

\begin{lemma}[The time derivative lies in $L^{3/2}(J;Z_{B})$]\label{lem:dtu-dual}
Let $(u,p)$ be a suitable weak solution on $\mathcal{O}$ with force
$f\in L^{q}_{loc}$, $q>5/2$, and let $B,J$ be as in \eqref{eq:BJ-compact}. Then
the distribution $\partial_{t}u$ on $B\times J$ belongs to
$L^{3/2}(J;Z_{B})$, where $Z_{B}=(H^{2}_{0}(B;\R^{3}))^{*}$ as in
\eqref{eq:Z-space}: there is a function
$G\in L^{3/2}(J;Z_{B})$ such that
\begin{equation}\label{eq:dtu-weak}
  -\int_{J}\int_{B}u\cdot\varphi\,\theta'(t)\,dx\,dt
  \;=\;\int_{J}\bigl\langle G(t),\varphi\bigr\rangle\,\theta(t)\,dt
  \qquad
  \text{for all }\varphi\in H^{2}_{0}(B;\R^{3}),\ \theta\in C_{c}^{\infty}(J),
\end{equation}
and, with $C_{1}$ the constant of Lemma~\ref{lem:H2-embeddings},
\begin{multline}\label{eq:dtu-bound}
  \norm{G}_{L^{3/2}(J;Z_{B})}
  \le\sqrt{3}C_{1}\Bigl[
      \abs{J}^{1/6}\abs{B}^{1/3}\norm{\nabla u}_{L^{2}(B\times J)}
      +\abs{J}^{1/15}\abs{B}^{7/30}\norm{u}_{L^{10/3}(B\times J)}^{2}\\
      +\abs{B}^{1/6}\norm{p}_{L^{3/2}(B\times J)}
      +\abs{J}^{\frac23-\frac1q}\abs{B}^{1-\frac1q}\,\norm{f}_{L^{q}(B\times J)}
     \Bigr].
\end{multline}
In particular, if $\abs{B}\le1$ and $\abs{J}\le1$ then
\begin{equation}\label{eq:dtu-bound-simple}
  \norm{G}_{L^{3/2}(J;Z_{B})}
  \le C_{2}\Bigl[\norm{\nabla u}_{L^{2}(B\times J)}
    +\norm{u}_{L^{10/3}(B\times J)}^{2}
    +\norm{p}_{L^{3/2}(B\times J)}
    +\norm{f}_{L^{q}(B\times J)}\Bigr],
\end{equation}
where $C_{2}=\sqrt{3}C_{1}$ is absolute.
\end{lemma}

\begin{proof}
\emph{Step 1: the candidate for $G$.} For $t\in J$ define a linear functional on
$H^{2}_{0}(B;\R^{3})$ by
\begin{equation}\label{eq:G-def}
  \bigl\langle G(t),\varphi\bigr\rangle
  \;=\;\int_{B}\Bigl[\,u_{i}(x,t)u_{j}(x,t)\,\partial_{j}\varphi_{i}(x)
      -\nabla u(x,t):\nabla\varphi(x)
      +p(x,t)\,\dv\varphi(x)
      +f(x,t)\cdot\varphi(x)\,\Bigr]dx .
\end{equation}
By (S1) and $f\in L^{q}_{loc}$ each of the four integrands is in $L^{1}(B)$ for
almost every $t\in J$, so $\langle G(t),\varphi\rangle$ is defined for a.e.\
$t$, and it is measurable in $t$ for each fixed $\varphi$. Since
$H^{2}_{0}(B;\R^{3})$ is a separable Hilbert space, so is its dual $Z_{B}$;
hence, by Pettis' measurability theorem
(External Input~\ref{ext:pettis}), this weak measurability implies that
$t\mapsto G(t)$ is strongly (Bochner) measurable, which is what membership in
the Bochner space $L^{3/2}(J;Z_{B})$ requires.

\emph{Step 2: the four bounds.} Fix $t$ for which all four integrands are
integrable, and let $\varphi\in H^{2}_{0}(B;\R^{3})$ with
$\norm{\varphi}_{H^{2}(B)}\le1$.

\emph{(a) Dissipation.} By Cauchy--Schwarz and
Lemma~\ref{lem:H2-embeddings} (with $s=2\le6$, so
$\abs{B}^{1/2-1/6}\le\abs{B}^{1/3}$),
\[
  \Bigl|\int_{B}\nabla u:\nabla\varphi\Bigr|
  \le\norm{\nabla u(\cdot,t)}_{L^{2}(B)}\norm{\nabla\varphi}_{L^{2}(B)}
  \le C_{1}\abs{B}^{1/3}\norm{\nabla u(\cdot,t)}_{L^{2}(B)} .
\]

\emph{(b) Transport.} Since $H^{2}_{0}(B)\not\hookrightarrow W^{1,\infty}(B)$ in
three dimensions, the convective term must be paired against
$\nabla\varphi\in L^{6}$ rather than against $\nabla\varphi\in L^{\infty}$. By
H\"older with $\tfrac{5}{6}+\tfrac16=1$ and Lemma~\ref{lem:H2-embeddings},
\[
  \Bigl|\int_{B}u_{i}u_{j}\,\partial_{j}\varphi_{i}\Bigr|
  \le\norm{u(\cdot,t)}_{L^{12/5}(B)}^{2}\,
     \norm{\nabla\varphi}_{L^{6}(B)}
  \le C_{1}\norm{u(\cdot,t)}_{L^{12/5}(B)}^{2} ,
\]
where we used $\norm{u_{i}u_{j}}_{L^{6/5}(B)}\le\norm{u}_{L^{12/5}(B)}^{2}$.

\emph{(c) Pressure.} By H\"older with $\tfrac{2}{3}+\tfrac13=1$ and
Lemma~\ref{lem:H2-embeddings} with $s=3\le6$, using
$\abs{\dv\varphi}\le\sqrt{3}\abs{\nabla\varphi}$,
\[
  \Bigl|\int_{B}p\,\dv\varphi\Bigr|
  \le\sqrt{3}\norm{p(\cdot,t)}_{L^{3/2}(B)}\,\norm{\nabla\varphi}_{L^{3}(B)}
  \le\sqrt{3}C_{1}\abs{B}^{1/6}\norm{p(\cdot,t)}_{L^{3/2}(B)} .
\]

\emph{(d) Force.} By H\"older with $\tfrac1q+\tfrac1{q'}=1$ and
Lemma~\ref{lem:H2-embeddings},
\[
  \Bigl|\int_{B}f\cdot\varphi\Bigr|
  \le\norm{f(\cdot,t)}_{L^{q}(B)}\norm{\varphi}_{L^{q'}(B)}
  \le C_{1}\abs{B}^{1/q'}\norm{f(\cdot,t)}_{L^{q}(B)} .
\]
Adding,
\[
\begin{aligned}
&\norm{G(t)}_{Z_{B}}\le\sqrt{3}C_{1}\bigl[\abs{B}^{1/3}\norm{\nabla u(\cdot,t)}_{L^{2}(B)}
\\ &\quad+\norm{u(\cdot,t)}_{L^{12/5}(B)}^{2}+\abs{B}^{1/6}\norm{p(\cdot,t)}_{L^{3/2}(B)}
\\ &\quad+\abs{B}^{1/q'}\norm{f(\cdot,t)}_{L^{q}(B)}\bigr].\end{aligned}\]

\emph{Step 3: time integrability of each of the four terms.} We show that each
of the four functions of $t$ on the right lies in $L^{3/2}(J)$.

\emph{(a$'$)} By (S1), $t\mapsto\norm{\nabla u(\cdot,t)}_{L^{2}(B)}$ belongs
to $L^{2}(J)$, and therefore to $L^{3/2}(J)$, with
$\norm{\cdot}_{L^{3/2}(J)}\le\abs{J}^{1/6}\norm{\cdot}_{L^{2}(J)}$ by H\"older,
since $\tfrac1{3/2}-\tfrac12=\tfrac16$.

\emph{(b)} Interpolating $L^{12/5}$ between $L^{2}$ and $L^{6}$ with
$\tfrac{5}{12}=\tfrac{3/4}{2}+\tfrac{1/4}{6}$ gives
\[
\norm{u}_{L^{12/5}(B)}\le\norm{u}_{L^{2}(B)}^{3/4}\norm{u}_{L^{6}(B)}^{1/4}
,
\]
so
\[
\norm{u(\cdot,t)}_{L^{12/5}(B)}^{2}
 \le\norm{u(\cdot,t)}_{L^{2}(B)}^{3/2}\norm{u(\cdot,t)}_{L^{6}(B)}^{1/2}
.
\]
By (S1) and External Input~\ref{ext:sobolev-ball}(ii),
$t\mapsto\norm{u(\cdot,t)}_{L^{6}(B)}$ lies in $L^{2}(J)$ and
$t\mapsto\norm{u(\cdot,t)}_{L^{2}(B)}$ lies in $L^{\infty}(J)$; hence the
product lies in $L^{4}(J)\subset L^{3/2}(J)$. Equivalently, and this is the form
recorded in \eqref{eq:dtu-bound}, the interpolation
$L^{\infty}(J;L^{2})\cap L^{2}(J;L^{6})\subset L^{10/3}(B\times J)$
(External Input~\ref{ext:sobolev-ball}(iv)) gives
\[
\norm{u}^{2}_{L^{12/5}(B)}\le\norm{u}^{2}_{L^{10/3}(B)}\abs{B}^{2(5/12-3/10)}
\]

and $t\mapsto\norm{u(\cdot,t)}^{2}_{L^{10/3}(B)}\in L^{5/3}(J)\subset L^{3/2}(J)$.

\emph{(c)} $t\mapsto\norm{p(\cdot,t)}_{L^{3/2}(B)}$ lies in $L^{3/2}(J)$ by
(S1), with $L^{3/2}(J)$-norm equal to $\norm{p}_{L^{3/2}(B\times J)}$. This is
the term that fixes the exponent $3/2$ in the statement: none of the other three
terms forces an exponent below $3/2$, and this one is exactly $3/2$.

\emph{(d)} $t\mapsto\norm{f(\cdot,t)}_{L^{q}(B)}$ lies in $L^{q}(J)$ with norm
$\norm{f}_{L^{q}(B\times J)}$, and $q>5/2>3/2$, so it lies in $L^{3/2}(J)$ with
$\norm{\cdot}_{L^{3/2}(J)}\le\abs{J}^{2/3-1/q}\norm{f}_{L^{q}(B\times J)}$.

Collecting the four contributions gives \eqref{eq:dtu-bound}: the first is
$C_{1}\abs{B}^{1/3}\abs{J}^{1/6}\norm{\nabla u}_{L^{2}(B\times J)}$ by (a) and
(a$'$); the second is
$C_{1}\abs{B}^{7/30}\abs{J}^{1/15}\norm{u}^{2}_{L^{10/3}(B\times J)}$, since
\[
\norm{u(\cdot,t)}^{2}_{L^{12/5}(B)}\le\abs{B}^{2(\frac5{12}-\frac3{10})}
 \norm{u(\cdot,t)}^{2}_{L^{10/3}(B)}
\]
 with $2(\tfrac5{12}-\tfrac3{10})=\tfrac7{30}$
and
\[
\norm{\,\cdot\,}_{L^{3/2}(J)}\le\abs{J}^{\frac23-\frac35}\norm{\,\cdot\,}_{L^{5/3}(J)}
\]

with $\tfrac23-\tfrac35=\tfrac1{15}$; the third is
\[
\sqrt{3}C_{1}\abs{B}^{1/6}\norm{p}_{L^{3/2}(B\times J)}
\]
 by (c); and the fourth is
\[
C_{1}\abs{B}^{1/q'}\abs{J}^{\frac23-\frac1q}\norm{f}_{L^{q}(B\times J)}
\]
 by (d),
with $1/q'=1-1/q$. In particular $G\in L^{3/2}(J;Z_{B})$, and
\eqref{eq:dtu-bound-simple} follows by discarding the factors
$\abs{B}^{a},\abs{J}^{b}\le1$.

\emph{Step 4: $G$ represents $\partial_{t}u$.} Let
$\varphi\in C_{c}^{\infty}(B;\R^{3})$ and $\theta\in C_{c}^{\infty}(J)$, and
apply (S3) with test field $(x,t)\mapsto\theta(t)\varphi(x)$, which belongs to
$C^{\infty}_{c}(\mathcal{O};\R^{3})$ by \eqref{eq:BJ-compact}. Since
$\partial_{t}\bigl(\theta\varphi\bigr)=\theta'\varphi$ and all spatial
derivatives fall on $\varphi$, \eqref{eq:weak-momentum} reads
\[
  -\int_{J}\int_{B}u\cdot\varphi\,\theta'\,dx\,dt
  =\int_{J}\theta(t)\int_{B}\Bigl[u_{i}u_{j}\partial_{j}\varphi_{i}
     -\nabla u:\nabla\varphi+p\,\dv\varphi+f\cdot\varphi\Bigr]dx\,dt
  =\int_{J}\langle G(t),\varphi\rangle\,\theta(t)\,dt ,
\]
which is \eqref{eq:dtu-weak} for $\varphi\in C_{c}^{\infty}(B;\R^{3})$. Both
sides are continuous in $\varphi$ with respect to the $H^{2}(B)$ norm (the
left side because
\[
\begin{aligned}
&\abs{\int_{J}\int_{B}u\cdot\varphi\,\theta'}
\\ &\le\norm{\theta'}_{L^{\infty}}\norm{u}_{L^{1}(B\times J)}\norm{\varphi}_{L^{\infty}(B)}
\\ &\le C_{1}\norm{\theta'}_{L^{\infty}}\norm{u}_{L^{1}(B\times J)}\norm{\varphi}_{H^{2}(B)}
\end{aligned}
,
\]
the right side by Step 2) and $C_{c}^{\infty}(B;\R^{3})$ is dense in
$H^{2}_{0}(B;\R^{3})$ by definition. Hence \eqref{eq:dtu-weak} holds for all
$\varphi\in H^{2}_{0}(B;\R^{3})$.
\end{proof}

\begin{lemma}[Weak continuity in time]\label{lem:weak-time-continuity}
Let $(u,p)$ be a suitable weak solution on $\mathcal{O}$ with force
$f\in L^{q}_{loc}$, $q>5/2$, and let $B,J$ be as in \eqref{eq:BJ-compact}. Then
there is a set $N\subset J$ of Lebesgue measure zero and a map
$\bar u:\overline{J}\to L^{2}(B;\R^{3})$ such that
\begin{enumerate}
\item[(i)] $\bar u(t)=u(\cdot,t)$ in $L^{2}(B;\R^{3})$ for every $t\in J\setminus N$;
\item[(ii)] $\sup_{t\in\overline{J}}\norm{\bar u(t)}_{L^{2}(B)}
  \le\norm{u}_{L^{\infty}(J;L^{2}(B))}=:M_{B,J}$;
\item[(iii)] for every $w\in L^{2}(B;\R^{3})$ the function
  $t\mapsto\int_{B}\bar u(t)\cdot w\,dx$ is continuous on $\overline{J}$.
\end{enumerate}
We say that $u$ belongs to $C(\overline{J};L^{2}_{w}(B))$ and we always work
with the representative $\bar u$.
\end{lemma}

\begin{proof}
\emph{Step 1: a continuous $Z_{B}$-valued representative.} By
Lemma~\ref{lem:dtu-dual}, $G\in L^{3/2}(J;Z_{B})\subset L^{1}(J;Z_{B})$, so
$t\mapsto\int_{t_{1}}^{t}G(s)\,ds$ is a continuous $Z_{B}$-valued function on
$\overline{J}$ for any fixed $t_{1}\in J$. Set
$v(t)=\int_{t_{1}}^{t}G(s)\,ds\in Z_{B}$. Regarding $u(\cdot,t)\in L^{2}(B)$ as
an element of $Z_{B}$ through $\langle u(\cdot,t),\varphi\rangle
=\int_{B}u(\cdot,t)\cdot\varphi$ (a bounded map by
Lemma~\ref{lem:H2-embeddings}), identity \eqref{eq:dtu-weak} says that the
$Z_{B}$-valued distribution $u-v$ on $J$ has vanishing time derivative, i.e.\
$\int_{J}(u(\cdot,t)-v(t))\theta'(t)\,dt=0$ in $Z_{B}$ for every
$\theta\in C_{c}^{\infty}(J)$. By the vector-valued du Bois-Reymond lemma
(External Input~\ref{ext:dubois-reymond}) applied in the Banach space $Z_{B}$,
there is $\xi\in Z_{B}$ with $u(\cdot,t)=\xi+v(t)$ in $Z_{B}$ for a.e.\ $t\in J$;
let $N$ contain the null set where this fails and the null set where the
$L^{2}$ bound $\norm{u(\cdot,t)}_{L^{2}(B)}\le M_{B,J}$ fails. Define
$\bar u(t)=\xi+v(t)\in Z_{B}$ for all $t\in\overline{J}$; then
$\bar u:\overline{J}\to Z_{B}$ is continuous and agrees with $u(\cdot,t)$ for
$t\in J\setminus N$. This is (i), once we know $\bar u(t)\in L^{2}$.

\emph{Step 2: $\bar u(t)\in L^{2}(B)$ with the bound (ii).} Let
$t_{0}\in\overline{J}$ and choose $t_{k}\in J\setminus N$ with $t_{k}\to t_{0}$;
this is possible because $J\setminus N$ has full measure in $J$ and
$\overline{J}$ is the closure of $J$. By (S1),
$\norm{\bar u(t_{k})}_{L^{2}(B)}=\norm{u(\cdot,t_{k})}_{L^{2}(B)}\le M_{B,J}$
for every $k$, by the choice of $N$, and
$\bar u(t_{k})\to\bar u(t_{0})$ in $Z_{B}$ by Step 1. We claim that a sequence
$g_{k}\in L^{2}(B;\R^{3})$ with $\norm{g_{k}}_{L^{2}}\le M$ that converges in
$Z_{B}$ to some $g\in Z_{B}$ satisfies $g\in L^{2}(B;\R^{3})$,
$\norm{g}_{L^{2}}\le M$, and $g_{k}\rightharpoonup g$ weakly in $L^{2}$. Indeed,
for $\varphi\in C^{\infty}_{c}(B;\R^{3})$,
$\abs{\langle g,\varphi\rangle}=\lim_{k}\abs{\int_{B}g_{k}\cdot\varphi}
\le M\norm{\varphi}_{L^{2}(B)}$, so $g$ extends to a bounded functional on the
$L^{2}$-closure of $C_{c}^{\infty}(B;\R^{3})$, namely $L^{2}(B;\R^{3})$, of norm
at most $M$; by the Riesz representation theorem $g\in L^{2}(B;\R^{3})$ with
$\norm{g}_{L^{2}}\le M$. For weak convergence, let $w\in L^{2}(B;\R^{3})$ and
$\varepsilon>0$, and choose $\varphi\in C_{c}^{\infty}(B;\R^{3})$ with
$\norm{w-\varphi}_{L^{2}(B)}\le\varepsilon/(4M+1)$. Then
\[
  \Bigl|\int_{B}(g_{k}-g)\cdot w\Bigr|
  \le\Bigl|\int_{B}(g_{k}-g)\cdot(w-\varphi)\Bigr|
    +\Bigl|\int_{B}(g_{k}-g)\cdot\varphi\Bigr|
  \le 2M\,\frac{\varepsilon}{4M+1}
    +\abs{\langle g_{k}-g,\varphi\rangle} ,
\]
and the last term tends to $0$ because $g_{k}\to g$ in $Z_{B}$ and
$\varphi\in H^{2}_{0}(B;\R^{3})$. Hence
\[
\limsup_{k}\abs{\int_{B}(g_{k}-g)\cdot w}\le\varepsilon/2<\varepsilon
.
\]
Applying the claim with $g_{k}=\bar u(t_{k})$, $g=\bar u(t_{0})$, $M=M_{B,J}$
gives $\bar u(t_{0})\in L^{2}(B;\R^{3})$ with
$\norm{\bar u(t_{0})}_{L^{2}(B)}\le M_{B,J}$, which is (ii).

\emph{Step 3: weak continuity (iii).} Let $t_{0}\in\overline{J}$ and let
$t_{k}\to t_{0}$ in $\overline{J}$ be arbitrary. By Step 2 each
$\bar u(t_{k})\in L^{2}(B)$ with norm at most $M_{B,J}$, and
$\bar u(t_{k})\to\bar u(t_{0})$ in $Z_{B}$ by Step 1. The claim of Step 2 gives
$\bar u(t_{k})\rightharpoonup\bar u(t_{0})$ weakly in $L^{2}(B;\R^{3})$, i.e.\
$\int_{B}\bar u(t_{k})\cdot w\to\int_{B}\bar u(t_{0})\cdot w$ for every
$w\in L^{2}(B;\R^{3})$. Since $t_{k}\to t_{0}$ was arbitrary, (iii) follows.
\end{proof}

\begin{lemma}[A common null set for the divergence-free condition]
\label{lem:divfree-slice}
Let $u\in L^{1}_{loc}(\mathcal{O};\R^{3})$ satisfy $\dv u=0$ in
$\mathcal{D}'(\mathcal{O})$, let $B=\Ball{\rho}{x_{0}}$ with
$\overline{B}\subset\Omega$, let $J$ be an interval with $\overline{J}\subset I$,
and let $E\subset B$ be compact. Then there is a Lebesgue-null set $N\subset J$
such that
\begin{equation}\label{eq:divfree-common}
  \int_{B}u(y,s)\cdot\nabla\psi(y)\,dy=0
  \qquad\text{for every }s\in J\setminus N\text{ and every }
  \psi\in C^{1}_{c}(B)\text{ with }\supp\psi\subset E .
\end{equation}
\end{lemma}

\begin{proof}
Let $E'$ be a compact neighbourhood of $E$ in $B$. Restriction of a
function and its three first derivatives embeds
$C^{1}_{E}=\{\psi\in C^{1}(\R^{3}):\supp\psi\subset E\}$ isometrically
into $C(E')^{4}$. Since $E'$ is compact metrizable, this space and its
subspaces are separable. Let
$\{\psi_{k}\}_{k\in\N}\subset C_{c}^{\infty}(B)$ with $\supp\psi_{k}\subset E'$
be a family whose $C^{1}$ closure contains $C^{1}_{E}$. To construct it,
mollify a countable dense subset of $C^{1}_{E}$ at a countable sequence of
sufficiently small radii. The supports then lie in $E'$, and the mollifications
converge in $C^{1}$ (External Input~\ref{ext:mollify}). Testing (S2) with
$\psi_{k}(x)\theta(t)$, $\theta\in C_{c}^{\infty}(J)$, gives
$\int_{J}\theta(s)\bigl(\int_{B}u(\cdot,s)\cdot\nabla\psi_{k}\bigr)ds=0$ for all
$\theta$, hence a null set $N_{k}\subset J$ off which
$\int_{B}u(\cdot,s)\cdot\nabla\psi_{k}=0$. Let
$N=\bigl(\bigcup_{k}N_{k}\bigr)\cup\{s:\norm{u(\cdot,s)}_{L^{1}(B)}=\infty\}$,
a null set by Fubini. For $s\notin N$ and $\psi\in C^{1}_{E}$, choose
$\psi_{k}\to\psi$ in $C^{1}$; then
$\bigl|\int_{B}u(\cdot,s)\cdot\nabla(\psi-\psi_{k})\bigr|
 \le\norm{u(\cdot,s)}_{L^{1}(B)}\norm{\nabla\psi-\nabla\psi_{k}}_{L^{\infty}}\to0$,
which gives \eqref{eq:divfree-common}.
\end{proof}

\begin{lemma}[Pointwise-in-time local energy inequality]\label{lem:lei-pointwise}
Let $(u,p)$ be a suitable weak solution on $\mathcal{O}$ with force
$f\in L^{q}_{loc}$, $q>5/2$. Let $\phi\in C_{c}^{\infty}(\mathcal{O})$
with $\phi\ge0$, and let $J=(t_{-},t_{+})$ be a bounded open interval with
$\overline{J}\subset I$ and
\begin{equation}\label{eq:phi-vanishes}
  \phi(\cdot,s)=0\quad\text{for \emph{all} }s\le t_{-}+\varepsilon'
  \qquad\text{for some }\varepsilon'>0 ,
\end{equation}
so that in particular $\supp\phi\subset\Omega\times(t_{-},\infty)$. Then for
almost every $t\in J$,
\begin{multline}\label{eq:lei-pointwise}
  \int_{\Omega}\abs{u(x,t)}^{2}\phi(x,t)\,dx
  \;+\;2\int_{t_{-}}^{t}\!\!\int_{\Omega}\abs{\nabla u}^{2}\phi\,dx\,ds\\
  \le\;\int_{t_{-}}^{t}\!\!\int_{\Omega}
    \Bigl[\abs{u}^{2}\bigl(\partial_{s}\phi+\Delta\phi\bigr)
      +\bigl(\abs{u}^{2}+2p\bigr)u\cdot\nabla\phi
      +2(f\cdot u)\phi\Bigr]dx\,ds .
\end{multline}
\end{lemma}

\begin{proof}
Fix $t\in J$ and let $K=\supp\phi$, a compact subset of $\mathcal{O}$.
Cover its spatial projection by finitely many balls whose closures lie in
$\Omega$, and let $B$ denote their union. Choose an open interval $J'$
with $\overline{J'}\subset I$, $K\subset B\times J'$ and
$\overline{J}\subset J'$. The local integrability bounds used below hold
on each ball and hence on their finite union. All spatial integrals below
may be restricted to $B$, since the test function vanishes outside it.

\emph{Step 1: a Lipschitz cut-off in time.} For $0<h<t-t_{-}$ let
$\chi_{h}:\R\to[0,1]$ be the piecewise affine function
\begin{equation}\label{eq:time-cutoff}
  \chi_{h}(s)=\begin{cases}
    1,& s\le t-h,\\[2pt]
    \dfrac{t-s}{h},& t-h<s<t,\\[6pt]
    0,& s\ge t .
  \end{cases}
\end{equation}
Then $\chi_{h}$ is Lipschitz with
$\chi_{h}'=-h^{-1}\mathbf{1}_{(t-h,t)}$ almost everywhere. Fix
$\zeta\in C_{c}^{\infty}((-1,1))$ with $\zeta\ge0$ and $\int_{\R}\zeta=1$, put
$\zeta_{\varepsilon}(s)=\varepsilon^{-1}\zeta(s/\varepsilon)$, and set
$\chi_{h,\varepsilon}=\chi_{h}*\zeta_{\varepsilon}$, which is smooth, takes
values in $[0,1]$, equals $1$ for $s\le t-h-\varepsilon$ and $0$ for
$s\ge t+\varepsilon$. For $\varepsilon$ small enough that
$t+\varepsilon<t_{+}$, the function
$\phi\,\chi_{h,\varepsilon}$ belongs to $C_{c}^{\infty}(\mathcal{O})$ and is
nonnegative, so \eqref{eq:lei} applies to it. Using
$\partial_{s}(\phi\chi_{h,\varepsilon})
 =(\partial_{s}\phi)\chi_{h,\varepsilon}+\phi\,\chi_{h,\varepsilon}'$
and $\Delta(\phi\chi_{h,\varepsilon})=(\Delta\phi)\chi_{h,\varepsilon}$,
inequality \eqref{eq:lei} becomes
\begin{multline}\label{eq:lei-cut}
  2\iint\abs{\nabla u}^{2}\phi\,\chi_{h,\varepsilon}
  \;-\;\iint\abs{u}^{2}\phi\,\chi_{h,\varepsilon}'\\
  \le\;\iint\Bigl[\abs{u}^{2}(\partial_{s}\phi+\Delta\phi)
    +(\abs{u}^{2}+2p)u\cdot\nabla\phi+2(f\cdot u)\phi\Bigr]\chi_{h,\varepsilon}.
\end{multline}

\emph{Step 2: $\varepsilon\to0$ at fixed $h$.} Each integrand in
\eqref{eq:lei-cut} other than the one containing $\chi'$ is in
$L^{1}(B\times J')$: indeed $\abs{\nabla u}^{2}\phi\in L^{1}$ by (S1),
$\abs{u}^{2}(\partial_{s}\phi+\Delta\phi)\in L^{1}$ because
$u\in L^{10/3}(B\times J')\subset L^{2}$, $\abs{u}^{2}u\cdot\nabla\phi\in L^{1}$
because $u\in L^{10/3}\subset L^{3}$, $p\,u\cdot\nabla\phi\in L^{1}$ by
H\"older with $\tfrac23+\tfrac13=1$ and $p\in L^{3/2}$, $u\in L^{3}$, and
$(f\cdot u)\phi\in L^{1}$ by H\"older with $f\in L^{q}$, $q>5/2$, and
$u\in L^{10/3}$ (note $\tfrac1q+\tfrac{3}{10}<\tfrac25+\tfrac3{10}<1$). Since
$0\le\chi_{h,\varepsilon}\le1$ and $\chi_{h,\varepsilon}\to\chi_{h}$ pointwise
away from the two corner points as $\varepsilon\to0$, dominated convergence
applies to all these terms. For the remaining term, $\chi_{h,\varepsilon}'
=\chi_{h}'*\zeta_{\varepsilon}$ is bounded by $h^{-1}$ and supported in
$(t-h-\varepsilon,t+\varepsilon)$, and $\chi_{h,\varepsilon}'\to\chi_{h}'$ in
$L^{1}(\R)$; since $s\mapsto\int\abs{u(x,s)}^{2}\phi(x,s)\,dx$ is in
$L^{\infty}(J')$ by (S1), we get
$\iint\abs{u}^{2}\phi\,\chi_{h,\varepsilon}'\to
 -h^{-1}\int_{t-h}^{t}\int\abs{u}^{2}\phi\,dx\,ds$. Hence
\begin{multline}\label{eq:lei-h}
  2\int_{t_{-}}^{t}\!\!\int\abs{\nabla u}^{2}\phi\,\chi_{h}
  \;+\;\frac1h\int_{t-h}^{t}\!\!\int\abs{u}^{2}\phi\,dx\,ds\\
  \le\;\int_{t_{-}}^{t}\!\!\int\Bigl[\abs{u}^{2}(\partial_{s}\phi+\Delta\phi)
    +(\abs{u}^{2}+2p)u\cdot\nabla\phi+2(f\cdot u)\phi\Bigr]\chi_{h} .
\end{multline}

\emph{Step 3: $h\to0$.} On the right of \eqref{eq:lei-h} we have
$0\le\chi_{h}\le\mathbf{1}_{(-\infty,t)}$ and $\chi_{h}\to\mathbf{1}_{(-\infty,t)}$
pointwise, so dominated convergence gives the right side of
\eqref{eq:lei-pointwise}. In the first term on the left, $\chi_{h}$ increases to
$\mathbf{1}_{(-\infty,t)}$ and the integrand is nonnegative, so monotone
convergence gives $2\int_{t_{-}}^{t}\int\abs{\nabla u}^{2}\phi$. It remains to
prove
\begin{equation}\label{eq:lsc-average}
  \int\abs{u(x,t)}^{2}\phi(x,t)\,dx
  \;\le\;\liminf_{h\to0^{+}}\frac1h\int_{t-h}^{t}\!\!\int\abs{u(x,s)}^{2}\phi(x,s)\,dx\,ds .
\end{equation}
Write $\Lambda(s)=\int|u(x,s)|^2\phi(x,s)\,dx$.
This function is locally integrable by (S1). At every one-sided
Lebesgue point $t$ of $\Lambda$,
\[
 \lim_{h\downarrow0}\frac1h\int_{t-h}^t\Lambda(s)\,ds=\Lambda(t).
\]
Almost every $t$ has this property. This gives \eqref{eq:lsc-average}
(in fact equality) at precisely the times required by the statement.
The remaining terms converge by the preceding integrability and
monotone-convergence arguments.
\end{proof}

\begin{remark}\label{rem:lei-direction}
At the Lebesgue times used above the average converges to the slice
energy. Only the corresponding lower bound is needed for the left side
of \eqref{eq:lei-pointwise}. Every subsequent energy estimate uses an
essential supremum or an integral in time, so it needs no assertion at
an exceptional time. The earlier weak-time-continuity discussion provides
a stronger optional representative convention, not an input to this route.
\end{remark}

The local product-domain formulation loses no generality for local
regularity on an open spacetime set: cover that set by relatively compact
products $\Omega'\times J$. All weak-equation and energy pairings are
integrable on the support of a test, by (S1), local interpolation, and the
force assumption.

\subsection{Scale-invariant quantities}\label{ss:quantities}

Let $(u,p)$ be a suitable weak solution on $\mathcal{O}$ with force $f$, let
$z=(x,t)\in\mathcal{O}$ and let $r>0$ be such that
$\overline{\Cyl{r}{z}}\subset\mathcal{O}$. The five quantities we work with are those of
Kukavica \cite[\S2]{Kukavica2008}:
\begin{equation}\label{eq:alpha-beta}
\begin{aligned}
  \alpha(z,r)&=\Bigl(\frac1r\operatorname*{ess\,sup}_{s\in(t-r^{2},t)}
      \int_{\Ball{r}{x}}\abs{u(y,s)}^{2}dy\Bigr)^{1/2},
  &\quad
  \beta(z,r)&=\Bigl(\frac1r\iint_{\Cyl{r}{z}}\abs{\nabla u}^{2}dz'\Bigr)^{1/2},\\[4pt]
  \gamma(z,r)&=\Bigl(\frac1{r^{2}}\iint_{\Cyl{r}{z}}\abs{u}^{3}dz'\Bigr)^{1/3},
  &\quad
  \delta(z,r)&=\Bigl(\frac1{r^{2}}\iint_{\Cyl{r}{z}}\abs{p}^{3/2}dz'\Bigr)^{1/3},
\end{aligned}
\end{equation}
\begin{equation}\label{eq:lambda}
  \lambda(z,r)=r^{\sigma}\Bigl(\iint_{\Cyl{r}{z}}\abs{f}^{q}dz'\Bigr)^{1/q}
              =r^{\sigma}\norm{f}_{L^{q}(\Cyl{r}{z})},
  \qquad \sigma=3-\tfrac5q>1 .
\end{equation}
When the base point is the origin we write $\alpha(r)=\alpha((0,0),r)$, and so
on. All five are dimensionless, in the precise sense of
Lemma~\ref{lem:scaling-quantities}: each is unchanged when $(u,p,f)$ is replaced
by the rescaled triple \eqref{eq:rescaling} and $r$ by $r/\mu$. For $q=2$ one has
$\sigma=1/2$ and $\lambda(z,r)=\bigl(r\iint_{\Cyl{r}{z}}\abs{f}^{2}\bigr)^{1/2}$,
which is the force functional of \cite[\S2]{Kukavica2008}.

\begin{lemma}[Dictionary to the Lin, CKN and Ladyzhenskaya--Seregin notation]
\label{lem:dictionary}
Put
\begin{equation}\label{eq:ABCDE}
\begin{aligned}
  A(z,r)&=\frac1r\operatorname*{ess\,sup}_{s\in(t-r^{2},t)}\int_{\Ball{r}{x}}\abs{u}^{2},
  &\quad
  B(z,r)&=\frac1r\iint_{\Cyl{r}{z}}\abs{\nabla u}^{2},
  &\quad
  C(z,r)&=\frac1{r^{2}}\iint_{\Cyl{r}{z}}\abs{u}^{3},\\
  D(z,r)&=\frac1{r^{2}}\iint_{\Cyl{r}{z}}\abs{p}^{3/2},
  &\quad
  E_{q}(z,r)&=\lambda(z,r).
\end{aligned}
\end{equation}
Then
\begin{equation}\label{eq:dictionary}
  \alpha=A^{1/2},\qquad\beta=B^{1/2},\qquad\gamma=C^{1/3},
  \qquad\delta=D^{1/3},\qquad\lambda=E_{q} .
\end{equation}
\end{lemma}

\begin{proof}
Each identity is the definition \eqref{eq:alpha-beta} with the outer exponent
removed, compared with \eqref{eq:ABCDE}.
\end{proof}

\begin{remark}[Letters used differently by different authors]\label{rem:letters}
$A$, $B$, $C$, $D$ are Lin's notation \cite[\S3]{Lin1998}, which is also that of
\cite{CKN1982}. Ladyzhenskaya--Seregin \cite[\S5]{LadyzhenskayaSeregin1999} and
Tsai \cite[\S6.3]{Tsai2018} write $E(r)$ for what is here $B(z,r)$, i.e.\ for the
scale-invariant Dirichlet energy; we never use the letter $E$ alone, precisely
to avoid this collision, and we write $E_{q}$ only as an alias for $\lambda$ when
comparing with \cite{Tsai2018}. Tsai \cite[\S6.3]{Tsai2018} also uses
$\widetilde C$ for the velocity excess, which agrees with \eqref{eq:excess-CD}
below.
\end{remark}

We shall also use the two \emph{excess} (oscillation) quantities, in which the
velocity is measured against its space-time mean and the pressure against its
spatial mean at each time,
\begin{equation}\label{eq:excess-CD}
  \widetilde C(z,r)
    =\frac1{r^{2}}\iint_{\Cyl{r}{z}}\bigl|u-\stavg{u}{\Cyl{r}{z}}\bigr|^{3}dz',
  \qquad
  \widetilde D(z,r)
    =\frac1{r^{2}}\iint_{\Cyl{r}{z}}
       \bigl|p(y,s)-\spavg{p}{\Ball{r}{x}}(s)\bigr|^{3/2}dy\,ds ,
\end{equation}
the \emph{drift functional}
\begin{equation}\label{eq:drift}
  \Psi(z,r)\;=\;r\,\bigl|\stavg{u}{\Cyl{r}{z}}\bigr| ,
\end{equation}
and the \emph{excess} of $(u,p)$ on $\Cyl{r}{z}$, in the normalization of
\cite[(6.14)]{Tsai2018},
\begin{equation}\label{eq:excess}\begin{aligned}
  &\varphi(z,r)
  \;=\;\widetilde C(z,r)^{1/3}+\widetilde D(z,r)^{2/3}
  \\ &\;=\;\Bigl(\frac1{r^{2}}\iint_{\Cyl{r}{z}}\bigl|u-\stavg{u}{\Cyl{r}{z}}\bigr|^{3}\Bigr)^{1/3}
      +\Bigl(\frac1{r^{2}}\iint_{\Cyl{r}{z}}\bigl|p-\spavg{p}{\Ball{r}{x}}\bigr|^{3/2}\Bigr)^{2/3}.
\end{aligned}\end{equation}
Both $\varphi$ and $\Psi$ are dimensionless
(Lemma~\ref{lem:scaling-quantities}).

\begin{remark}[Why the two means are different]\label{rem:two-means}
In \eqref{eq:excess} the velocity is compared with a \emph{space-time} mean and
the pressure with a \emph{spatial} mean at each time. These choices serve
different purposes.

The pressure of a suitable weak solution is determined only up to an additive
function of time: if $(u,p)$ is a suitable weak solution with force $f$ and
$c\in L^{3/2}_{loc}(I)$, then $(u,p+c)$ is one as well, since $\nabla c=0$ and
$c\,u\cdot\nabla\phi$ integrates to $0$ against $\dv u=0$ after an integration
by parts. Hence no statement about $p$ can be invariant unless a spatial mean is
subtracted at each time.

The velocity excess uses a space-time mean because the parabolic Campanato
characterization of H\"older continuity (External Input~\ref{ext:campanato})
measures oscillation in both space and time. A spatial mean alone removes
all information about a spatially constant, time-dependent velocity. This
observation explains the choice of excess; it is not a counterexample to
regularity under the full suitable-solution hypotheses.
\end{remark}

\begin{lemma}[Comparison of the excess with $\gamma,\delta$]\label{lem:excess-vs}
For every $z$ and every admissible $r$,
\begin{equation}\label{eq:excess-vs-CD}
  \widetilde C(z,r)\le 8\,C(z,r),
  \qquad
  \widetilde D(z,r)\le 2^{3/2}D(z,r),
  \qquad
  \varphi(z,r)\le2\gamma(z,r)+2\delta(z,r)^{2},
\end{equation}
and $\Psi(z,r)\le C_{3}\,\gamma(z,r)$ with $C_{3}=(4\pi/3)^{-1/3}$, an absolute
constant.
\end{lemma}

\begin{proof}
Jensen's inequality, applied to $\tau\mapsto\tau^{3}$ with the normalized
Lebesgue measure on $\Cyl{r}{z}$, gives
\[
\bigl|\stavg{u}{\Cyl{r}{z}}\bigr|^{3}
 \le\bigl(\abs{\Cyl{r}{z}}^{-1}\iint\abs{u}\bigr)^{3}
 \le\abs{\Cyl{r}{z}}^{-1}\iint\abs{u}^{3}
,
\] hence
\[
\iint_{\Cyl{r}{z}}\bigl|\stavg{u}{\Cyl{r}{z}}\bigr|^{3}\le\iint_{\Cyl{r}{z}}\abs{u}^{3}
.
\]
Combining with $\abs{a-b}^{3}\le4(\abs{a}^{3}+\abs{b}^{3})$ gives
$\widetilde C\le4(C+C)=8C$. The same argument with $\tau\mapsto\tau^{3/2}$,
applied for each fixed $s$ on $\Ball{r}{x}$, gives
\[
\int_{\Ball{r}{x}}\bigl|\spavg{p}{\Ball{r}{x}}(s)\bigr|^{3/2}dy
 \le\int_{\Ball{r}{x}}\abs{p(\cdot,s)}^{3/2}dy
,
\] and with
$\abs{a-b}^{3/2}\le2^{1/2}(\abs{a}^{3/2}+\abs{b}^{3/2})$ this gives
$\widetilde D\le2^{1/2}\cdot2D=2^{3/2}D$. Taking the $1/3$ and $2/3$ powers,
$\widetilde C^{1/3}\le8^{1/3}C^{1/3}=2\gamma$ and
$\widetilde D^{2/3}\le(2^{3/2})^{2/3}D^{2/3}=2\delta^{2}$, which is the third
inequality. Finally, by Jensen again,
\[
\begin{aligned}
&\abs{\stavg{u}{\Cyl{r}{z}}}
\\ &\le\bigl(\abs{\Cyl{r}{z}}^{-1}\iint\abs{u}^{3}\bigr)^{1/3}
\\ &=\bigl(\tfrac{4\pi}{3}\bigr)^{-1/3}r^{-5/3}\bigl(r^{2}\gamma^{3}\bigr)^{1/3}
\\ &=(\tfrac{4\pi}{3})^{-1/3}r^{-1}\gamma
\end{aligned}
,
\] so
$\Psi=r\abs{\stavg{u}{\Cyl{r}{z}}}\le(\tfrac{4\pi}{3})^{-1/3}\gamma$.
\end{proof}

\begin{remark}[Normalizations that are \emph{not} used]\label{rem:no-lin-normalization}
Lin \cite{Lin1998} writes $v_{\theta}=\theta^{-5}\iint_{Q_{\theta}}v$ and
$P_{\theta}(t)=\theta^{-3}\int_{B_{\theta}}P$. These differ from the genuine
averages \eqref{eq:spatial-average}, \eqref{eq:spacetime-average} by the fixed
factors $\abs{Q_{1}}=\tfrac{4\pi}{3}$ and $\abs{B_{1}}=\tfrac{4\pi}{3}$. We use
only genuine averages: every mean in this paper divides by the measure of the
set it is taken over. This matters for the Campanato step
(External Input~\ref{ext:campanato}), which is stated for genuine averages, and
for the Poincar\'e and Jensen inequalities, which are false with a mismatched
normalizing factor.
\end{remark}

\section{Statements}\label{sec:theorems}

\subsection{Regular and singular points}\label{ss:regular}

\begin{definition}[Parabolic H\"older spaces]\label{def:holder}
Let $E\subset\R^{3}\times\R$ and $\gamma\in(0,1]$. A function
$w:E\to\R^{m}$ belongs to $C^{\gamma,\gamma/2}(E)$ if it is bounded and
\[
  [w]_{C^{\gamma,\gamma/2}(E)}
  \;=\;\sup_{z\neq z'\in E}\frac{\abs{w(z)-w(z')}}{\dpar(z,z')^{\gamma}}
  \;<\;\infty ,
\]
with $\dpar$ as in \eqref{eq:parabolic-metric}. We set
$\norm{w}_{C^{\gamma,\gamma/2}(E)}
 =\norm{w}_{L^{\infty}(E)}+[w]_{C^{\gamma,\gamma/2}(E)}$.
\end{definition}

\begin{definition}[Regular and singular points]\label{def:regular}
Let $(u,p)$ be a suitable weak solution on $\mathcal{O}=\Omega\times I$ with
force $f$. A point $z_{0}\in\mathcal{O}$ is a \emph{regular point of $u$} if
there exist an open set $\mathcal{N}$ with $z_{0}\in\mathcal{N}\subset\mathcal{O}$,
an exponent $\gamma\in(0,1]$, and a function $w\in C^{\gamma,\gamma/2}(\mathcal{N})$
such that $u=w$ almost everywhere in $\mathcal{N}$. Otherwise $z_{0}$ is a
\emph{singular point}. The \emph{singular set} is
\[
  S \;=\; S(u)\;=\;\{z_{0}\in\mathcal{O}: z_{0}\text{ is a singular point of }u\} .
\]
\end{definition}

\begin{lemma}\label{lem:S-closed}
$S$ is relatively closed in $\mathcal{O}$.
\end{lemma}

\begin{proof}
Let $R=\mathcal{O}\setminus S$ be the set of regular points and let
$z_{0}\in R$, with $\mathcal{N},\gamma,w$ as in Definition~\ref{def:regular}.
Every $z_{1}\in\mathcal{N}$ is then regular, witnessed by the same
$\mathcal{N}$, $\gamma$ and $w$. Hence $\mathcal{N}\subset R$, so $R$ is open in
$\mathcal{O}$ and $S$ is relatively closed.
\end{proof}

\begin{definition}[Parabolic Hausdorff measure]\label{def:parabolic-hausdorff}
For $\alpha\ge0$, $\delta>0$ and $E\subset\R^{3}\times\R$ put
\[
\begin{aligned}
  &\Phaus^{\alpha}_{\delta}(E)
  =\inf\Bigl\{\sum_{j\in J}r_{j}^{\alpha}\ :\
     J\subset\N,\quad
     E\subset\bigcup_{j\in J}\mathfrak{B}_{r_{j}}(z_{j}),\
     0<r_{j}\le\delta\ (j\in J)\Bigr\},\\
  &\Phaus^{\alpha}(E)=\lim_{\delta\to0^{+}}\Phaus^{\alpha}_{\delta}(E)
  =\sup_{\delta>0}\Phaus^{\alpha}_{\delta}(E),
\end{aligned}
\]
where empty and finite $J$ are allowed, an empty sum is zero, and
$\mathfrak{B}_{r}(z)$ is the $\dpar$-ball \eqref{eq:parabolic-ball}.
For $\alpha=0$ each nonempty covering ball has cost $1$.
The same index convention is used for the cylinder gauge below. Thus $\Phaus^{\alpha}$ is the \emph{spherical} $\alpha$-dimensional Hausdorff
measure of the metric space $(\R^{3}\times\R,\dpar)$: the gauge is the radius of
a covering ball rather than the diameter of an arbitrary covering set
(External Input~\ref{ext:hausdorff}). The two differ by a factor between $1$ and
$2^{\alpha}$ and in particular have the same null sets, which is all that is
used.
\end{definition}

\begin{remark}[Gauge independence]\label{rem:hausdorff-gauge}
Definition~\ref{def:parabolic-hausdorff} uses the \emph{radius} of a covering
ball as gauge. The Hausdorff measure of $(\R^{3}\times\R,\dpar)$ built, as in
the usual measure-theoretic convention, from covers by arbitrary sets $E_{j}$ of
small $\dpar$-diameter and the gauge $(\diam_{\dpar}E_{j})^{\alpha}$, denoted by
$\Phaus^{\alpha}_{\mathrm{diam}}$, satisfies
\begin{equation}\label{eq:gauge-comparison}
  \Phaus^{\alpha}_{\mathrm{diam}}(E)\;\le\;2^{\alpha}\,\Phaus^{\alpha}(E)
  \qquad\text{and}\qquad
  \Phaus^{\alpha}(E)\;\le\;2^{\alpha}\,\Phaus^{\alpha}_{\mathrm{diam}}(E)
\end{equation}
for every $E$. Indeed, a ball $\mathfrak{B}_{r}(z)$ has
$\diam_{\dpar}\le2r$, which gives the first inequality; conversely a non-empty set $E_{j}$
of diameter $d_{j}$ containing a point $z_{j}$ satisfies
$E_{j}\subset\mathfrak{B}_{r_{j}}(z_{j})$ for every $r_{j}>d_{j}$, and letting
$r_{j}\downarrow d_{j}$ gives the second, in fact with factor $1$ for
$\alpha>0$: choose the radii with arbitrarily small summable excess in
$\sum_j r_j^\alpha$, including arbitrarily small positive radii when $d_j=0$. At $\alpha=0$ both
measures are counting measure: a finite set of $n$ points requires $n$
sufficiently small covering sets, and every infinite set contains finite
subsets of arbitrarily large cardinality. Empty covering sets are discarded
at zero cost in the diameter convention. Thus both comparisons include
$\alpha=0$ and $E=\varnothing$. In particular the two have exactly the
same null sets, and Theorem~\ref{thm:C} (an assertion that a set is
$\Phaus^{1}$-null) is unaffected by the choice of gauge. The same applies to
the cylinder gauge of Lemma~\ref{lem:P1-cylinders}.
\end{remark}

\begin{lemma}[Cylinders and parabolic balls give the same null sets]\label{lem:P1-cylinders}
Define $\widehat{\Phaus}^{\alpha}$ exactly as in
Definition~\ref{def:parabolic-hausdorff} but with the covering sets
$\mathfrak{B}_{r_{j}}(z_{j})$ replaced by cylinders $\Cyl{r_{j}}{z_{j}}$. Then
for every $E$ and every $\alpha\ge0$,
\[
  \Phaus^{\alpha}(E)\;\le\;\widehat{\Phaus}^{\alpha}(E)
  \;\le\;2^{\alpha}\,\Phaus^{\alpha}(E) .
\]
In particular $\Phaus^{\alpha}(E)=0$ if and only if
$\widehat{\Phaus}^{\alpha}(E)=0$.
\end{lemma}

\begin{proof}
If $E\subset\bigcup_{j}\Cyl{r_{j}}{z_{j}}$ with $r_{j}\le\delta$ then
$E\subset\bigcup_{j}\mathfrak{B}_{r_{j}}(z_{j})$ by
Lemma~\ref{lem:parabolic-metric}, with the same radii; hence
$\Phaus^{\alpha}_{\delta}(E)\le\widehat{\Phaus}^{\alpha}_{\delta}(E)$.
Conversely, if $E\subset\bigcup_{j}\mathfrak{B}_{r_{j}}(z_{j})$ with
$z_{j}=(x_{j},t_{j})$ and $r_{j}\le\delta$, then by
Lemma~\ref{lem:parabolic-metric}
$\mathfrak{B}_{r_{j}}(z_{j})\subset\Cyl{2r_{j}}{(x_{j},t_{j}+r_{j}^{2})}$, so
$E$ is covered by cylinders of radii $2r_{j}\le2\delta$ and
$\widehat{\Phaus}^{\alpha}_{2\delta}(E)\le\sum_{j}(2r_{j})^{\alpha}
 =2^{\alpha}\sum_{j}r_{j}^{\alpha}$. Taking the infimum and then
$\delta\to0^{+}$ gives the second inequality.
\end{proof}

\subsection{The three theorems}\label{ss:main-theorems}

\begin{theorem}[$\varepsilon$-regularity, $L^{3}$ form]\label{thm:A}
For every $q>5/2$ there exist $\varepsilon_{0}=\varepsilon_{0}(q)>0$,
$\gamma_{0}=\gamma_{0}(q)\in(0,\tfrac23]$ and $C_{4}=C_{4}(q)<\infty$ with the
following property. Let $(u,p)$ be a suitable weak solution on
$\mathcal{O}=\Omega\times I$, with $\overline{Q_{1}}\subset\mathcal{O}$, and force $f\in L^{q}(Q_{1};\R^{3})$, and
suppose that, with the notation \eqref{eq:alpha-beta},
\begin{equation}\label{eq:thmA-hyp}
  \iint_{Q_{1}}\bigl(\abs{u}^{3}+\abs{p}^{3/2}+\abs{f}^{q}\bigr)\,dz
  \;\le\;\varepsilon_{0},
  \qquad\text{equivalently}\qquad
  \gamma(1)^{3}+\delta(1)^{3}+\lambda(1)^{q}\le\varepsilon_{0}.
\end{equation}
Then $u$ agrees almost everywhere on $Q_{1/2}$ with a function
$w\in C^{\gamma_{0},\gamma_{0}/2}(\overline{Q_{1/2}};\R^{3})$ satisfying
$\norm{w}_{C^{\gamma_{0},\gamma_{0}/2}(\overline{Q_{1/2}})}\le C_{4}$. In
particular every point of the \emph{open} cylinder
$\Ball{1/2}{0}\times(-\tfrac14,0)$ is a regular point of $u$.
\end{theorem}

\begin{remark}[The top face]\label{rem:top-face}
Definition~\ref{def:regular} requires an \emph{open} neighbourhood, whereas
$Q_{1/2}$ contains its top face $\{t=0\}$; Theorem~\ref{thm:A} therefore does
not by itself declare the points of that face regular. This costs nothing,
because Theorem~\ref{thm:B} is not deduced from Theorem~\ref{thm:A} by applying
it at $z_{0}$: it is proved in \S\ref{ss:assembly} from
Proposition~\ref{prop:morrey-decay}, whose conclusion
\eqref{eq:morrey-decay} holds at every centre of the \emph{open}
$\dpar$-ball $\mathfrak{B}_{r_{2}}(z_{0})$, and from
Theorem~\ref{thm:endgame}, which produces a H\"older representative on an open
$\dpar$-ball. The upper-semicontinuity Lemma~\ref{lem:alpha-usc} is exactly what
makes the passage to later base points legitimate.
\end{remark}

\begin{theorem}[$\varepsilon$-regularity, gradient form]\label{thm:B}
For every $q>5/2$ there exists $\varepsilon_{1}=\varepsilon_{1}(q)>0$ with the
following property. Let $(u,p)$ be a suitable weak solution on
$\mathcal{O}=\Omega\times I$ with force $f\in L^{q}_{loc}(\mathcal{O};\R^{3})$,
and let $z_{0}\in\mathcal{O}$. Suppose that the following inequality holds,
with both limsups taken in $[0,\infty]$:\footnote{Suitability and interior
geometry make each sufficiently small cylinder mass finite, so the squared
real $\beta$ expression agrees with the nonnegative integral expression.}
\begin{equation}\label{eq:thmB-hyp}
  \limsup_{r\to0^{+}}\;\beta(z_{0},r)^{2}
  \;=\;\limsup_{r\to0^{+}}\;\frac1r\iint_{\Cyl{r}{z_{0}}}\abs{\nabla u}^{2}\,dz
  \;<\;\varepsilon_{1}^{2}.
\end{equation}
Then $z_{0}$ is a regular point of $u$.
\end{theorem}

\begin{theorem}[Caffarelli-Kohn-Nirenberg]\label{thm:C}
Let $(u,p)$ be a suitable weak solution on $\mathcal{O}=\Omega\times I$ with
force $f\in L^{q}_{loc}(\mathcal{O};\R^{3})$ for some $q>5/2$. Then the singular
set $S$ satisfies $\Phaus^{1}(S)=0$.
\end{theorem}

Theorems~\ref{thm:A} and \ref{thm:B} are proved in \S\ref{ss:assembly};
Theorem~\ref{thm:C} is deduced from Theorem~\ref{thm:B} in
Section~\ref{sec:covering}.

\begin{remark}[Strictness in \eqref{eq:thmB-hyp}]\label{rem:strict}
The inequality in \eqref{eq:thmB-hyp} is strict, as in
\cite[Theorem 2.2]{LadyzhenskayaSeregin1999} and
\cite[Theorem 1]{Kukavica2008}; the proof produces a threshold
$\varepsilon_{1}$ and needs $B(z_{0},r)<\varepsilon_{1}^{2}$ for all sufficiently
small $r$, which is what the strict $\limsup$ inequality supplies. For
$0<\varepsilon_{1}'<\varepsilon_{1}$, the nonstrict hypothesis
$\limsup_{r\downarrow0}B(z_{0},r)\le(\varepsilon_{1}')^{2}<\varepsilon_{1}^{2}$ suffices.
\end{remark}

\section{Invariances and interpolation}\label{sec:scaling}

\subsection{Translation and parabolic rescaling}\label{ss:scaling}

\begin{definition}\label{def:rescaling}
Let $z_{0}=(x_{0},t_{0})\in\R^{3}\times\R$ and $\mu>0$. For functions
$u,p,f$ defined on $\mathcal{O}=\Omega\times I$ set
\begin{equation}\label{eq:rescaling}\begin{gathered}
  u^{\mu,z_{0}}(y,s)=\mu\,u(x_{0}+\mu y,\;t_{0}+\mu^{2}s),
  \\
  p^{\mu,z_{0}}(y,s)=\mu^{2}\,p(x_{0}+\mu y,\;t_{0}+\mu^{2}s),
  \\
  f^{\mu,z_{0}}(y,s)=\mu^{3}\,f(x_{0}+\mu y,\;t_{0}+\mu^{2}s),\end{gathered}
\end{equation}
defined on
$\mathcal{O}^{\mu,z_{0}}
 =\bigl(\mu^{-1}(\Omega-x_{0})\bigr)\times\bigl(\mu^{-2}(I-t_{0})\bigr)$.
\end{definition}

\begin{lemma}[The class is invariant]\label{lem:scaling-invariance}
Let $q>5/2$. If $(u,p)$ is a suitable weak solution on $\mathcal{O}$ with force
$f\in L^{q}_{loc}(\mathcal{O};\R^{3})$, then for every $z_{0}$ and every
$\mu>0$ the pair $(u^{\mu,z_{0}},p^{\mu,z_{0}})$ is a suitable weak solution on
$\mathcal{O}^{\mu,z_{0}}$ with force $f^{\mu,z_{0}}$.
\end{lemma}

\begin{proof}
Throughout write $u^{\mu}=u^{\mu,z_{0}}$ and let
$T(y,s)=(x_{0}+\mu y,t_{0}+\mu^{2}s)$, so that $u^{\mu}=\mu\,(u\circ T)$, and
$T$ maps $\mathcal{O}^{\mu,z_{0}}$ onto $\mathcal{O}$ with Jacobian
$\mu^{5}$ in $(y,s)$.

\emph{(S1).} If $\Omega''\times J''$ has compact closure in
$\mathcal{O}^{\mu,z_{0}}$ then $T(\Omega''\times J'')$ has compact closure in
$\mathcal{O}$, and
\[
  \int_{\Omega''}\abs{u^{\mu}(y,s)}^{2}dy
  =\mu^{2}\mu^{-3}\int_{T_{x}(\Omega'')}\abs{u(\cdot,t_{0}+\mu^{2}s)}^{2}dx
  =\mu^{-1}\int_{T_{x}(\Omega'')}\abs{u}^{2}dx ,
\]
where $T_{x}(y)=x_{0}+\mu y$; by the same change of variables,
\[
\iint\abs{\nabla_{y}u^{\mu}}^{2}dy\,ds
 =\mu^{2}\mu^{2}\mu^{-5}\iint\abs{\nabla u}^{2}dz=\mu^{-1}\iint\abs{\nabla u}^{2}
,
\]
$\iint\abs{u^{\mu}}^{2}=\mu^{-3}\iint\abs{u}^{2}$,
\[
\iint\abs{p^{\mu}}^{3/2}=\mu^{3}\mu^{-5}\iint\abs{p}^{3/2}=\mu^{-2}\iint\abs{p}^{3/2}
\]

and
$\iint\abs{f^{\mu}}^{q}=\mu^{3q-5}\iint\abs{f}^{q}$. All are finite, so (S1)
holds for $(u^{\mu},p^{\mu})$.

\emph{(S2) and (S3).} Given $\psi\in C^{\infty}_{c}(\mathcal{O}^{\mu,z_{0}})$,
put $\widehat\psi=\psi\circ T^{-1}\in C_{c}^{\infty}(\mathcal{O})$. Then
$\nabla_{y}\psi=\mu(\nabla\widehat\psi)\circ T$ and
$\iint u^{\mu}\cdot\nabla_{y}\psi\,dy\,ds
 =\mu\cdot\mu\cdot\mu^{-5}\iint u\cdot\nabla\widehat\psi\,dz
 =\mu^{-3}\iint u\cdot\nabla\widehat\psi\,dz=0$, which is (S2). For (S3), let
$\varphi\in C_{c}^{\infty}(\mathcal{O}^{\mu,z_{0}};\R^{3})$ and put
$\widehat\varphi=\varphi\circ T^{-1}$. Using
$\partial_{s}\varphi=\mu^{2}(\partial_{t}\widehat\varphi)\circ T$,
$\nabla_{y}\varphi=\mu(\nabla\widehat\varphi)\circ T$, one computes the five
terms of \eqref{eq:weak-momentum} for $(u^{\mu},p^{\mu},f^{\mu},\varphi)$:
\[
\begin{aligned}
  \iint u^{\mu}\cdot\partial_{s}\varphi
    &=\mu\cdot\mu^{2}\cdot\mu^{-5}\iint u\cdot\partial_{t}\widehat\varphi
     =\mu^{-2}\iint u\cdot\partial_{t}\widehat\varphi, \\
  \iint u^{\mu}_{i}u^{\mu}_{j}\partial_{j}\varphi_{i}
    &=\mu^{2}\cdot\mu\cdot\mu^{-5}\iint u_{i}u_{j}\partial_{j}\widehat\varphi_{i}
     =\mu^{-2}\iint u_{i}u_{j}\partial_{j}\widehat\varphi_{i}, \\
  \iint\nabla_{y}u^{\mu}:\nabla_{y}\varphi
    &=\mu^{2}\cdot\mu\cdot\mu^{-5}\iint\nabla u:\nabla\widehat\varphi
     =\mu^{-2}\iint\nabla u:\nabla\widehat\varphi, \\
  \iint p^{\mu}\,\dv_{y}\varphi
    &=\mu^{2}\cdot\mu\cdot\mu^{-5}\iint p\,\dv\widehat\varphi
     =\mu^{-2}\iint p\,\dv\widehat\varphi, \\
  \iint f^{\mu}\cdot\varphi
    &=\mu^{3}\cdot\mu^{-5}\iint f\cdot\widehat\varphi
     =\mu^{-2}\iint f\cdot\widehat\varphi .
\end{aligned}
\]
Every term carries the same factor $\mu^{-2}$, so \eqref{eq:weak-momentum} for
$(u^{\mu},p^{\mu},f^{\mu})$ is $\mu^{-2}$ times \eqref{eq:weak-momentum} for
$(u,p,f)$ with test field $\widehat\varphi$, hence holds.

\emph{(S4).} Let $\phi\in C_{c}^{\infty}(\mathcal{O}^{\mu,z_{0}})$, $\phi\ge0$,
and $\widehat\phi=\phi\circ T^{-1}\ge0$. Then
$\partial_{s}\phi+\Delta_{y}\phi=\mu^{2}\bigl((\partial_{t}\widehat\phi+\Delta\widehat\phi)\circ T\bigr)$,
so
\[
\begin{aligned}
  \iint\abs{\nabla_{y}u^{\mu}}^{2}\phi
    &=\mu^{-1}\iint\abs{\nabla u}^{2}\widehat\phi,
  \\
  \iint\abs{u^{\mu}}^{2}(\partial_{s}\phi+\Delta_{y}\phi)
    &=\mu^{2}\mu^{2}\mu^{-5}\iint\abs{u}^{2}(\partial_{t}\widehat\phi+\Delta\widehat\phi),\\
  \iint\abs{u^{\mu}}^{2}u^{\mu}\!\cdot\!\nabla_{y}\phi
    &=\mu^{3}\mu\mu^{-5}\iint\abs{u}^{2}u\cdot\nabla\widehat\phi,
  \\
  \iint p^{\mu}u^{\mu}\!\cdot\!\nabla_{y}\phi
    &=\mu^{2}\mu\mu\mu^{-5}\iint p\,u\cdot\nabla\widehat\phi,\\
  \iint (f^{\mu}\cdot u^{\mu})\phi
    &=\mu^{3}\mu\mu^{-5}\iint(f\cdot u)\widehat\phi .
\end{aligned}
\]
Each of the five quantities equals $\mu^{-1}$ times the corresponding quantity
for $(u,p,f,\widehat\phi)$. Since $\mu^{-1}>0$, inequality \eqref{eq:lei} for
$(u,p,f)$ with test function $\widehat\phi$ is equivalent to \eqref{eq:lei} for
$(u^{\mu},p^{\mu},f^{\mu})$ with test function $\phi$.
\end{proof}

\begin{lemma}[Transformation of the scale-invariant quantities]\label{lem:scaling-quantities}
Let $(u,p,f)$, $z_{0}$, $\mu$ be as in Lemma~\ref{lem:scaling-invariance}, let
$z=(x,t)$, and let $\zeta=(\mu^{-1}(x-x_{0}),\mu^{-2}(t-t_{0}))$ be the
corresponding point for the rescaled solution. Then for every $r>0$ with
$\overline{\Cyl{\mu r}{z}}\subset\mathcal{O}$,
\begin{equation}\label{eq:scaling-quantities}
\begin{aligned}
  A(\zeta,r;u^{\mu,z_{0}})&=A(z,\mu r;u),
  &\quad B(\zeta,r;u^{\mu,z_{0}})&=B(z,\mu r;u),
  &\quad C(\zeta,r;u^{\mu,z_{0}})&=C(z,\mu r;u),\\
  D(\zeta,r;p^{\mu,z_{0}})&=D(z,\mu r;p),
  &\quad E_{q}(\zeta,r;f^{\mu,z_{0}})&=E_{q}(z,\mu r;f),
  &\quad \Psi(\zeta,r;u^{\mu,z_{0}})&=\Psi(z,\mu r;u),
\end{aligned}
\end{equation}
and likewise $\widetilde C(\zeta,r)=\widetilde C(z,\mu r)$,
$\widetilde D(\zeta,r)=\widetilde D(z,\mu r)$ and
$\varphi(\zeta,r)=\varphi(z,\mu r)$.
\end{lemma}

\begin{proof}
All six identities are the same computation. The map
$T(y,s)=(x_{0}+\mu y,t_{0}+\mu^{2}s)$ carries $\Cyl{r}{\zeta}$ onto $\Cyl{\mu r}{z}$
with $dz=\mu^{5}dy\,ds$ and carries $\Ball{r}{\cdot}$ onto $\Ball{\mu r}{\cdot}$
with $dx=\mu^{3}dy$. Hence, writing $v=u^{\mu,z_{0}}$,
\[\begin{aligned}
  &A(\zeta,r;v)=\frac1r\operatorname*{ess\,sup}_{s}\int_{\Ball{r}{\cdot}}\abs{v}^{2}dy
  \\ &=\frac1r\,\mu^{2}\mu^{-3}\operatorname*{ess\,sup}\int_{\Ball{\mu r}{\cdot}}\abs{u}^{2}dx
  \\ &=\frac{1}{\mu r}\operatorname*{ess\,sup}\int_{\Ball{\mu r}{\cdot}}\abs{u}^{2}dx
  =A(z,\mu r;u),
\end{aligned}\]
\[
  B(\zeta,r;v)=\frac1r\mu^{2}\mu^{2}\mu^{-5}\iint_{\Cyl{\mu r}{z}}\abs{\nabla u}^{2}
   =\frac{1}{\mu r}\iint_{\Cyl{\mu r}{z}}\abs{\nabla u}^{2}=B(z,\mu r;u),
\]
\[
  C(\zeta,r;v)=\frac1{r^{2}}\mu^{3}\mu^{-5}\iint_{\Cyl{\mu r}{z}}\abs{u}^{3}
   =\frac{1}{(\mu r)^{2}}\iint_{\Cyl{\mu r}{z}}\abs{u}^{3}=C(z,\mu r;u),
\]
\[
  D(\zeta,r;p^{\mu,z_{0}})=\frac1{r^{2}}\mu^{3}\mu^{-5}\iint_{\Cyl{\mu r}{z}}\abs{p}^{3/2}
   =D(z,\mu r;p),
\]
\[
  E_{q}(\zeta,r;f^{\mu,z_{0}})
   =r^{3-\frac5q}\bigl(\mu^{3q}\mu^{-5}\bigr)^{1/q}\Bigl(\iint_{\Cyl{\mu r}{z}}\abs{f}^{q}\Bigr)^{1/q}
   =(\mu r)^{3-\frac5q}\norm{f}_{L^{q}(\Cyl{\mu r}{z})}=E_{q}(z,\mu r;f).
\]
For $\Psi$, the space-time mean transforms as
$\stavg{v}{\Cyl{r}{\zeta}}=\mu\,\stavg{u}{\Cyl{\mu r}{z}}$ (the Jacobian cancels
against the normalization), so
$\Psi(\zeta,r;v)=r\mu\abs{\stavg{u}{\Cyl{\mu r}{z}}}=\Psi(z,\mu r;u)$. The same
transformation of the mean, together with the computation for $C$, gives
$\widetilde C(\zeta,r)=\widetilde C(z,\mu r)$; and since
$\spavg{p^{\mu,z_{0}}}{\Ball{r}{\cdot}}(s)=\mu^{2}\spavg{p}{\Ball{\mu r}{\cdot}}(t_{0}+\mu^{2}s)$,
the computation for $D$ gives $\widetilde D(\zeta,r)=\widetilde D(z,\mu r)$.
Adding the appropriate powers yields the statement for $\varphi$.
\end{proof}

\subsection{Interpolation on balls and cylinders}\label{ss:interpolation}

We use the Sobolev and Poincar\'e inequalities on a ball in the following
scale-explicit form; see External Input~\ref{ext:sobolev-ball} and
External Input~\ref{ext:poincare}.

\begin{lemma}[Sobolev and Poincar\'e on a ball]\label{lem:sobolev-poincare-ball}
There is an absolute constant $C_{5}$ such that for every $r>0$, every
$x\in\R^{3}$ and every $v\in H^{1}(\Ball{r}{x})$,
\begin{align}
  \bigl\|v-\spavg{v}{\Ball{r}{x}}\bigr\|_{L^{6}(\Ball{r}{x})}
    &\le C_{5}\,\norm{\nabla v}_{L^{2}(\Ball{r}{x})},
    \label{eq:sobolev-poincare-6}\\
  \norm{v}_{L^{6}(\Ball{r}{x})}
    &\le C_{5}\Bigl(\norm{\nabla v}_{L^{2}(\Ball{r}{x})}
       +r^{-1}\norm{v}_{L^{2}(\Ball{r}{x})}\Bigr),
    \label{eq:sobolev-ball}
\end{align}
and, for every $g\in W^{1,1}(\Ball{r}{x})$,
\begin{equation}\label{eq:poincare-L1}
  \bigl\|g-\spavg{g}{\Ball{r}{x}}\bigr\|_{L^{1}(\Ball{r}{x})}
  \le C_{5}\,r\,\norm{\nabla g}_{L^{1}(\Ball{r}{x})} .
\end{equation}
\end{lemma}

\begin{proof}
Inequality \eqref{eq:sobolev-poincare-6} for $r=1$, $x=0$ is
External Input~\ref{ext:poincare}(ii); the general case follows from the change
of variables $v(y)=\psi((y-x)/r)$, under which
$\norm{v-\spavg{v}{\Ball{r}{x}}}_{L^{6}(\Ball{r}{x})}
 =r^{1/2}\norm{\psi-\spavg{\psi}{B_{1}}}_{L^{6}(B_{1})}$ and
$\norm{\nabla v}_{L^{2}(\Ball{r}{x})}=r^{1/2}\norm{\nabla\psi}_{L^{2}(B_{1})}$,
so that both sides carry the same power $r^{1/2}$ and the constant is unchanged.
Inequality \eqref{eq:poincare-L1} is External Input~\ref{ext:poincare}(i)
rescaled in the same way: both sides carry the factor $r^{3}$ after the
substitution, once the factor $r$ on the right is accounted for. For
\eqref{eq:sobolev-ball}, write $v=(v-\spavg{v}{\Ball{r}{x}})+\spavg{v}{\Ball{r}{x}}$
and use \eqref{eq:sobolev-poincare-6} together with
\[\begin{aligned}
  &\bigl\|\spavg{v}{\Ball{r}{x}}\bigr\|_{L^{6}(\Ball{r}{x})}
  =\abs{\Ball{r}{x}}^{1/6}\bigl|\spavg{v}{\Ball{r}{x}}\bigr|
  \\ &\le\abs{\Ball{r}{x}}^{1/6}\abs{\Ball{r}{x}}^{-1/2}\norm{v}_{L^{2}(\Ball{r}{x})}
  =\bigl(\tfrac{4\pi}{3}\bigr)^{-1/3}r^{-1}\norm{v}_{L^{2}(\Ball{r}{x})},
\end{aligned}\]
where the middle step is the Cauchy--Schwarz inequality. Enlarging the constant
to cover all three inequalities gives $C_{5}$.
\end{proof}

\begin{lemma}[Interpolation on a ball]\label{lem:interpolation-ball}
Let $2\le \ell\le6$ and set
\begin{equation}\label{eq:a-exponent}
  a=a(\ell)=\tfrac34(\ell-2)\in[0,3],
  \qquad\text{so that}\qquad
  \tfrac \ell2-a=\tfrac{6-\ell}{4}\in[0,1].
\end{equation}
There is an absolute constant $C_{6}$ such that for every $r>0$, every
$x\in\R^{3}$ and every $v\in H^{1}(\Ball{r}{x})$,
\begin{equation}\label{eq:interpolation-ball}
  \int_{\Ball{r}{x}}\abs{v}^{\ell}\,dy
  \;\le\;C_{6}\Bigl(\int_{\Ball{r}{x}}\abs{\nabla v}^{2}\Bigr)^{a}
      \Bigl(\int_{\Ball{r}{x}}\abs{v}^{2}\Bigr)^{\frac \ell2-a}
     \;+\;\frac{C_{6}}{r^{2a}}\Bigl(\int_{\Ball{r}{x}}\abs{v}^{2}\Bigr)^{\frac \ell2}.
\end{equation}
\end{lemma}

\begin{proof}
Write $B=\Ball{r}{x}$. Let $\theta\in[0,1]$ be defined by
$\tfrac1\ell=\tfrac{1-\theta}{2}+\tfrac{\theta}{6}$, that is
$\theta=3\bigl(\tfrac12-\tfrac1\ell\bigr)=\tfrac{3(\ell-2)}{2\ell}$; this is in $[0,1]$
precisely because $2\le \ell\le6$. The H\"older interpolation inequality between
$L^{2}$ and $L^{6}$ gives
$\norm{v}_{L^{\ell}(B)}\le\norm{v}_{L^{2}(B)}^{1-\theta}\norm{v}_{L^{6}(B)}^{\theta}$,
hence
\begin{equation}\label{eq:interp-step1}
  \int_{B}\abs{v}^{\ell}
  =\norm{v}_{L^{\ell}(B)}^{\ell}
  \le\norm{v}_{L^{2}(B)}^{\ell(1-\theta)}\norm{v}_{L^{6}(B)}^{\ell\theta} .
\end{equation}
The two exponents are
\[
  \ell\theta=\tfrac{3(\ell-2)}{2}=2a,
  \qquad
  \ell(1-\theta)=\ell-\tfrac{3(\ell-2)}{2}=\tfrac{6-\ell}{2}=2\Bigl(\tfrac \ell2-a\Bigr).
\]
Inserting \eqref{eq:sobolev-ball} into \eqref{eq:interp-step1} and using
$(s_{1}+s_{2})^{2a}\le2^{2a}(s_{1}^{2a}+s_{2}^{2a})\le 64\,(s_{1}^{2a}+s_{2}^{2a})$
for $s_{1},s_{2}\ge0$ and $0\le a\le3$,
\[\begin{aligned}
  &\int_{B}\abs{v}^{\ell}
  \le 64\,C_{5}^{2a}\norm{v}_{L^{2}(B)}^{2(\frac \ell2-a)}
     \Bigl(\norm{\nabla v}_{L^{2}(B)}^{2a}+r^{-2a}\norm{v}_{L^{2}(B)}^{2a}\Bigr)
  \\ &=64\,C_{5}^{2a}\Bigl[
     \Bigl(\int_{B}\abs{\nabla v}^{2}\Bigr)^{a}\Bigl(\int_{B}\abs{v}^{2}\Bigr)^{\frac \ell2-a}
     +r^{-2a}\Bigl(\int_{B}\abs{v}^{2}\Bigr)^{\frac \ell2}\Bigr].
\end{aligned}\]
Since $0\le a\le3$, the factor $C_{5}^{2a}$ is bounded by $\max\{1,C_{5}^{6}\}$;
take $C_{6}=64\max\{1,C_{5}^{6}\}$, which is absolute.
\end{proof}

\begin{remark}[Two instances, and a defect in the source]\label{rem:interp-instances}
The two cases used below are
\begin{equation}\label{eq:interp-q3}
  \ell=3,\quad a=\tfrac34:\qquad
  \int_{\Ball{r}{x}}\abs{v}^{3}
   \le C_{6}\Bigl(\int\abs{\nabla v}^{2}\Bigr)^{3/4}\Bigl(\int\abs{v}^{2}\Bigr)^{3/4}
     +\frac{C_{6}}{r^{3/2}}\Bigl(\int\abs{v}^{2}\Bigr)^{3/2},
\end{equation}
\begin{equation}\label{eq:interp-q103}
  \ell=\tfrac{10}{3},\quad a=1:\qquad
  \int_{\Ball{r}{x}}\abs{v}^{10/3}
   \le C_{6}\Bigl(\int\abs{\nabla v}^{2}\Bigr)\Bigl(\int\abs{v}^{2}\Bigr)^{2/3}
     +\frac{C_{6}}{r^{2}}\Bigl(\int\abs{v}^{2}\Bigr)^{5/3}.
\end{equation}
Integrating \eqref{eq:interp-q103} in time over $(t-r^{2},t)$ and using (S1)
gives the embedding $L^{\infty}L^{2}\cap L^{2}H^{1}\subset L^{10/3}$ of
External Input~\ref{ext:sobolev-ball}(iv).

The inequality \cite[(2.1)]{Lin1998} is printed with the exponents $a$ and
$\ell/2-a$ interchanged, i.e.\ with
$\bigl(\int\abs{\nabla v}^{2}\bigr)^{\ell/2-a}\bigl(\int\abs{v}^{2}\bigr)^{a}$ in
place of the first term of \eqref{eq:interpolation-ball}. That form agrees with
\eqref{eq:interpolation-ball} when $\ell=3$, the only case Lin uses, but it is
false for every $\ell\in(3,6]$. Indeed, take $r=1$, fix
$\psi\in C_{c}^{\infty}(B_{1})$ with $\psi\not\equiv0$, and set
$v_{\varepsilon}(y)=\varepsilon^{-3/2}\psi(y/\varepsilon)$ for
$0<\varepsilon<1$. Then
\[
  \int_{B_{1}}\abs{v_{\varepsilon}}^{2}=\int_{B_{1}}\abs{\psi}^{2},
  \qquad
  \int_{B_{1}}\abs{\nabla v_{\varepsilon}}^{2}
    =\varepsilon^{-2}\int_{B_{1}}\abs{\nabla\psi}^{2},
  \qquad
  \int_{B_{1}}\abs{v_{\varepsilon}}^{\ell}
    =\varepsilon^{3-\frac{3\ell}{2}}\int_{B_{1}}\abs{\psi}^{\ell},
\]
so the printed right side is
$O\bigl(\varepsilon^{-2(\ell/2-a)}\bigr)+O(1)=O\bigl(\varepsilon^{(\ell-6)/2}\bigr)+O(1)$
while the left side is of exact order $\varepsilon^{3-3\ell/2}$; since
$3-\tfrac{3\ell}{2}<\min\{\tfrac{\ell-6}{2},0\}$ for $\ell>3$, the printed inequality
fails as $\varepsilon\to0$. With the exponents as in
\eqref{eq:interpolation-ball} the first term is of order
$\varepsilon^{-2a}=\varepsilon^{3-3\ell/2}$, matching the left side exactly.
\end{remark}

\begin{lemma}[Interpolation on cylinders]\label{lem:interpolation-cylinder}
There is an absolute constant $C_{7}$ with the following property. Let $(u,p)$
satisfy (S1), let $z=(x,t)$ and let $0<r\le\rho$ with
$\overline{\Cyl{\rho}{z}}\subset\mathcal{O}$. Then
\begin{equation}\label{eq:lemma-C-decay}
  C(z,r)\;\le\;C_{7}\Bigl[\Bigl(\frac r\rho\Bigr)^{3}A(z,\rho)^{3/2}
     +\Bigl(\frac\rho r\Bigr)^{3}A(z,\rho)^{3/4}B(z,\rho)^{3/4}\Bigr].
\end{equation}
Equivalently, in Kukavica's variables,
$\gamma(z,r)^{3}\le C_{7}\bigl[(r/\rho)^{3}\alpha(z,\rho)^{3}
 +(\rho/r)^{3}\alpha(z,\rho)^{3/2}\beta(z,\rho)^{3/2}\bigr]$.
\end{lemma}

\begin{proof}
Write $A=A(z,\rho)$, $B=B(z,\rho)$, $B_{r}=\Ball{r}{x}$, $B_{\rho}=\Ball{\rho}{x}$,
and for $s\in(t-\rho^{2},t)$ put
\[
  Y(s)=\int_{B_{r}}\abs{u(y,s)}^{2}dy,
  \qquad
  G(s)=\int_{B_{\rho}}\abs{\nabla u(y,s)}^{2}dy .
\]
All statements below hold for almost every such $s$; we suppress ``a.e.''.

\emph{Step 1: a bound for $Y(s)$.} Split
\[
  \int_{B_{r}}\abs{u}^{2}
  =\int_{B_{r}}\Bigl(\abs{u}^{2}-\spavg{\abs{u}^{2}}{B_{\rho}}(s)\Bigr)
   +\abs{B_{r}}\,\spavg{\abs{u}^{2}}{B_{\rho}}(s)
  \le\int_{B_{\rho}}\Bigl|\abs{u}^{2}-\spavg{\abs{u}^{2}}{B_{\rho}}(s)\Bigr|
   +\Bigl(\frac r\rho\Bigr)^{3}\int_{B_{\rho}}\abs{u}^{2},
\]
where we used $B_{r}\subset B_{\rho}$ in the first integral and
$\abs{B_{r}}/\abs{B_{\rho}}=(r/\rho)^{3}$ in the second. Apply
\eqref{eq:poincare-L1} to $g=\abs{u(\cdot,s)}^{2}$, which lies in
$W^{1,1}(B_{\rho})$ with $\abs{\nabla g}\le2\abs{u}\abs{\nabla u}$ because
$u(\cdot,s)\in H^{1}(B_{\rho})$ and $u(\cdot,s)\in L^{4}(B_{\rho})$ by
\eqref{eq:interpolation-ball} with $q=4$:
\[
  \int_{B_{\rho}}\Bigl|\abs{u}^{2}-\spavg{\abs{u}^{2}}{B_{\rho}}\Bigr|
  \le 2C_{5}\rho\int_{B_{\rho}}\abs{u}\abs{\nabla u}
  \le 2C_{5}\rho\,\norm{u(\cdot,s)}_{L^{2}(B_{\rho})}G(s)^{1/2} .
\]
Since $\norm{u(\cdot,s)}_{L^{2}(B_{\rho})}^{2}\le\rho A$ by definition of $A$,
\begin{equation}\label{eq:Y-bound}
  Y(s)\;\le\;2C_{5}\,\rho^{3/2}A^{1/2}G(s)^{1/2}
     \;+\;\frac{r^{3}}{\rho^{2}}\,A .
\end{equation}

\emph{Step 2: the pointwise-in-time $L^{3}$ bound.} Apply \eqref{eq:interp-q3}
on $B_{r}$:
\[
  \int_{B_{r}}\abs{u}^{3}
  \le C_{6}\Bigl(\int_{B_{r}}\abs{\nabla u}^{2}\Bigr)^{3/4}Y(s)^{3/4}
    +\frac{C_{6}}{r^{3/2}}Y(s)^{3/2} .
\]
In the first term use
\[
 \int_{B_{r}}\abs{\nabla u}^{2}\le G(s),
 \qquad Y(s)\le\int_{B_{\rho}}\abs{u}^{2}\le\rho A,
\]
obtaining
$C_{6}\rho^{3/4}A^{3/4}G(s)^{3/4}$. In the second term insert
\eqref{eq:Y-bound} and use
$(s_{1}+s_{2})^{3/2}\le2^{3/2}(s_{1}^{3/2}+s_{2}^{3/2})$:
\[
  \frac{C_{6}}{r^{3/2}}Y(s)^{3/2}
  \le\frac{2^{3/2}C_{6}}{r^{3/2}}
    \Bigl[(2C_{5})^{3/2}\rho^{9/4}A^{3/4}G(s)^{3/4}
      +\frac{r^{9/2}}{\rho^{3}}A^{3/2}\Bigr].
\]
Hence, with $C_{8}=C_{6}\bigl(1+2^{3/2}(2C_{5})^{3/2}+2^{3/2}\bigr)$,
\begin{equation}\label{eq:L3-pointwise}
  \int_{B_{r}}\abs{u(\cdot,s)}^{3}
  \;\le\;C_{8}\Bigl(\rho^{3/4}+\frac{\rho^{9/4}}{r^{3/2}}\Bigr)A^{3/4}G(s)^{3/4}
    \;+\;C_{8}\,\frac{r^{3}}{\rho^{3}}\,A^{3/2}.
\end{equation}

\emph{Step 3: integration in time.} Integrate \eqref{eq:L3-pointwise} over
$s\in(t-r^{2},t)$. For the first term use H\"older with exponents
$\tfrac43$ and $4$ and then $\int_{t-\rho^{2}}^{t}G(s)\,ds=\rho B$:
\[
  \int_{t-r^{2}}^{t}G(s)^{3/4}ds
  \le\bigl(r^{2}\bigr)^{1/4}\Bigl(\int_{t-r^{2}}^{t}G(s)\,ds\Bigr)^{3/4}
  \le r^{1/2}(\rho B)^{3/4} .
\]
The second term contributes $C_{8}r^{5}\rho^{-3}A^{3/2}$. Therefore
\[\begin{aligned}
  &\iint_{\Cyl{r}{z}}\abs{u}^{3}
  \le C_{8}\Bigl(\rho^{3/4}+\frac{\rho^{9/4}}{r^{3/2}}\Bigr)
       A^{3/4}\,r^{1/2}\rho^{3/4}B^{3/4}
  \\ &+C_{8}\frac{r^{5}}{\rho^{3}}A^{3/2}
  \\ &= C_{8}\Bigl(\rho^{3/2}r^{1/2}+\frac{\rho^{3}}{r}\Bigr)A^{3/4}B^{3/4}
    +C_{8}\frac{r^{5}}{\rho^{3}}A^{3/2}.
\end{aligned}\]
Dividing by $r^{2}$,
\[
  C(z,r)\le C_{8}\Bigl[\Bigl(\frac\rho r\Bigr)^{3/2}
     +\Bigl(\frac\rho r\Bigr)^{3}\Bigr]A^{3/4}B^{3/4}
    +C_{8}\Bigl(\frac r\rho\Bigr)^{3}A^{3/2},
\]
and since $r\le\rho$ gives $(\rho/r)^{3/2}\le(\rho/r)^{3}$, the claim follows
with $C_{7}=2C_{8}$.
\end{proof}

\begin{corollary}[Gagliardo--Nirenberg in Kukavica's variables]\label{cor:gagliardo}
With $C_{7}$ as in Lemma~\ref{lem:interpolation-cylinder}, for every $z$ and
every $\rho>0$ with $\overline{\Cyl{\rho}{z}}\subset\mathcal{O}$,
\begin{equation}\label{eq:gagliardo}
  \gamma(z,\rho)\;\le\;C_{9}\,\alpha(z,\rho)^{1/2}\beta(z,\rho)^{1/2}
    \;+\;C_{9}\,\alpha(z,\rho),
  \qquad C_{9}=(2C_{7})^{1/3},
\end{equation}
The same cylinder interpolation estimate also gives
\[
C(z,\rho)\le 2C_{7}\bigl[A(z,\rho)^{3/2}+A(z,\rho)^{3/4}B(z,\rho)^{3/4}\bigr]
.
\]
\end{corollary}

\begin{proof}
Take $r=\rho$ in \eqref{eq:lemma-C-decay}:
$C(z,\rho)\le C_{7}[A^{3/2}+A^{3/4}B^{3/4}]$. Using
$\gamma=C^{1/3}$, $\alpha=A^{1/2}$, $\beta=B^{1/2}$ from
Lemma~\ref{lem:dictionary} and
$(s_{1}+s_{2})^{1/3}\le s_{1}^{1/3}+s_{2}^{1/3}$,
\[
  \gamma(z,\rho)=C(z,\rho)^{1/3}
  \le C_{7}^{1/3}\bigl(A^{1/2}+A^{1/4}B^{1/4}\bigr)
  = C_{7}^{1/3}\bigl(\alpha+\alpha^{1/2}\beta^{1/2}\bigr).
\]
Since $C_{7}^{1/3}\le(2C_{7})^{1/3}=C_{9}$, \eqref{eq:gagliardo} follows.
\end{proof}

\begin{lemma}[Monotonicity in the radius]\label{lem:monotonicity}
Let $z\in\mathcal{O}$ and $0<\rho_{1}<\rho_{2}$ with
$\overline{\Cyl{\rho_{2}}{z}}\subset\mathcal{O}$. Then
\begin{equation}\label{eq:monotonicity}
\begin{aligned}
  \alpha(z,\rho_{1})&\le\Bigl(\frac{\rho_{2}}{\rho_{1}}\Bigr)^{1/2}\alpha(z,\rho_{2}),
  &\qquad
  \beta(z,\rho_{1})&\le\Bigl(\frac{\rho_{2}}{\rho_{1}}\Bigr)^{1/2}\beta(z,\rho_{2}),\\
  \gamma(z,\rho_{1})&\le\Bigl(\frac{\rho_{2}}{\rho_{1}}\Bigr)^{2/3}\gamma(z,\rho_{2}),
  &\qquad
  \delta(z,\rho_{1})^{2}&\le\Bigl(\frac{\rho_{2}}{\rho_{1}}\Bigr)^{4/3}\delta(z,\rho_{2})^{2},
\end{aligned}
\end{equation}
and $\lambda(z,\rho_{1})\le(\rho_{1}/\rho_{2})^{\sigma}\lambda(z,\rho_{2})$.
\end{lemma}

\begin{proof}
All five follow from $\Cyl{\rho_{1}}{z}\subset\Cyl{\rho_{2}}{z}$,
$\Ball{\rho_{1}}{x}\subset\Ball{\rho_{2}}{x}$ and
$(t-\rho_{1}^{2},t)\subset(t-\rho_{2}^{2},t)$, which make every integral in
\eqref{eq:alpha-beta} monotone in the radius, combined with the explicit power of
$r$ in the prefactor. For instance
\[
\begin{aligned}
&\alpha(z,\rho_{1})^{2}
\\ &=\rho_{1}^{-1}\operatorname*{ess\,sup}_{(t-\rho_{1}^{2},t)}\int_{\Ball{\rho_{1}}{x}}\abs{u}^{2}
\\ &\le\rho_{1}^{-1}\operatorname*{ess\,sup}_{(t-\rho_{2}^{2},t)}\int_{\Ball{\rho_{2}}{x}}\abs{u}^{2}
\\ &=(\rho_{2}/\rho_{1})\alpha(z,\rho_{2})^{2}
\end{aligned}
;
\] enlarging the domain of
integration in the same way in the three remaining integrals gives
$\beta(z,\rho_{1})^{2}\le(\rho_{2}/\rho_{1})\beta(z,\rho_{2})^{2}$,
$\gamma(z,\rho_{1})^{3}\le(\rho_{2}/\rho_{1})^{2}\gamma(z,\rho_{2})^{3}$ and
$\delta(z,\rho_{1})^{3}\le(\rho_{2}/\rho_{1})^{2}\delta(z,\rho_{2})^{3}$;
raising the last two to the powers $1/3$ and $2/3$ gives the stated forms. For
$\lambda$, $\norm{f}_{L^{q}(\Cyl{\rho_{1}}{z})}\le\norm{f}_{L^{q}(\Cyl{\rho_{2}}{z})}$
gives
\[
\lambda(z,\rho_{1})\le\rho_{1}^{\sigma}\norm{f}_{L^{q}(\Cyl{\rho_{2}}{z})}
=(\rho_{1}/\rho_{2})^{\sigma}\lambda(z,\rho_{2})
.
\]
\end{proof}

\begin{remark}\label{rem:kukavica-monotonicity}
Kukavica \cite[(2.23) and the line above it]{Kukavica2008} states the fourth
inequality of \eqref{eq:monotonicity} as
$\delta(\rho_{1})\le(\rho_{2}/\rho_{1})^{4/3}\delta(\rho_{2})$. The exponent
$4/3$ belongs to $\delta^{2}$, not to $\delta$: the correct statements are
$\delta(\rho_{1})\le(\rho_{2}/\rho_{1})^{2/3}\delta(\rho_{2})$ and
$\delta(\rho_{1})^{2}\le(\rho_{2}/\rho_{1})^{4/3}\delta(\rho_{2})^{2}$. Since
$\rho_{2}/\rho_{1}\ge1$ the printed inequality is weaker than the truth, so the
argument of \cite{Kukavica2008} is unaffected.
\end{remark}

\section{The pressure}\label{sec:pressure}

\begin{remark}[Operator and source conventions]
The pressure identities below use weak derivatives and completed $L^p$
Riesz operators. One may first prove the Newtonian derivative adjoint
identity on compact smooth tests by heat-kernel subordination and then
pass to nonsmooth compact sources by density. Exterior kernel
representations are justified on sets separated from the source support.
Local potential integrability and growth at infinity identify the
harmonic residual by Liouville's theorem; they do not turn a singular
Hessian kernel into an everywhere absolutely integrable convolution.

The singly centred tensor $-u_i(u_j-c_j)$ used in the pressure-decay
estimate is distinct from the doubly centred tensor
$-(u_i-c_i)(u_j-c_j)$ used in Proposition~\ref{prop:lin34}.
For the latter, subtract the two distributional pairings and use
incompressibility and Hessian symmetry to cancel the constant terms.
The nine scalar Riesz estimates give the factor $9C_{11}$.
A source obtained by differentiating the velocity tensor contains no
force term: the force potentials $p_7,p_8$ are added separately.
\end{remark}

Throughout this section $(u,p)$ is a suitable weak solution on $\mathcal{O}$
with force $f\in L^{q}_{loc}$, $q>5/2$, and $z_{0}=(x_{0},t_{0})\in\mathcal{O}$
and $\rho>0$ are such that $\overline{\Cyl{\rho}{z_{0}}}\subset\mathcal{O}$. We write
$B_{\rho}=\Ball{\rho}{x_{0}}$, $J_{\rho}=(t_{0}-\rho^{2},t_{0})$, and
\begin{equation}\label{eq:kappa}
  \kappa=\frac r\rho\in\bigl(0,\tfrac12\bigr]
\end{equation}
for the two radii $0<r\le\rho/2$ appearing in the estimates. The symbol
$\kappa$ keeps this meaning up to and including Lemma~\ref{lem:theta-decay},
where it is still a free parameter; from Convention~\ref{conv:kappa} onwards it
denotes the \emph{fixed} absolute constant defined there, and the functional
$\theta$ of \eqref{eq:theta} is understood with that fixed value.

\subsection{The spatial Poisson equation for the pressure}\label{ss:poisson}

\begin{lemma}[$\Delta p$ at almost every time]\label{lem:delta-p}
For almost every $t\in J_{\rho}$ and every $\psi\in C_{c}^{\infty}(B_{\rho})$,
\begin{equation}\label{eq:delta-p}
  \int_{B_{\rho}}p(y,t)\,\Delta\psi(y)\,dy
  \;=\;\int_{B_{\rho}}U_{ij}(y,t)\,\partial_{i}\partial_{j}\psi(y)\,dy
     \;-\;\int_{B_{\rho}}f(y,t)\cdot\nabla\psi(y)\,dy,
\end{equation}
where
\begin{equation}\label{eq:Uij}
  U_{ij}(y,t)\;=\;-\,u_{i}(y,t)\bigl(u_{j}(y,t)-\spavg{u_{j}}{B_{\rho}}(t)\bigr).
\end{equation}
Equivalently, $\Delta p=\partial_{i}\partial_{j}U_{ij}+\dv f$ in
$\mathcal{D}'(B_{\rho})$ for a.e.\ $t\in J_{\rho}$.
\end{lemma}

\begin{proof}
\emph{Step 1: a null set that works for all $\psi$.} Let
$\mathcal{S}\subset C_{c}^{\infty}(B_{\rho})$ be a countable set that is dense in
$C_{c}^{\infty}(B_{\rho})$ for the $C^{2}$ norm together with uniform control of
supports, for instance the set of finite rational-coefficient combinations of a
fixed countable family of bump functions. It suffices to prove
\eqref{eq:delta-p} for each fixed $\psi\in\mathcal{S}$ outside a null set
$N_{\psi}$ depending on $\psi$, since then \eqref{eq:delta-p} holds for all
$\psi\in\mathcal{S}$ outside $N=\bigcup_{\psi\in\mathcal{S}}N_{\psi}$, and both
sides of \eqref{eq:delta-p} are continuous in $\psi$ for the $C^{2}$ norm on a
fixed compact support (each integrand is in $L^{1}(B_{\rho})$ for a.e.\ $t$).

\emph{Step 2: the mean of $u$ against a gradient.} Applying (S2) to
$\psi(x)\theta(t)$ with $\theta\in C_{c}^{\infty}(J_{\rho})$ gives
$\int_{J_{\rho}}\theta(t)\bigl(\int_{B_{\rho}}u(\cdot,t)\cdot\nabla\psi\bigr)dt=0$
for all such $\theta$, hence
\begin{equation}\label{eq:divfree-ae}
  \int_{B_{\rho}}u(y,t)\cdot\nabla\psi(y)\,dy=0
  \qquad\text{for a.e.\ }t\in J_{\rho} .
\end{equation}

\emph{Step 3: testing the momentum equation with a gradient.} Fix
$\psi\in\mathcal{S}$ and $\theta\in C_{c}^{\infty}(J_{\rho})$, and apply (S3)
with $\varphi(x,t)=\theta(t)\nabla\psi(x)$, which lies in
$C_{c}^{\infty}(\mathcal{O};\R^{3})$. The five terms of
\eqref{eq:weak-momentum} are:
\[
  \iint u\cdot\partial_{t}\varphi
  =\int\theta'(t)\Bigl(\int_{B_{\rho}}u\cdot\nabla\psi\Bigr)dt=0
\]
by \eqref{eq:divfree-ae};
\[
  \iint\nabla u:\nabla\varphi
  =\int\theta(t)\Bigl(\int_{B_{\rho}}\partial_{j}u_{i}\,\partial_{j}\partial_{i}\psi\Bigr)dt
  =\int\theta(t)\Bigl(-\int_{B_{\rho}}u_{i}\,\partial_{i}(\Delta\psi)\Bigr)dt=0 ,
\]
where the middle equality is an integration by parts in $x_{j}$, licit because
$u(\cdot,t)\in H^{1}(B_{\rho})$ and $\psi\in C_{c}^{\infty}(B_{\rho})$, and the
last equality is \eqref{eq:divfree-ae} applied to $\Delta\psi\in
C_{c}^{\infty}(B_{\rho})$; and
\[
\begin{aligned}
  &\iint u_{i}u_{j}\partial_{j}\varphi_{i}
   =\int\theta\int u_{i}u_{j}\,\partial_{i}\partial_{j}\psi,\\
  &\iint p\,\dv\varphi=\int\theta\int p\,\Delta\psi,
  \qquad
  \iint f\cdot\varphi=\int\theta\int f\cdot\nabla\psi .
\end{aligned}
\]
Thus \eqref{eq:weak-momentum} reads
$\int\theta(t)\bigl[-\int u_{i}u_{j}\partial_{i}\partial_{j}\psi
 -\int p\,\Delta\psi-\int f\cdot\nabla\psi\bigr]dt=0$ for every
$\theta\in C_{c}^{\infty}(J_{\rho})$, whence for a.e.\ $t$
\begin{equation}\label{eq:delta-p-raw}
  \int_{B_{\rho}}p\,\Delta\psi
  =-\int_{B_{\rho}}u_{i}u_{j}\,\partial_{i}\partial_{j}\psi
   -\int_{B_{\rho}}f\cdot\nabla\psi .
\end{equation}

\emph{Step 4: replacing $u_{j}$ by $u_{j}-\spavg{u_{j}}{B_{\rho}}(t)$.} Let
$c\in\R^{3}$ be a constant vector. Then, for a.e.\ $t$,
\[
  \int_{B_{\rho}}u_{i}c_{j}\,\partial_{i}\partial_{j}\psi
  =c_{j}\int_{B_{\rho}}u_{i}\,\partial_{i}(\partial_{j}\psi)
  =0
\]
by \eqref{eq:divfree-ae} applied to $\partial_{j}\psi\in C_{c}^{\infty}(B_{\rho})$.
Taking $c=\spavg{u}{B_{\rho}}(t)$, which is a constant vector for each fixed $t$,
turns $-\int u_{i}u_{j}\partial_{i}\partial_{j}\psi$ into
$-\int u_{i}(u_{j}-c_{j})\partial_{i}\partial_{j}\psi
 =\int U_{ij}\partial_{i}\partial_{j}\psi$. Substituting into
\eqref{eq:delta-p-raw} gives \eqref{eq:delta-p}.
\end{proof}

\subsection{The cut-off and the decomposition}\label{ss:pressure-decomposition}

\begin{lemma}[The spatial cut-off]\label{lem:cutoff}
There is a family of constants $C_{10}(k)$, $k\in\N_{0}$, absolute in the sense
that they depend on $k$ only, with the following property. For every $x_{0}$ and
every $\rho>0$ there is $\eta=\eta_{x_{0},\rho}\in C_{c}^{\infty}(\R^{3})$ with
\begin{equation}\label{eq:cutoff}
  0\le\eta\le1,\quad
  \eta\equiv1\ \text{on}\ \Ball{13\rho/20}{x_{0}},\quad
  \supp\eta\subset\Ball{3\rho/4}{x_{0}},\quad
  \abs{\partial^{a}\eta}\le C_{10}(\abs{a})\rho^{-\abs{a}}\ \text{on }\R^{3}.
\end{equation}
Moreover $\nabla\eta$ and all higher derivatives of $\eta$ vanish outside the
annulus
\begin{equation}\label{eq:annulus}
  \mathcal{A}_{\rho}=\Ball{3\rho/4}{x_{0}}\setminus\overline{\Ball{13\rho/20}{x_{0}}},
\end{equation}
and for every $x\in\Ball{r}{x_{0}}$ with $r\le\rho/2$ and every
$y\in\mathcal{A}_{\rho}$ one has $\abs{x-y}\ge\tfrac{3\rho}{20}$.
\end{lemma}

\begin{proof}
Let $\zeta\in C_{c}^{\infty}(B_{1})$ be a fixed nonnegative radial mollifier with
$\int_{\R^{3}}\zeta=1$, and for $\varepsilon>0$ put
$\zeta_{\varepsilon}(y)=\varepsilon^{-3}\zeta(y/\varepsilon)$.
Since $\supp\zeta$ is compact in $B_1$, choose $a_*<1$ with
$\supp\zeta\subset\overline{B_{a_*}}$. The convolution below then has
support in $\overline{\Ball{(7/10+a_*/20)\rho}{x_0}}$, which is a compact
subset of $\Ball{3\rho/4}{x_0}$. Set
$\eta=\mathbf{1}_{\Ball{7\rho/10}{x_{0}}}*\zeta_{\rho/20}$. Then
$0\le\eta\le1$; if $\abs{y-x_{0}}\le\tfrac{7\rho}{10}-\tfrac{\rho}{20}
=\tfrac{13\rho}{20}$ then the whole ball $\Ball{\rho/20}{y}$ lies in
$\Ball{7\rho/10}{x_{0}}$ and $\eta(y)=1$; if
$\abs{y-x_{0}}\ge\tfrac{7\rho}{10}+\tfrac{\rho}{20}=\tfrac{3\rho}{4}$ then
$\Ball{\rho/20}{y}\cap\Ball{7\rho/10}{x_{0}}=\emptyset$ and $\eta(y)=0$. Also
$\partial^{a}\eta=\mathbf{1}_{\Ball{7\rho/10}{x_{0}}}*\partial^{a}\zeta_{\rho/20}$,
so
$\abs{\partial^{a}\eta}\le\norm{\partial^{a}\zeta_{\rho/20}}_{L^{1}}
 =(20/\rho)^{\abs{a}}\norm{\partial^{a}\zeta}_{L^{1}(\R^{3})}$, giving
\eqref{eq:cutoff} with
$C_{10}(k)=20^{k}\max_{\abs{a}=k}\norm{\partial^{a}\zeta}_{L^{1}(\R^{3})}$. Since
$\eta$ is constant on $\Ball{13\rho/20}{x_{0}}$ and on the complement of
$\overline{\Ball{3\rho/4}{x_{0}}}$, all its derivatives of order $\ge1$ vanish
off $\mathcal{A}_{\rho}$. Finally if $\abs{x-x_{0}}<r\le\rho/2$ and
$\abs{y-x_{0}}\ge\tfrac{13\rho}{20}$ then
$\abs{x-y}\ge\tfrac{13\rho}{20}-\tfrac{\rho}{2}=\tfrac{3\rho}{20}$.
\end{proof}

\begin{proposition}[Local decomposition of the pressure]\label{prop:pressure-decomposition}
Let $\eta$ be as in Lemma~\ref{lem:cutoff} and let
$N(y)=-\tfrac1{4\pi\abs{y}}$ be the Newtonian kernel of $\R^{3}$, so that
$\Delta N=\delta_{0}$. Then for almost every $t\in J_{\rho}$ the identity
\begin{equation}\label{eq:laplace-etap}
  \Delta(\eta p)
  =\partial_{i}\partial_{j}(\eta U_{ij})
   +(\partial_{i}\partial_{j}\eta)U_{ij}
   -\partial_{j}\bigl(U_{ij}\partial_{i}\eta\bigr)
   -\partial_{i}\bigl(U_{ij}\partial_{j}\eta\bigr)
   -p\,\Delta\eta
   +2\,\partial_{j}\bigl((\partial_{j}\eta)p\bigr)
   +\eta\,\dv f
\end{equation}
holds in $\mathcal{D}'(\R^{3})$, all functions being extended by zero outside
$B_{\rho}$, and consequently
\begin{equation}\label{eq:pressure-decomposition}
  \eta p \;=\; p_{1}+p_{2}+p_{3}+p_{4}+p_{5}+p_{6}+p_{7}+p_{8}
  \qquad\text{a.e.\ in }\R^{3},
\end{equation}
where, with $R_{i}$ the $i$-th Riesz transform of $\R^{3}$
(External Input~\ref{ext:CZ}),
\begin{equation}\label{eq:pk}
\begin{aligned}
  p_{1}&=-R_{i}R_{j}\bigl(\eta U_{ij}\bigr), &\qquad
  p_{2}&=N*\bigl((\partial_{i}\partial_{j}\eta)U_{ij}\bigr),\\
  p_{3}&=-\partial_{j}N*\bigl(U_{ij}\partial_{i}\eta\bigr), &\qquad
  p_{4}&=-\partial_{i}N*\bigl(U_{ij}\partial_{j}\eta\bigr),\\
  p_{5}&=-N*\bigl(p\,\Delta\eta\bigr), &\qquad
  p_{6}&=2\,\partial_{j}N*\bigl((\partial_{j}\eta)p\bigr),\\
  p_{7}&=\partial_{j}N*\bigl(\eta f_{j}\bigr), &\qquad
  p_{8}&=-N*\bigl((\partial_{j}\eta)f_{j}\bigr).
\end{aligned}
\end{equation}
If, in addition, $\dv f=0$ locally in $\mathcal D'(\mathcal O)$, then
$p_{7}+p_{8}=0$ and the last two terms may be omitted. This scalar-test
condition is not part of Definition~\ref{def:sws}; it imposes no
whole-space divergence condition on a zero extension of $f$.
\end{proposition}

\begin{proof}
\emph{Step 1: the identity \eqref{eq:laplace-etap}.} Fix $t$ for which
Lemma~\ref{lem:delta-p} holds. For $\psi\in C_{c}^{\infty}(\R^{3})$ we have
$\eta\psi\in C_{c}^{\infty}(B_{\rho})$ by \eqref{eq:cutoff}, so all pairings
below are legitimate. Three elementary distributional identities are used. First,
for a scalar distribution $T$ and smooth $\eta$,
\begin{equation}\label{eq:leibniz-lap}
  \Delta(\eta T)=\eta\,\Delta T+2\,\partial_{j}\eta\,\partial_{j}T+T\,\Delta\eta .
\end{equation}
Second, for a matrix field $U$,
\begin{equation}\label{eq:commute}
  \eta\,\partial_{i}\partial_{j}U_{ij}
  =\partial_{i}\partial_{j}(\eta U_{ij})
   +(\partial_{i}\partial_{j}\eta)U_{ij}
   -\partial_{j}\bigl(U_{ij}\partial_{i}\eta\bigr)
   -\partial_{i}\bigl(U_{ij}\partial_{j}\eta\bigr),
\end{equation}
which is verified by expanding the right side:
$\partial_{i}\partial_{j}(\eta U_{ij})
 =\eta\partial_{i}\partial_{j}U_{ij}+\partial_{i}\eta\,\partial_{j}U_{ij}
  +\partial_{j}\eta\,\partial_{i}U_{ij}+(\partial_{i}\partial_{j}\eta)U_{ij}$,
while
$\partial_{j}(U_{ij}\partial_{i}\eta)
 =\partial_{j}U_{ij}\,\partial_{i}\eta+U_{ij}\partial_{i}\partial_{j}\eta$
and
$\partial_{i}(U_{ij}\partial_{j}\eta)
 =\partial_{i}U_{ij}\,\partial_{j}\eta+U_{ij}\partial_{i}\partial_{j}\eta$;
substituting and cancelling gives \eqref{eq:commute}. Third,
$2\partial_{j}\eta\,\partial_{j}p
 =2\partial_{j}\bigl((\partial_{j}\eta)p\bigr)-2p\,\Delta\eta$.
Now combine: by \eqref{eq:leibniz-lap} and
Lemma~\ref{lem:delta-p},
\[
  \Delta(\eta p)
  =\eta\bigl(\partial_{i}\partial_{j}U_{ij}+\dv f\bigr)
   +2\partial_{j}\bigl((\partial_{j}\eta)p\bigr)-2p\,\Delta\eta+p\,\Delta\eta ,
\]
and inserting \eqref{eq:commute} yields exactly \eqref{eq:laplace-etap}.

\emph{Step 2: the representation \eqref{eq:pressure-decomposition}.} The
function $\eta p(\cdot,t)$ belongs to $L^{3/2}(\R^{3})$ and has compact support
in $\overline{\Ball{3\rho/4}{x_{0}}}$. We apply
External Input~\ref{ext:newtonian} with
\[
  G_{ij}=\eta U_{ij},\qquad
  H_{j}=-U_{ij}\partial_{i}\eta-U_{ji}\partial_{i}\eta+2(\partial_{j}\eta)p+\eta f_{j},
  \qquad
  K=(\partial_{i}\partial_{j}\eta)U_{ij}-p\,\Delta\eta-(\partial_{j}\eta)f_{j}.
\]
All three have compact support in $\overline{\Ball{3\rho/4}{x_{0}}}$;
$G\in L^{3/2}$ because $U(\cdot,t)\in L^{3/2}(B_{\rho})$ by H\"older
($\abs{U}\le\abs{u}\abs{u-\spavg{u}{B_{\rho}}}$ with $u(\cdot,t)\in L^{2}\cap L^{6}$),
$H\in L^{m}$ with $m=\tfrac54\in(1,3)$, and $K\in L^{m'}$ with $m'=\tfrac54\in(1,3/2)$:
indeed each of $U(\cdot,t)$, $p(\cdot,t)$ and $f(\cdot,t)$ lies in
$L^{3/2}(B_{\rho})$ (for $f$ because $q>5/2>3/2$), all the densities are
supported in the bounded set $\overline{\Ball{3\rho/4}{x_{0}}}$, and
$L^{3/2}(B_{\rho})\subset L^{5/4}(B_{\rho})$ by H\"older. The exponent $5/4$
lies strictly inside the admissible range of
External Input~\ref{ext:newtonian}. The identity $\Delta(\eta p)=\partial_{i}\partial_{j}G_{ij}
+\partial_{j}H_{j}+K$ is \eqref{eq:laplace-etap}, and
External Input~\ref{ext:newtonian} then gives
$\eta p=-R_{i}R_{j}G_{ij}+\partial_{j}N*H_{j}+N*K$, which is
\eqref{eq:pressure-decomposition} with the terms grouped as in \eqref{eq:pk}.

\emph{Step 3: the solenoidal case.} If $\dv f=0$ in
$\mathcal{D}'(\mathcal{O})$ then, arguing as in Step 2 of the proof of
Lemma~\ref{lem:delta-p}, $\int_{B_{\rho}}f(\cdot,t)\cdot\nabla\psi=0$ for a.e.\
$t$ and every $\psi\in C_{c}^{\infty}(B_{\rho})$, i.e.\ $\dv f(\cdot,t)=0$ in
$\mathcal{D}'(B_{\rho})$. Hence the last term of \eqref{eq:laplace-etap}
vanishes, and since
$\eta\,\dv f=\dv(\eta f)-(\nabla\eta)\cdot f$ we get
$\partial_{j}N*(\eta f_{j})-N*((\partial_{j}\eta)f_{j})=N*(\eta\,\dv f)=0$, i.e.\
$p_{7}+p_{8}=0$.
\end{proof}

\subsection{Bounds for the eight terms}\label{ss:pressure-bounds}

Throughout this subsection $0<r\le\rho/2$, $\kappa=r/\rho$, $B_{r}=\Ball{r}{x_{0}}$,
and for a.e.\ $t\in J_{\rho}$ we abbreviate
\begin{equation}\label{eq:slice-norms}
\begin{aligned}
  &\mathfrak{u}(t)=\norm{u(\cdot,t)}_{L^{2}(B_{\rho})},\quad
  \mathfrak{g}(t)=\norm{\nabla u(\cdot,t)}_{L^{2}(B_{\rho})},\\
  &\mathfrak{p}(t)=\norm{p(\cdot,t)}_{L^{3/2}(B_{\rho})},\quad
  \mathfrak{f}(t)=\norm{f(\cdot,t)}_{L^{q}(B_{\rho})}.
\end{aligned}
\end{equation}
By \eqref{eq:alpha-beta},
\begin{equation}\label{eq:slice-norm-bounds}
  \operatorname*{ess\,sup}_{J_{\rho}}\mathfrak{u}\le\rho^{1/2}\alpha(\rho),
  \quad
  \norm{\mathfrak{g}}_{L^{2}(J_{\rho})}=\rho^{1/2}\beta(\rho),
  \quad
  \norm{\mathfrak{p}}_{L^{3/2}(J_{\rho})}=\rho^{4/3}\delta(\rho)^{2},
  \quad
  \norm{\mathfrak{f}}_{L^{q}(J_{\rho})}=\rho^{-\sigma}\lambda(\rho),
\end{equation}
where $\alpha(\rho)=\alpha(z_{0},\rho)$ and so on. We use three instances of
H\"older's inequality in time, valid for any measurable $g\ge0$ on $J_{\rho}$
and any $0<r\le\rho/2$:
\begin{equation}\label{eq:time-holder}
  \Bigl(\int_{t_{0}-r^{2}}^{t_{0}}g^{3/2}\Bigr)^{2/3}
  \le r^{1/3}\norm{g}_{L^{2}(J_{\rho})},
  \qquad
  \Bigl(\int_{t_{0}-r^{2}}^{t_{0}}g^{3/2}\Bigr)^{2/3}
  \le \norm{g}_{L^{3/2}(J_{\rho})},
\end{equation}
\begin{equation}\label{eq:time-holder-q}
  \Bigl(\int_{t_{0}-r^{2}}^{t_{0}}g^{3/2}\Bigr)^{2/3}
  \le r^{2(\frac23-\frac1q)}\norm{g}_{L^{q}(J_{\rho})}
  \qquad(q\ge\tfrac32).
\end{equation}
The first follows from H\"older with exponents $\tfrac43$ and $4$
(so that the measure factor is $(r^{2})^{2/3-1/2}=r^{1/3}$), the third from
H\"older with exponents $\tfrac{2q}{3}$ and its conjugate; the second is
monotonicity of the integral.

\begin{lemma}[Algebraic bounds on $U$]\label{lem:U-bounds}
Let $U$ be as in \eqref{eq:Uij}, with its Euclidean matrix norm
$\abs{U}=\bigl(\sum_{i,j}U_{ij}^{2}\bigr)^{1/2}$. For almost every $t\in J_{\rho}$,
\begin{equation}\label{eq:U-bounds}
  \norm{U(\cdot,t)}_{L^{3/2}(B_{\rho})}\le C_{5}\,\mathfrak{u}(t)\,\mathfrak{g}(t),
  \qquad
  \norm{U(\cdot,t)}_{L^{1}(B_{\rho})}
    \le\bigl(\tfrac{4\pi}{3}\bigr)^{1/3}C_{5}\,\rho\,\mathfrak{u}(t)\,\mathfrak{g}(t).
\end{equation}
\end{lemma}

\begin{proof}
Pointwise $\abs{U}\le\abs{u}\,\bigl|u-\spavg{u}{B_{\rho}}\bigr|$. By H\"older
with $\tfrac23=\tfrac12+\tfrac16$ and \eqref{eq:sobolev-poincare-6},
\[
  \norm{U(\cdot,t)}_{L^{3/2}(B_{\rho})}
  \le\norm{u(\cdot,t)}_{L^{2}(B_{\rho})}
     \bigl\|u(\cdot,t)-\spavg{u}{B_{\rho}}(t)\bigr\|_{L^{6}(B_{\rho})}
  \le C_{5}\,\mathfrak{u}(t)\mathfrak{g}(t).
\]
For the $L^{1}$ bound, H\"older with $1=\tfrac23+\tfrac13$ gives
\[
\norm{U}_{L^{1}(B_{\rho})}\le\abs{B_{\rho}}^{1/3}\norm{U}_{L^{3/2}(B_{\rho})}
 =(\tfrac{4\pi}{3})^{1/3}\rho\norm{U}_{L^{3/2}(B_{\rho})}
.
\]
\end{proof}

All space--time norms of the pressure terms below use jointly measurable
representatives. For the singular term, the localized tensor is a strongly
measurable function of time with values in the separable space $L^{3/2}$;
the bounded spatial Riesz operator preserves this measurability and admits
a jointly measurable representative. The remaining localized potentials
are handled by measurable truncated kernel integrals and their a.e. limits.
In particular, their spatial power integrals and the slice norms and
majorants used below are measurable functions of time, so that Fubini and
the time H\"older inequalities apply.

\begin{lemma}[Term-by-term bounds]\label{lem:pk-bounds}
There is an absolute constant $C_{12}$, and a constant $C_{13}=C_{13}(q)$, such
that for every $0<r\le\rho/2$, with $\kappa=r/\rho$,
\begin{align}
  \frac{1}{r^{4/3}}\norm{p_{1}}_{L^{3/2}(\Cyl{r}{z_{0}})}
    &\le C_{12}\,\kappa^{-1}\alpha(\rho)\beta(\rho),
    \label{eq:p1-bound}\\
  \frac{1}{r^{4/3}}\sum_{k=2}^{4}\norm{p_{k}}_{L^{3/2}(\Cyl{r}{z_{0}})}
    &\le C_{12}\,\kappa\,\alpha(\rho)\beta(\rho),
    \label{eq:p234-bound}\\
  \frac{1}{r^{4/3}}\sum_{k=5}^{6}\norm{p_{k}}_{L^{3/2}(\Cyl{r}{z_{0}})}
    &\le C_{12}\,\kappa^{2/3}\,\delta(\rho)^{2},
    \label{eq:p56-bound}\\
  \frac{1}{r^{4/3}}\sum_{k=7}^{8}\norm{p_{k}}_{L^{3/2}(\Cyl{r}{z_{0}})}
    &\le C_{13}\,\kappa\,\lambda(\rho).
    \label{eq:p78-bound}
\end{align}
\end{lemma}

\begin{proof}
Throughout, $c$ denotes a quantity bounded by an absolute constant times
$1+C_{5}+C_{10}(1)+C_{10}(2)+C_{11}$, where $C_{11}$ is the constant of
External Input~\ref{ext:CZ}; at the end we let $C_{12}$ be the largest of the
finitely many such quantities.

\emph{(a) The term $p_{1}$.} By External Input~\ref{ext:CZ} applied for each
fixed $t$ to each of the nine functions $\eta U_{ij}(\cdot,t)\in L^{3/2}(\R^{3})$,
\[
  \norm{p_{1}(\cdot,t)}_{L^{3/2}(\R^{3})}
  \le C_{11}\sum_{i,j}\norm{\eta U_{ij}(\cdot,t)}_{L^{3/2}(\R^{3})}
  \le 9C_{11}\norm{U(\cdot,t)}_{L^{3/2}(B_{\rho})}
  \le 9C_{11}C_{5}\,\mathfrak{u}(t)\mathfrak{g}(t),
\]
using $0\le\eta\le1$, $\supp\eta\subset B_{\rho}$ and
Lemma~\ref{lem:U-bounds}. Restricting to $B_{r}$, raising to the power $3/2$,
integrating over $(t_{0}-r^{2},t_{0})$ and using the first inequality of
\eqref{eq:time-holder} with $g=\mathfrak{u}\mathfrak{g}$ together with
\eqref{eq:slice-norm-bounds},
\[\begin{aligned}
  &\norm{p_{1}}_{L^{3/2}(\Cyl{r}{z_{0}})}
  \le 9C_{11}C_{5}\,r^{1/3}\norm{\mathfrak{u}\mathfrak{g}}_{L^{2}(J_{\rho})}
  \\ &\le 9C_{11}C_{5}\,r^{1/3}\bigl(\operatorname*{ess\,sup}\mathfrak{u}\bigr)
      \norm{\mathfrak{g}}_{L^{2}(J_{\rho})}
  \le 9C_{11}C_{5}\,r^{1/3}\rho\,\alpha(\rho)\beta(\rho).
\end{aligned}\]
Dividing by $r^{4/3}$ gives \eqref{eq:p1-bound}.

\emph{(b) The terms $p_{2},p_{3},p_{4}$.} The densities
$(\partial_{i}\partial_{j}\eta)U_{ij}$, $U_{ij}\partial_{i}\eta$ and
$U_{ij}\partial_{j}\eta$ vanish outside the annulus $\mathcal{A}_{\rho}$ of
\eqref{eq:annulus}, and for $x\in B_{r}$, $y\in\mathcal{A}_{\rho}$ we have
$\abs{x-y}\ge\tfrac{3\rho}{20}$ by Lemma~\ref{lem:cutoff}, hence
\[
  \abs{N(x-y)}=\frac{1}{4\pi\abs{x-y}}\le\frac{5}{3\pi\rho}\le\frac2\rho,
  \qquad
  \abs{\nabla N(x-y)}=\frac{1}{4\pi\abs{x-y}^{2}}
    \le\frac{100}{9\pi\rho^{2}}\le\frac4{\rho^{2}} .
\]
By Cauchy--Schwarz over the indices,
$\sum_{i,j}\abs{\partial_{i}\partial_{j}\eta}\abs{U_{ij}}\le3C_{10}(2)\rho^{-2}\abs{U}$
and $\sum_{j}\bigl|\sum_{i}U_{ij}\partial_{i}\eta\bigr|
 \le3C_{10}(1)\rho^{-1}\abs{U}$. Therefore, for $x\in B_{r}$,
\[
  \abs{p_{2}(x,t)}\le\frac2\rho\cdot\frac{3C_{10}(2)}{\rho^{2}}
     \norm{U(\cdot,t)}_{L^{1}(B_{\rho})},
  \qquad
  \abs{p_{3}(x,t)}+\abs{p_{4}(x,t)}\le2\cdot\frac{4}{\rho^{2}}\cdot
     \frac{3C_{10}(1)}{\rho}\norm{U(\cdot,t)}_{L^{1}(B_{\rho})} ,
\]
so by Lemma~\ref{lem:U-bounds}
\[
\sum_{k=2}^{4}\abs{p_{k}(x,t)}\le c\,\rho^{-2}\mathfrak{u}(t)\mathfrak{g}(t)
.
\]
Hence
\[
\norm{p_{k}(\cdot,t)}_{L^{3/2}(B_{r})}\le\abs{B_{r}}^{2/3}
 \norm{p_{k}(\cdot,t)}_{L^{\infty}(B_{r})}
 \le c\,r^{2}\rho^{-2}\mathfrak{u}(t)\mathfrak{g}(t)
,
\] and by
\eqref{eq:time-holder} and \eqref{eq:slice-norm-bounds},
\[
  \sum_{k=2}^{4}\norm{p_{k}}_{L^{3/2}(\Cyl{r}{z_{0}})}
  \le c\,\frac{r^{2}}{\rho^{2}}\,r^{1/3}\rho\,\alpha(\rho)\beta(\rho)
  = c\,\frac{r^{7/3}}{\rho}\alpha(\rho)\beta(\rho),
\]
and dividing by $r^{4/3}$ gives \eqref{eq:p234-bound}.

\emph{(c) The terms $p_{5},p_{6}$.} The densities $p\,\Delta\eta$ and
$(\partial_{j}\eta)p$ also vanish outside $\mathcal{A}_{\rho}$, so the same kernel
bounds give, for $x\in B_{r}$,
\[
\begin{aligned}
&\abs{p_{5}(x,t)}+\abs{p_{6}(x,t)}
\\ &\le\Bigl(\frac2\rho\cdot\frac{3C_{10}(2)}{\rho^{2}}
     +2\cdot\frac{4}{\rho^{2}}\cdot\frac{3C_{10}(1)}{\rho}\Bigr)
     \norm{p(\cdot,t)}_{L^{1}(B_{\rho})}
\\ &\le c\,\rho^{-3}\cdot\bigl(\tfrac{4\pi}{3}\bigr)^{1/3}\rho\,\mathfrak{p}(t)
  = c\,\rho^{-2}\mathfrak{p}(t),
\end{aligned}
\]
where $\norm{p}_{L^{1}(B_{\rho})}\le\abs{B_{\rho}}^{1/3}\mathfrak{p}$ by
H\"older. Then
$\norm{p_{5}(\cdot,t)+p_{6}(\cdot,t)}_{L^{3/2}(B_{r})}
 \le c\,r^{2}\rho^{-2}\mathfrak{p}(t)$, and by the second inequality of
\eqref{eq:time-holder} and \eqref{eq:slice-norm-bounds},
\[
  \sum_{k=5}^{6}\norm{p_{k}}_{L^{3/2}(\Cyl{r}{z_{0}})}
  \le c\,\frac{r^{2}}{\rho^{2}}\norm{\mathfrak{p}}_{L^{3/2}(J_{\rho})}
  = c\,\frac{r^{2}}{\rho^{2}}\rho^{4/3}\delta(\rho)^{2}
  = c\,r^{2}\rho^{-2/3}\delta(\rho)^{2} ,
\]
and dividing by $r^{4/3}$ gives $c\,(r/\rho)^{2/3}\delta(\rho)^{2}$, which is
\eqref{eq:p56-bound}.

\emph{(d) The term $p_{7}$.} By External Input~\ref{ext:riesz} with
$\tfrac1{15}=\tfrac25-\tfrac13$, for a.e.\ $t$,
$\norm{p_{7}(\cdot,t)}_{L^{15}(\R^{3})}
 \le C'_{11}\sum_{j=1}^{3}\norm{\eta f_j(\cdot,t)}_{L^{5/2}(\R^{3})}
 \le 3C'_{11}\norm{f(\cdot,t)}_{L^{5/2}(B_{\rho})}$. By H\"older on $B_{\rho}$
(legitimate since $q>5/2$),
$\norm{f(\cdot,t)}_{L^{5/2}(B_{\rho})}
 \le\abs{B_{\rho}}^{\frac25-\frac1q}\mathfrak{f}(t)
 = c\,\rho^{\frac65-\frac3q}\mathfrak{f}(t)$, and by H\"older on $B_{r}$,
\[
  \norm{p_{7}(\cdot,t)}_{L^{3/2}(B_{r})}
  \le\abs{B_{r}}^{\frac23-\frac1{15}}\norm{p_{7}(\cdot,t)}_{L^{15}(\R^{3})}
  \le c\,r^{9/5}\rho^{\frac65-\frac3q}\mathfrak{f}(t).
\]
Applying \eqref{eq:time-holder-q} and \eqref{eq:slice-norm-bounds},
\[
  \norm{p_{7}}_{L^{3/2}(\Cyl{r}{z_{0}})}
  \le c\,r^{9/5}\rho^{\frac65-\frac3q}\,r^{\frac43-\frac2q}\,
       \rho^{-\sigma}\lambda(\rho)
  = c\,r^{\frac{47}{15}-\frac2q}\,\rho^{-\frac95+\frac2q-\frac43+\frac43}
       \lambda(\rho),
\]
where we used $\tfrac65-\tfrac3q-\sigma=\tfrac65-\tfrac3q-3+\tfrac5q
=-\tfrac95+\tfrac2q$. Dividing by $r^{4/3}$ and recalling
$\tfrac{47}{15}-\tfrac43=\tfrac{27}{15}=\tfrac95$,
\[
  \frac{1}{r^{4/3}}\norm{p_{7}}_{L^{3/2}(\Cyl{r}{z_{0}})}
  \le c\,\kappa^{\frac95-\frac2q}\lambda(\rho)\le c\,\kappa\,\lambda(\rho),
\]
the last step because $q>5/2$ gives $\tfrac95-\tfrac2q>\tfrac95-\tfrac45=1$ and
$\kappa\le\tfrac12<1$.

\emph{(e) The term $p_{8}$.} Its density $(\partial_{j}\eta)f_{j}$ is supported in
$\mathcal{A}_{\rho}$, so for $x\in B_{r}$,
\[
  \abs{p_{8}(x,t)}\le\frac2\rho\cdot\frac{3C_{10}(1)}{\rho}
      \norm{f(\cdot,t)}_{L^{1}(B_{\rho})}
  \le c\,\rho^{-2}\abs{B_{\rho}}^{1-\frac1q}\mathfrak{f}(t)
  = c\,\rho^{1-\frac3q}\mathfrak{f}(t).
\]
Hence $\norm{p_{8}(\cdot,t)}_{L^{3/2}(B_{r})}\le c\,r^{2}\rho^{1-\frac3q}\mathfrak{f}(t)$
and, by \eqref{eq:time-holder-q},
\[
  \frac{1}{r^{4/3}}\norm{p_{8}}_{L^{3/2}(\Cyl{r}{z_{0}})}
  \le c\,r^{2-\frac43+\frac43-\frac2q}\rho^{1-\frac3q-\sigma}\lambda(\rho)
  = c\,r^{2-\frac2q}\rho^{-2+\frac2q}\lambda(\rho)
  = c\,\kappa^{2-\frac2q}\lambda(\rho)\le c\,\kappa\,\lambda(\rho),
\]
because $q>5/2$ gives $2-\tfrac2q>\tfrac65>1$. Adding (d) and (e) and taking
$C_{13}(q)$ to be the resulting constant proves \eqref{eq:p78-bound}.
\end{proof}

\subsection{The pressure decay lemma}\label{ss:pressure-decay}

\begin{theorem}[Pressure decay]\label{thm:pressure-decay}
There is an absolute constant $C_{14}$ and a constant $C_{15}=C_{15}(q)$ such
that for every suitable weak solution as above, every $z_{0}$, every $\rho$ with
$\Cyl{\rho}{z_{0}}\subset\mathcal{O}$, and every $0<r\le\rho/2$, writing
$\kappa=r/\rho$,
\begin{equation}\label{eq:pressure-decay}
  \delta(z_{0},r)
  \;\le\;C_{14}\,\kappa^{-1/2}\alpha(z_{0},\rho)^{1/2}\beta(z_{0},\rho)^{1/2}
   \;+\;C_{14}\,\kappa^{1/3}\delta(z_{0},\rho)
   \;+\;C_{15}\,\kappa^{1/2}\lambda(z_{0},\rho)^{1/2} .
\end{equation}
If moreover $\dv f=0$ then the last term may be omitted; that hypothesis is
never used below.
\end{theorem}

\begin{proof}
Since $r\le\rho/2<\tfrac{13\rho}{20}$, Lemma~\ref{lem:cutoff} gives $\eta\equiv1$
on $B_{r}$, so $p=\eta p$ on $\Cyl{r}{z_{0}}$ and
\eqref{eq:pressure-decomposition} gives
$p=\sum_{k=1}^{8}p_{k}$ there. By definition of $\delta$,
\[
  \delta(z_{0},r)^{2}
  =\Bigl(\frac1{r^{2}}\iint_{\Cyl{r}{z_{0}}}\abs{p}^{3/2}\Bigr)^{2/3}
  =\frac{1}{r^{4/3}}\norm{p}_{L^{3/2}(\Cyl{r}{z_{0}})}
  \le\frac{1}{r^{4/3}}\sum_{k=1}^{8}\norm{p_{k}}_{L^{3/2}(\Cyl{r}{z_{0}})} ,
\]
by the triangle inequality in $L^{3/2}(\Cyl{r}{z_{0}})$. Lemma~\ref{lem:pk-bounds}
and $\kappa\le\tfrac12$, which gives $\kappa\le\kappa^{-1}$, yield
\[
  \delta(z_{0},r)^{2}
  \le 2C_{12}\,\kappa^{-1}\alpha(\rho)\beta(\rho)
    +C_{12}\,\kappa^{2/3}\delta(\rho)^{2}
    +C_{13}\,\kappa\,\lambda(\rho).
\]
Taking square roots and using
$(s_{1}+s_{2}+s_{3})^{1/2}\le s_{1}^{1/2}+s_{2}^{1/2}+s_{3}^{1/2}$ for
$s_{i}\ge0$ gives \eqref{eq:pressure-decay} with
$C_{14}=(2C_{12})^{1/2}$ and $C_{15}=C_{13}^{1/2}$. Under the additional
hypothesis $\dv f=0$ one has $p_{7}+p_{8}=0$ by
Proposition~\ref{prop:pressure-decomposition} and the term with $C_{13}$ is
absent.
\end{proof}

\subsection{The harmonic part and its interior estimates}\label{ss:harmonic}

\begin{corollary}[Calder\'on--Zygmund part, harmonic part, force part]
\label{cor:CZ-harmonic}
Set
\begin{equation}\label{eq:CZ-harmonic}
  p_{\mathrm{CZ}}=p_{1}=-R_{i}R_{j}(\eta U_{ij}),
  \qquad
  p_{\mathrm{har}}=p_{2}+p_{3}+p_{4}+p_{5}+p_{6},
  \qquad
  p_{\mathrm{f}}=p_{7}+p_{8},
\end{equation}
so that $p=p_{\mathrm{CZ}}+p_{\mathrm{har}}+p_{\mathrm{f}}$ on
$\Ball{13\rho/20}{x_{0}}$ for a.e.\ $t\in J_{\rho}$, with
$p_{\mathrm{f}}=0$ when $\dv f=0$ (which Definition~\ref{def:sws} does not
assume). Then for a.e.\ $t\in J_{\rho}$ the function
$p_{\mathrm{har}}(\cdot,t)$ is harmonic on $\Ball{13\rho/20}{x_{0}}$, and for
every multi-index $a$ there is a constant $C_{16}(\abs{a})$, depending on
$\abs{a}$ and on $C_{10}(1),C_{10}(2)$ only, with
\begin{equation}\label{eq:har-Ck}
  \sup_{x\in\Ball{\rho/2}{x_{0}}}
    \bigl|\partial^{a}p_{\mathrm{har}}(x,t)\bigr|
  \;\le\;C_{16}(\abs{a})\,\rho^{-2-\abs{a}}
      \Bigl(\alpha(z_{0},\rho)^{2}+\norm{p(\cdot,t)}_{L^{3/2}(B_{\rho})}\Bigr).
\end{equation}
The exponent depends on $\abs{a}$ and not only on its bound: replacing
$\rho^{-2-\abs{a}}$ by $\rho^{-2-k}$ for all $\abs{a}\le k$ would be false for
$\rho>1$, as the suitable weak solution $u\equiv0$, $p\equiv1$, $f\equiv0$
shows, for which $p_{\mathrm{har}}\equiv1$ on $\Ball{\rho/2}{x_{0}}$ while the
right side with $k=1$ is $O(\rho^{-1})$.
\end{corollary}

\begin{proof}
Each of $p_{2},\dots,p_{6}$ is of the form $N*g$ or $\partial_{j}N*g$ with $g$
supported in $\overline{\mathcal{A}_{\rho}}\subset\R^{3}\setminus\Ball{13\rho/20}{x_{0}}$
and $g(\cdot,t)\in L^{1}(\R^{3})$. For $x$ in the open set
$\Ball{13\rho/20}{x_{0}}$ the kernel $y\mapsto N(x-y)$ is smooth and harmonic in
$x$ on $\supp g$, and $\abs{\partial_{x}^{a}N(x-y)}\le c_{a}\abs{x-y}^{-1-\abs{a}}$
is integrable against $g(\cdot,t)\,dy$ locally uniformly in $x$; hence
differentiation under the integral sign is licit and
$\Delta_{x}(N*g)(x)=\int\Delta_{x}N(x-y)g(y)\,dy=0$ there. The same applies to
$\partial_{j}N*g$. So $p_{\mathrm{har}}(\cdot,t)$ is harmonic on
$\Ball{13\rho/20}{x_{0}}$.

For \eqref{eq:har-Ck}, let $x\in\Ball{\rho/2}{x_{0}}$ and
$y\in\mathcal{A}_{\rho}$, so $\abs{x-y}\ge\tfrac{3\rho}{20}$ by
Lemma~\ref{lem:cutoff}. Differentiating under the integral,
$\abs{\partial^{a}p_{\mathrm{har}}(x,t)}$ is bounded by
$c_{a}\rho^{-1-\abs{a}}$ times the total mass of the densities, one extra factor
$\rho^{-1}$ being present for $p_{3},p_{4},p_{6}$ because the kernel is
$\partial_{j}N$; together with
$\abs{\partial\eta}\le C_{10}(1)\rho^{-1}$,
$\abs{\partial^{2}\eta}\le C_{10}(2)\rho^{-2}$ this gives
\[
  \abs{\partial^{a}p_{\mathrm{har}}(x,t)}
  \le c_{a}\bigl(C_{10}(1)+C_{10}(2)\bigr)\rho^{-3-\abs{a}}
     \Bigl(\norm{U(\cdot,t)}_{L^{1}(B_{\rho})}+\norm{p(\cdot,t)}_{L^{1}(B_{\rho})}\Bigr).
\]
Now $\norm{p(\cdot,t)}_{L^{1}(B_{\rho})}\le(\tfrac{4\pi}{3})^{1/3}\rho\,
\norm{p(\cdot,t)}_{L^{3/2}(B_{\rho})}$ by H\"older, and
\[
  \norm{U(\cdot,t)}_{L^{1}(B_{\rho})}
  \le\norm{u(\cdot,t)}_{L^{2}(B_{\rho})}
     \bigl\|u(\cdot,t)-\spavg{u}{B_{\rho}}(t)\bigr\|_{L^{2}(B_{\rho})}
  \le2\norm{u(\cdot,t)}_{L^{2}(B_{\rho})}^{2}
  \le2\rho\,\alpha(z_{0},\rho)^{2},
\]
where the middle step uses that $g\mapsto g-\spavg{g}{B_{\rho}}$ has operator
norm at most $2$ on $L^{2}(B_{\rho})$ (indeed at most $1$, by orthogonality), and
the last is \eqref{eq:slice-norm-bounds}. Substituting and absorbing the numeric
factors into $C_{16}(\abs{a})$ gives \eqref{eq:har-Ck}.
\end{proof}

\begin{lemma}[Interior estimates for harmonic functions in $L^{3/2}$]\label{lem:harmonic-interior}
There is an absolute constant $C_{17}$ with the following property. Let
$h$ be harmonic on $\Ball{\rho}{x_{0}}$ with $h\in L^{3/2}(\Ball{\rho}{x_{0}})$.
Then
\begin{align}
  \norm{h}_{L^{\infty}(\Ball{3\rho/4}{x_{0}})}
    &\le C_{17}\,\rho^{-2}\norm{h}_{L^{3/2}(\Ball{\rho}{x_{0}})},
    \label{eq:harm-sup}\\
  \norm{\nabla h}_{L^{\infty}(\Ball{\rho/2}{x_{0}})}
    &\le C_{17}\,\rho^{-3}\norm{h}_{L^{3/2}(\Ball{\rho}{x_{0}})},
    \label{eq:harm-grad}
\end{align}
and for every $0<r\le\rho/2$,
\begin{align}
  \frac1{r^{2}}\int_{\Ball{r}{x_{0}}}\abs{h}^{3/2}
    &\le C_{17}\,\frac r\rho\cdot\frac1{\rho^{2}}
       \int_{\Ball{\rho}{x_{0}}}\abs{h}^{3/2},
    \label{eq:harm-decay-plain}\\
  \frac1{r^{2}}\int_{\Ball{r}{x_{0}}}
     \bigl|h-\spavg{h}{\Ball{r}{x_{0}}}\bigr|^{3/2}
    &\le C_{17}\,\Bigl(\frac r\rho\Bigr)^{5/2}\cdot\frac1{\rho^{2}}
       \int_{\Ball{\rho}{x_{0}}}\abs{h}^{3/2}.
    \label{eq:harm-decay-osc}
\end{align}
\end{lemma}

\begin{proof}
We use the classical smooth harmonic representative of $h$.
If $h$ is initially given as an $L^{3/2}$ function with $\Delta h=0$
distributionally, Weyl's lemma gives an a.e.-equal smooth harmonic
representative on the ball. All supremum and gradient estimates below
refer to this representative; its $L^{3/2}$ norm agrees with that of $h$.

\emph{\eqref{eq:harm-sup}.} Let $x\in\Ball{3\rho/4}{x_{0}}$. Then
$\Ball{\rho/4}{x}\subset\Ball{\rho}{x_{0}}$, and the mean value property
(External Input~\ref{ext:harmonic}) gives
$h(x)=\abs{\Ball{\rho/4}{x}}^{-1}\int_{\Ball{\rho/4}{x}}h$. Hence, by H\"older
with $1=\tfrac23+\tfrac13$,
\[
  \abs{h(x)}\le\abs{\Ball{\rho/4}{x}}^{-1}\abs{\Ball{\rho}{x_{0}}}^{1/3}
     \norm{h}_{L^{3/2}(\Ball{\rho}{x_{0}})}
  =\bigl(\tfrac{4\pi}{3}\bigr)^{-2/3}4^{3}\,\rho^{-2}
     \norm{h}_{L^{3/2}(\Ball{\rho}{x_{0}})} .
\]

\emph{\eqref{eq:harm-grad}.} Each $\partial_{i}h$ is harmonic on
$\Ball{\rho}{x_{0}}$, and for $x\in\Ball{\rho/2}{x_{0}}$ the gradient estimate
(External Input~\ref{ext:harmonic}) on the ball $\Ball{\rho/4}{x}\subset
\Ball{3\rho/4}{x_{0}}$ gives
$\abs{\nabla h(x)}\le\tfrac{3}{\rho/4}\norm{h}_{L^{\infty}(\Ball{3\rho/4}{x_{0}})}$;
combining with \eqref{eq:harm-sup} proves \eqref{eq:harm-grad}.

\emph{\eqref{eq:harm-decay-plain}.} For $r\le\rho/2$,
\[
\begin{aligned}
&\int_{\Ball{r}{x_{0}}}\abs{h}^{3/2}
\\ &\le\abs{\Ball{r}{x_{0}}}
 \norm{h}_{L^{\infty}(\Ball{3\rho/4}{x_{0}})}^{3/2}
\\ &\le\tfrac{4\pi}{3}r^{3}\bigl(C_{17}'\rho^{-2}\bigr)^{3/2}
 \bigl(\int_{\Ball{\rho}{x_{0}}}\abs{h}^{3/2}\bigr)
\end{aligned}
\]
 by \eqref{eq:harm-sup},
where $C_{17}'$ is the constant there. Dividing by $r^{2}$ gives
$r\rho^{-3}$ times $\int_{\Ball{\rho}{x_{0}}}\abs{h}^{3/2}$, i.e.\
\eqref{eq:harm-decay-plain}.

\emph{\eqref{eq:harm-decay-osc}.} For $x\in\Ball{r}{x_{0}}$,
$\bigl|h(x)-\spavg{h}{\Ball{r}{x_{0}}}\bigr|
 \le\sup_{y,y'\in\Ball{r}{x_{0}}}\abs{h(y)-h(y')}
 \le2r\norm{\nabla h}_{L^{\infty}(\Ball{\rho/2}{x_{0}})}$, since
$\Ball{r}{x_{0}}\subset\Ball{\rho/2}{x_{0}}$ is convex. Hence, by
\eqref{eq:harm-grad},
\[
  \int_{\Ball{r}{x_{0}}}\bigl|h-\spavg{h}{\Ball{r}{x_{0}}}\bigr|^{3/2}
  \le\tfrac{4\pi}{3}r^{3}\bigl(2rC_{17}'\rho^{-3}\bigr)^{3/2}
     \int_{\Ball{\rho}{x_{0}}}\abs{h}^{3/2}
  = c\,r^{9/2}\rho^{-9/2}\int_{\Ball{\rho}{x_{0}}}\abs{h}^{3/2},
\]
and dividing by $r^{2}$ gives $c\,(r/\rho)^{5/2}\rho^{-2}
\int_{\Ball{\rho}{x_{0}}}\abs{h}^{3/2}$.
\end{proof}

\subsection{The oscillation form of the pressure estimate}\label{ss:lin34}

The estimates of Lemma~\ref{lem:pk-bounds} measure the velocity through
$\alpha\beta$. For the scale iteration one also needs a form in which the
velocity enters only through its oscillation. Define, for
$\Cyl{\rho}{z_{0}}\subset\mathcal{O}$,
\begin{equation}\label{eq:Chat}
  \widehat C(z_{0},\rho)
  \;=\;\frac1{\rho^{2}}\iint_{\Cyl{\rho}{z_{0}}}
       \bigl|u(y,s)-\spavg{u}{B_{\rho}}(s)\bigr|^{3}dy\,ds .
\end{equation}

\begin{lemma}[$\Delta p$ with the doubly centred nonlinearity]\label{lem:delta-p-centred}
For a.e.\ $t\in J_{\rho}$, $\Delta p=\partial_{i}\partial_{j}\widehat U_{ij}+\dv f$
in $\mathcal{D}'(B_{\rho})$, where
\begin{equation}\label{eq:Uhat}
  \widehat U_{ij}(y,t)
  =-\bigl(u_{i}(y,t)-\spavg{u_{i}}{B_{\rho}}(t)\bigr)
    \bigl(u_{j}(y,t)-\spavg{u_{j}}{B_{\rho}}(t)\bigr).
\end{equation}
\end{lemma}

\begin{proof}
Write $c=\spavg{u}{B_{\rho}}(t)\in\R^{3}$, a constant vector at fixed $t$. Expanding,
$(u_{i}-c_{i})(u_{j}-c_{j})=u_{i}u_{j}-c_{i}u_{j}-c_{j}u_{i}+c_{i}c_{j}$. For
$\psi\in C_{c}^{\infty}(B_{\rho})$ and a.e.\ $t$,
$\int c_{i}u_{j}\partial_{i}\partial_{j}\psi
 =c_{i}\int u_{j}\partial_{j}(\partial_{i}\psi)=0$ and
$\int c_{j}u_{i}\partial_{i}\partial_{j}\psi
 =c_{j}\int u_{i}\partial_{i}(\partial_{j}\psi)=0$ by \eqref{eq:divfree-ae},
while $\int c_{i}c_{j}\partial_{i}\partial_{j}\psi=0$ because $\psi$ has compact
support. Hence
$\int u_{i}u_{j}\partial_{i}\partial_{j}\psi
 =\int(u_{i}-c_{i})(u_{j}-c_{j})\partial_{i}\partial_{j}\psi$, and substituting
into \eqref{eq:delta-p-raw} gives the claim.
\end{proof}

\begin{proposition}[Oscillation form of the pressure estimate]
\label{prop:lin34}
There is an absolute constant $C_{18}$ such that for every
$z_{0}$ and $\rho>0$ with $\overline{\Cyl{\rho}{z_{0}}}\subset\mathcal{O}$ and every
$0<r\le\rho/2$, the following hold.
\begin{enumerate}
\item[(i)] \emph{(Pointwise in time, when $\dv f=0$.)} For a.e.\ $t\in J_{\rho}$,
\begin{equation}\label{eq:lin34-pointwise}
  \frac1{r^{2}}\int_{\Ball{r}{x_{0}}}\abs{p(y,t)}^{3/2}dy
  \;\le\;C_{18}\Bigl(\frac\rho r\Bigr)^{2}\frac1{\rho^{2}}
     \int_{B_{\rho}}\bigl|u(y,t)-\spavg{u}{B_{\rho}}(t)\bigr|^{3}dy
   \;+\;C_{18}\,\frac r\rho\cdot\frac1{\rho^{2}}
     \int_{B_{\rho}}\abs{p(y,t)}^{3/2}dy .
\end{equation}
\item[(ii-a)] \emph{(Integrated in time, when $\dv f=0$.)} If $\dv f=0$ in
$\mathcal{D}'(\mathcal{O})$ then
\begin{equation}\label{eq:lin35}
  D(z_{0},r)\;\le\;C_{18}\Bigl(\frac\rho r\Bigr)^{2}\widehat C(z_{0},\rho)
   \;+\;C_{18}\,\frac r\rho\,D(z_{0},\rho).
\end{equation}
\item[(ii-b)] \emph{(Integrated in time, general force.)} Without any hypothesis
on $\dv f$, with $C_{32}(q)=2^{1/2}\bigl(C_{18}+C_{13}(q)^{3/2}\bigr)$,
\begin{equation}\label{eq:lin35-force}
  D(z_{0},r)\;\le\;C_{32}(q)\Bigl[\Bigl(\frac\rho r\Bigr)^{2}\widehat C(z_{0},\rho)
   +\frac r\rho\,D(z_{0},\rho)
   +\Bigl(\frac r\rho\Bigr)^{3/2}\lambda(z_{0},\rho)^{3/2}\Bigr].
\end{equation}
\end{enumerate}
\end{proposition}

\begin{proof}
Run the decomposition of Proposition~\ref{prop:pressure-decomposition} with
$\widehat U$ of \eqref{eq:Uhat} in place of $U$, which is legitimate by
Lemma~\ref{lem:delta-p-centred}. In part (i) we assume $\dv f=0$, so that
$p_{7}+p_{8}=0$ and $p=\sum_{k=1}^{6}p_{k}$ on
$\Ball{13\rho/20}{x_{0}}\supset\Ball{r}{x_{0}}$; part (ii) adds the
contribution of $p_{7}+p_{8}$ at the end. Write
$w(\cdot,t)=u(\cdot,t)-\spavg{u}{B_{\rho}}(t)$, so that
$\abs{\widehat U}\le\abs{w}^{2}$ and
$\norm{\widehat U(\cdot,t)}_{L^{3/2}(B_{\rho})}\le\norm{w(\cdot,t)}^{2}_{L^{3}(B_{\rho})}$.

\emph{(i).} For $p_{1}$, External Input~\ref{ext:CZ} gives
$\norm{p_{1}(\cdot,t)}_{L^{3/2}(\R^{3})}\le9C_{11}\norm{w(\cdot,t)}^{2}_{L^{3}(B_{\rho})}$,
hence
\[
  \int_{\Ball{r}{x_{0}}}\abs{p_{1}(\cdot,t)}^{3/2}
  \le\bigl(9C_{11}\bigr)^{3/2}\int_{B_{\rho}}\abs{w(\cdot,t)}^{3},
\]
and dividing by $r^{2}$ yields the first term of \eqref{eq:lin34-pointwise}
with the factor $(\rho/r)^{2}$. For $p_{2},p_{3},p_{4}$, the far-field bounds of
Lemma~\ref{lem:pk-bounds}(b) with $U$ replaced by $\widehat U$ give, for
$x\in\Ball{r}{x_{0}}$,
\[
  \sum_{k=2}^{4}\abs{p_{k}(x,t)}
  \le c\,\rho^{-3}\norm{\widehat U(\cdot,t)}_{L^{1}(B_{\rho})}
  \le c\,\rho^{-3}\abs{B_{\rho}}^{1/3}\norm{w(\cdot,t)}^{2}_{L^{3}(B_{\rho})}
  = c\,\rho^{-2}\norm{w(\cdot,t)}^{2}_{L^{3}(B_{\rho})} ,
\]
so
$\int_{\Ball{r}{x_{0}}}\bigl(\sum_{k=2}^{4}\abs{p_{k}}\bigr)^{3/2}
 \le c\,r^{3}\rho^{-3}\int_{B_{\rho}}\abs{w(\cdot,t)}^{3}$, and dividing by
$r^{2}$ gives $c\,(r/\rho)\rho^{-2}\int_{B_{\rho}}\abs{w}^{3}
\le c(\rho/r)^{2}\rho^{-2}\int_{B_{\rho}}\abs{w}^{3}$. For $p_{5},p_{6}$,
Lemma~\ref{lem:pk-bounds}(c) gives
$\abs{p_{5}}+\abs{p_{6}}\le c\rho^{-2}\norm{p(\cdot,t)}_{L^{3/2}(B_{\rho})}$ on
$\Ball{r}{x_{0}}$, so
\[
  \int_{\Ball{r}{x_{0}}}\bigl(\abs{p_{5}}+\abs{p_{6}}\bigr)^{3/2}
  \le c\,r^{3}\rho^{-3}\int_{B_{\rho}}\abs{p(\cdot,t)}^{3/2},
\]
and dividing by $r^{2}$ gives $c\,(r/\rho)\rho^{-2}\int_{B_{\rho}}\abs{p}^{3/2}$,
the second term of \eqref{eq:lin34-pointwise}. Summing the three groups, using
$\bigl(\sum_{k}\abs{p_{k}}\bigr)^{3/2}\le3^{1/2}\sum_{k}\abs{p_{k}}^{3/2}$ for
three groups, and absorbing numeric factors into $C_{18}$ proves (i).

\emph{(ii-a).} Integrate \eqref{eq:lin34-pointwise} over
$t\in(t_{0}-r^{2},t_{0})$ and enlarge the two time integrals on the right from
$(t_{0}-r^{2},t_{0})$ to $J_{\rho}=(t_{0}-\rho^{2},t_{0})$, which is legitimate
since both integrands are nonnegative. This gives exactly \eqref{eq:lin35}.

\emph{(ii-b).} For a general force write $p=\bigl(\sum_{k=1}^{6}p_{k}\bigr)
+(p_{7}+p_{8})$ on $\Ball{13\rho/20}{x_{0}}$ and use
$\abs{a+b}^{3/2}\le2^{1/2}(\abs{a}^{3/2}+\abs{b}^{3/2})$. The first group is
bounded exactly as in (ii-a), and by Lemma~\ref{lem:pk-bounds}(d)--(e)
\[
  \frac1{r^{2}}\iint_{\Cyl{r}{z_{0}}}\abs{p_{7}+p_{8}}^{3/2}
  =\Bigl(\frac{1}{r^{4/3}}\norm{p_{7}+p_{8}}_{L^{3/2}(\Cyl{r}{z_{0}})}\Bigr)^{3/2}
  \le\bigl(C_{13}(q)\,\kappa\,\lambda(z_{0},\rho)\bigr)^{3/2},
\]
which gives \eqref{eq:lin35-force} with
$C_{32}(q)=2^{1/2}(C_{18}+C_{13}(q)^{3/2})$.
\end{proof}

\begin{remark}[The pressure-decay step in Lin's proof]\label{rem:lin34-defect}
Lin \cite[Lemma 3.4]{Lin1998} uses a coefficient $C_{\theta_0}$ depending
on the lower scale ratio $\theta_0$, not the small factor $C\theta_0$.
His Corollary 3.5 restricts $r>\theta_0\rho$, so the factor $(\rho/r)^2$
multiplying the velocity oscillation in \eqref{eq:lin34-pointwise} can
be absorbed into $C_{\theta_0}$. The velocity term of the corollary is
therefore consistent with that estimate.

The gap in the deduction from the printed lemma concerns the harmonic
pressure term. The useful unnormalized estimate retains the factor
$(r/\rho)^3$ in that term, supplied by the harmonic interior estimate.
After division by $r^2$ this becomes
$(r/\rho)\rho^{-2}\int_{B_\rho}|p|^{3/2}$, as in the corollary.
Without the cubic factor, direct rescaling gives
$(\rho/r)^2\rho^{-2}\int_{B_\rho}|p|^{3/2}$ instead.
Thus the lemma as printed is too weak to imply its corollary; the
missing cubic decay belongs to the harmonic term. We use
\eqref{eq:lin35}, retaining the explicit scale ratios throughout.
\end{remark}

\section{The \texorpdfstring{$\varepsilon$}{epsilon}-regularity core}\label{sec:core}

This section proves Theorem~\ref{thm:A} and Theorem~\ref{thm:B}. The
architecture is that of Lemarié-Rieusset \cite[\S13.8--13.9]{LemarieRieusset2016},
which combines the scale iteration of Kukavica \cite{Kukavica2008} with the
parabolic Morrey/potential endgame of O'Leary \cite{OLeary2003}. There are four
steps.
\begin{itemize}
\item[\textbf{Step 1}] (\S\ref{ss:step1}) The local energy inequality and the
  pressure decomposition of Section~\ref{sec:pressure} give a pair of coupled
  decay estimates for $\alpha,\beta,\delta$ at two scales. This is
  \cite[Lemma 13.3, (13.30)--(13.31), pp.~430--434]{LemarieRieusset2016}; we
  obtain it with Kukavica's backward-Gaussian test function.
\item[\textbf{Step 2}] (\S\ref{ss:step2}) Iterating these gives, on a
  neighbourhood of a point at which $\limsup_{r\to0}\beta(z_{0},r)$ is small,
  the parabolic Morrey bounds
  $\mathbf{1}_{Q}u\in\mathcal{M}^{3,\tau_{2}}$,
  $\mathbf{1}_{Q}\nabla u\in\mathcal{M}^{2,\tau_{3}}$,
  $\mathbf{1}_{Q}p\in\mathcal{M}^{3/2,\tau_{p}}$ with $\tau_{2}>5$. We use
  Kukavica's Lemma 3; compare
  \cite[Lemma 13.4, pp.~435--439]{LemarieRieusset2016}.
  The pressure exponents differ, as explained in Remark~\ref{rem:LR-pressure}.
\item[\textbf{Step 3}] (\S\ref{ss:step3}) A representation of $\phi u$ by
  parabolic Riesz potentials, together with the Adams--Hedberg inequality,
  improves the Morrey exponent of $u$ in finitely many rounds. This is
  \cite[Lemma 13.5, pp.~439--441]{LemarieRieusset2016}.
\item[\textbf{Step 4}] (\S\ref{ss:step4}) Once the exponents are large enough,
  the same representation and the parabolic Campanato lemma give H\"older
  continuity. This is
  \cite[Lemma 13.6 and Proposition 13.4]{LemarieRieusset2016}.
\end{itemize}

\begin{convention}[Standing choices in this section]\label{conv:step-params}
Throughout Section~\ref{sec:core}:
\begin{equation}\label{eq:standing}\begin{gathered}
  q>\tfrac52,\qquad \sigma=3-\tfrac5q>1,\qquad
  \varepsilon=\tfrac25,\qquad
  \tau_{2}=\frac{5}{1-\varepsilon}=\frac{25}{3},\\
  \frac1{\tau_{3}}=\frac1{\tau_{2}}+\frac15,\ \ \tau_{3}=\frac{25}{8},
  \qquad \tau_{p}=\frac{25}{8}.
\end{gathered}\end{equation}
The value $\varepsilon=2/5$ is admissible because
$\varepsilon<\min\{2/3,\sigma\}$ (Proposition~\ref{prop:morrey-decay}), and it
gives $\tau_{2}=25/3$ and $\tau_{3}=25/8$.
Steps 3 and 4 use these actual values in the product and potential
exponent identities, not merely $\tau_2>5$ and $\tau_3>5/2$. The pressure exponent $\tau_{p}=25/8$ is recorded for
completeness but is \emph{not} consumed downstream: the pressure enters Steps 3
and 4 through fixed-scale $L^{3/2}$ bounds for the harmonic remainder
in Lemma~\ref{lem:pressure-gradient-morrey}. See
Remark~\ref{rem:LR-pressure}. The pressure is taken in
$L^{3/2}_{loc}$, i.e.\ the exponent called $q_{0}$ in
\cite{LemarieRieusset2016} is $q_{0}=3/2$; see Remark~\ref{rem:LR-pressure}.
\end{convention}

\subsection{Parabolic Morrey spaces and Riesz potentials}\label{ss:morrey}

\begin{definition}[Parabolic Morrey space]\label{def:parabolic-morrey}
For $1\le P<\infty$ and $P\le\tau<\infty$ and $g\in L^{P}_{loc}(\R^{3}\times\R)$
put
\begin{equation}\label{eq:morrey-norm}
  \norm{g}_{\mathcal{M}^{P,\tau}}
  \;=\;\sup_{z\in\R^{3}\times\R}\ \sup_{r>0}\;
    r^{-5\left(\frac1P-\frac1\tau\right)}
    \Bigl(\iint_{\mathfrak{B}_{r}(z)}\abs{g}^{P}\,dz'\Bigr)^{1/P},
\end{equation}
with $\mathfrak{B}_{r}(z)$ the $\dpar$-ball \eqref{eq:parabolic-ball}, and
$\mathcal{M}^{P,\tau}=\{g:\norm{g}_{\mathcal{M}^{P,\tau}}<\infty\}$.
\end{definition}

The number $5$ is the homogeneous dimension of $(\R^{3}\times\R,\dpar)$: by
Lemma~\ref{lem:parabolic-metric},
$\abs{\mathfrak{B}_{r}(z)}=\tfrac{8\pi}{3}r^{5}$. Note
$L^{\tau}\subset\mathcal{M}^{P,\tau}$ and
$\mathcal{M}^{\tau,\tau}=L^{\tau}$.

\begin{lemma}[Cylinders, one-sided or two-sided, give the same space]
\label{lem:morrey-cylinders}
Let $1\le P\le\tau<\infty$ and put, for $g$ measurable,
\[
\begin{aligned}
&N_{1}(g)=\sup_{z,r}r^{-5(\frac1P-\frac1\tau)}
    \Bigl(\iint_{\Cyl{r}{z}}\abs{g}^{P}\Bigr)^{1/P},
  \\
 &  N_{2}(g)=\sup_{z,r}r^{-5(\frac1P-\frac1\tau)}
    \Bigl(\iint_{\mathfrak{B}_{r}(z)}\abs{g}^{P}\Bigr)^{1/P}
   =\norm{g}_{\mathcal{M}^{P,\tau}} .
\end{aligned}
\]
Then
\begin{equation}\label{eq:morrey-cyl}
  N_{1}(g)\;\le\;N_{2}(g)\;\le\;2^{5(\frac1P-\frac1\tau)}N_{1}(g).
\end{equation}
\end{lemma}

\begin{proof}
The first inequality is $\Cyl{r}{z}\subset\mathfrak{B}_{r}(z)$
(Lemma~\ref{lem:parabolic-metric}). For the second, fix $z=(x,t)$ and $r>0$
and set $z'=(x,t+r^{2})$. By Lemma~\ref{lem:parabolic-metric},
$\mathfrak{B}_{r}(z)\subset\Cyl{2r}{z'}$, so
$\iint_{\mathfrak{B}_{r}(z)}\abs{g}^{P}\le\iint_{\Cyl{2r}{z'}}\abs{g}^{P}
 \le N_{1}(g)^{P}\,(2r)^{5P(\frac1P-\frac1\tau)}$. Taking $P$-th roots and
dividing by $r^{5(\frac1P-\frac1\tau)}$ gives the claim.
\end{proof}

\begin{remark}[One-sided cylinders are what the energy inequality produces]
\label{rem:one-sided}
The local energy inequality is naturally an estimate on the one-sided cylinders
$\Cyl{r}{z}=\Ball{r}{x}\times(t-r^{2},t]$, and the hypothesis of
Theorem~\ref{thm:B} is stated on these cylinders. The original
\cite[Proposition 2]{CKN1982} instead uses a shifted backward cylinder;
we do not identify its cylinder convention with this one.
Lemarié-Rieusset \cite[\S13.8]{LemarieRieusset2016} uses the two-sided
cylinders $(t-r^{2},t+r^{2})\times\Ball{r}{x}$ throughout, so his hypothesis is
formally stronger than ours. Lemma~\ref{lem:morrey-cylinders} shows the
difference is immaterial for Morrey norms, and
Lemma~\ref{lem:step2-morrey-balls} below converts the one-sided bounds produced
by Step~2 into the two-sided form used by Steps 3 and 4: the point is that
Step~2 produces its bounds simultaneously at \emph{every} centre in a
neighbourhood, and $\mathfrak{B}_{r}(z)\subset\Cyl{\sqrt2 r}{(x,t+r^{2})}$.
\end{remark}

\begin{remark}[Normalization of local norms]
The metric ball has volume $(8\pi/3)r^5$; the backward cylinder has
volume $(4\pi/3)r^5$. Thus lowering the local integrability exponent from
$P$ to $P'$ contributes the volume of the unit domain to the power
$1/P'-1/P$, using the domain of integration. These factors, and the
comparison in \eqref{eq:morrey-cyl}, cannot be replaced by literal equality
of ball and cylinder norms.
\end{remark}

\begin{definition}[Parabolic Riesz potential]\label{def:riesz-potential}
For $a\in(0,5)$ and $g\in L^{1}_{loc}(\R^{3}\times\R)$ set
\begin{equation}\label{eq:riesz-potential}
  (I_{a}g)(z)\;=\;\iint_{\R^{3}\times\R}
     \frac{\abs{g(w)}}{\dpar(z,w)^{5-a}}\,dw ,
  \qquad z\in\R^{3}\times\R .
\end{equation}
\end{definition}

\begin{externalinput}[Parabolic Hardy--Littlewood maximal theorem]\label{ext:maximal}
For $g\in L^{1}_{loc}(\R^{3}\times\R)$ put
\[
  (\mathcal{M}g)(z)=\sup_{r>0}\frac{1}{\abs{\mathfrak{B}_{r}(z)}}
     \iint_{\mathfrak{B}_{r}(z)}\abs{g}\,dw .
\]
Then $\mathcal{M}g$ is lower semicontinuous, hence measurable. The doubling
property $\abs{\mathfrak{B}_{2r}(z)}=2^{5}\abs{\mathfrak{B}_{r}(z)}$ and the
Vitali covering lemma give the weak $(1,1)$ bound
$\abs{\{\mathcal{M}g>\varsigma\}}\le C\varsigma^{-1}\norm{g}_{L^{1}}$; the
trivial $L^{\infty}$ bound and Marcinkiewicz interpolation then give, for every
$P\in(1,\infty)$, a constant $C_{19}(P)$ with
$\norm{\mathcal{M}g}_{L^{P}}\le C_{19}(P)\norm{g}_{L^{P}}$. The same
construction yields the Lebesgue differentiation theorem on
$(\R^{3}\times\R,\dpar,dz)$, used in Lemma~\ref{lem:campanato}.
Reference: \cite[Ch.~I]{Stein1970}; the metric measure space
$(\R^{3}\times\R,\dpar,dz)$ is of homogeneous type.
\end{externalinput}

\begin{lemma}[Hedberg's inequality]\label{lem:hedberg}
Let $1\le P\le\tau<\infty$ and $0<a<5/\tau$. Set
\begin{equation}\label{eq:hedberg-lambda}
  \varkappa\;=\;\frac{a\tau}{5}\in(0,1).
\end{equation}
There is a constant $C_{20}(a,\tau)$ such that for every
$g\in\mathcal{M}^{P,\tau}$ and every $z$,
\begin{equation}\label{eq:hedberg}
  (I_{a}g)(z)\;\le\;C_{20}(a,\tau)\,
    \bigl(\mathcal{M}g(z)\bigr)^{1-\varkappa}\,
    \norm{g}_{\mathcal{M}^{P,\tau}}^{\varkappa}.
\end{equation}
\end{lemma}

\begin{proof}
Fix $z$ and $R>0$ and split $I_{a}g(z)=A_{R}+B_{R}$ with $A_{R}$ the integral
over $\{\dpar(z,w)<R\}=\mathfrak{B}_{R}(z)$ and $B_{R}$ the integral over the
complement.

\emph{Near part.} Decompose $\mathfrak{B}_{R}(z)$ into the dyadic shells
$S_{j}=\mathfrak{B}_{2^{-j}R}(z)\setminus\mathfrak{B}_{2^{-j-1}R}(z)$,
$j\ge0$. On $S_{j}$ we have $\dpar(z,w)^{-(5-a)}\le(2^{-j-1}R)^{-(5-a)}$, so
\[\begin{aligned}
  &A_{R}\le\sum_{j\ge0}(2^{-j-1}R)^{-(5-a)}\iint_{\mathfrak{B}_{2^{-j}R}(z)}\abs{g}
  \\ &\le\sum_{j\ge0}(2^{-j-1}R)^{-(5-a)}\,\tfrac{8\pi}{3}(2^{-j}R)^{5}\,\mathcal{M}g(z)
  \\ &=\tfrac{8\pi}{3}2^{5-a}R^{a}\,\mathcal{M}g(z)\sum_{j\ge0}2^{-ja},
\end{aligned}\]
so $A_{R}\le C R^{a}\mathcal{M}g(z)$ with
$C=\tfrac{8\pi}{3}2^{5-a}(1-2^{-a})^{-1}$.

\emph{Far part.} Decompose the complement into
$T_{j}=\mathfrak{B}_{2^{j+1}R}(z)\setminus\mathfrak{B}_{2^{j}R}(z)$, $j\ge0$. On
$T_{j}$, $\dpar(z,w)^{-(5-a)}\le(2^{j}R)^{-(5-a)}$. By H\"older's inequality on
$\mathfrak{B}_{2^{j+1}R}(z)$ with exponents $P$ and $P'$, and
Definition~\ref{def:parabolic-morrey},
\[\begin{aligned}
  &\iint_{\mathfrak{B}_{2^{j+1}R}(z)}\abs{g}
  \le\abs{\mathfrak{B}_{2^{j+1}R}}^{1-\frac1P}
     \Bigl(\iint_{\mathfrak{B}_{2^{j+1}R}}\abs{g}^{P}\Bigr)^{1/P}
  \\ &\le C'(2^{j+1}R)^{5(1-\frac1P)}(2^{j+1}R)^{5(\frac1P-\frac1\tau)}
      \norm{g}_{\mathcal{M}^{P,\tau}}
  \\ &= C'(2^{j+1}R)^{5(1-\frac1\tau)}\norm{g}_{\mathcal{M}^{P,\tau}},
\end{aligned}\]
with $C'=(\tfrac{8\pi}{3})^{1-1/P}$. Hence
\[
  B_{R}\le C'\,2^{5}\norm{g}_{\mathcal{M}^{P,\tau}}
    \sum_{j\ge0}(2^{j}R)^{a-\frac5\tau}
  = C''R^{a-\frac5\tau}\norm{g}_{\mathcal{M}^{P,\tau}},
\]
the geometric series converging because $a-\tfrac5\tau<0$ by hypothesis.

\emph{Optimisation.} If $\mathcal{M}g(z)=0$ then $g=0$ a.e.\ and there is
nothing to prove; if $\mathcal{M}g(z)=\infty$ the bound is trivial. Otherwise
choose $R$ with
$R^{5/\tau}=\norm{g}_{\mathcal{M}^{P,\tau}}/\mathcal{M}g(z)$, that is
$R=\bigl(\norm{g}_{\mathcal{M}^{P,\tau}}/\mathcal{M}g(z)\bigr)^{\tau/5}$. Then
$R^{a}\mathcal{M}g(z)
 =\norm{g}_{\mathcal{M}^{P,\tau}}^{a\tau/5}\mathcal{M}g(z)^{1-a\tau/5}$ and
$R^{a-5/\tau}\norm{g}_{\mathcal{M}^{P,\tau}}$ equals the same quantity. Adding
the two bounds gives \eqref{eq:hedberg} with $C_{20}=C+C''$.
\end{proof}

\begin{corollary}[Adams' theorem]\label{cor:adams}
Let $1<P\le\tau<\infty$ and $0<a<5/\tau$, and set $\varkappa=a\tau/5$,
$\lambda=1-\varkappa$, so that
\begin{equation}\label{eq:adams-exponents}
  \frac{1}{\tau/\lambda}=\frac1\tau-\frac a5 .
\end{equation}
There is $C_{21}(a,P,\tau)$ such that
\begin{equation}\label{eq:adams}
  \norm{I_{a}g}_{\mathcal{M}^{P/\lambda,\;\tau/\lambda}}
  \;\le\;C_{21}(a,P,\tau)\,\norm{g}_{\mathcal{M}^{P,\tau}}
  \qquad\text{for all }g\in\mathcal{M}^{P,\tau}.
\end{equation}
\end{corollary}

\begin{proof}
Write $\Lambda=\norm{g}_{\mathcal{M}^{P,\tau}}$ and let $z_{0}$, $r>0$. By
Lemma~\ref{lem:hedberg},
$\abs{I_{a}g}^{P/\lambda}\le C_{20}^{P/\lambda}
  (\mathcal{M}g)^{P}\Lambda^{\varkappa P/\lambda}$, since
$(1-\varkappa)P/\lambda=P$. Therefore
\[
  \iint_{\mathfrak{B}_{r}(z_{0})}\abs{I_{a}g}^{P/\lambda}
  \le C_{20}^{P/\lambda}\Lambda^{\varkappa P/\lambda}
     \iint_{\mathfrak{B}_{r}(z_{0})}(\mathcal{M}g)^{P} .
\]
It therefore suffices to prove the Morrey bound
\begin{equation}\label{eq:maximal-morrey}
  \iint_{\mathfrak{B}_{r}(z_{0})}(\mathcal{M}g)^{P}
  \;\le\;C\,r^{5(1-\frac P\tau)}\Lambda^{P}.
\end{equation}
Split $g=g_{1}+g_{2}$ with $g_{1}=g\mathbf{1}_{\mathfrak{B}_{2r}(z_{0})}$. For
$g_{1}$, External Input~\ref{ext:maximal} gives
$\iint(\mathcal{M}g_{1})^{P}\le C_{19}(P)^{P}\iint\abs{g_{1}}^{P}
 \le C_{19}(P)^{P}(2r)^{5(1-\frac P\tau)}\Lambda^{P}$.
For $g_{2}$ and $z\in\mathfrak{B}_{r}(z_{0})$, any ball
$\mathfrak{B}_{\varrho}(z)$ meeting $\supp g_{2}$ has $\varrho>r$, and then
$\mathfrak{B}_{\varrho}(z)\subset\mathfrak{B}_{2\varrho}(z_{0})$, so
\[
  \mathcal{M}g_{2}(z)\le\sup_{\varrho>r}
    \frac{1}{\tfrac{8\pi}{3}\varrho^{5}}\iint_{\mathfrak{B}_{2\varrho}(z_{0})}\abs{g}
  \le\sup_{\varrho>r}\frac{C}{\varrho^{5}}(2\varrho)^{5(1-\frac1\tau)}\Lambda
  = C'\,r^{-5/\tau}\Lambda ,
\]
using the H\"older step of the previous proof. Hence
$\iint_{\mathfrak{B}_{r}(z_{0})}(\mathcal{M}g_{2})^{P}
 \le\tfrac{8\pi}{3}r^{5}(C'r^{-5/\tau}\Lambda)^{P}
 =C''r^{5(1-\frac P\tau)}\Lambda^{P}$. Adding the two contributions, with the factor $2^{P-1}$ from
$(s_{1}+s_{2})^{P}\le2^{P-1}(s_{1}^{P}+s_{2}^{P})$ absorbed into $C$, gives
\eqref{eq:maximal-morrey}, and therefore
$\iint_{\mathfrak{B}_{r}(z_{0})}\abs{I_{a}g}^{P/\lambda}
 \le C\,r^{5(1-\frac P\tau)}\Lambda^{P/\lambda}$. Since
$5\bigl(1-\frac{P/\lambda}{\tau/\lambda}\bigr)=5\bigl(1-\frac P\tau\bigr)$, this
is exactly \eqref{eq:adams}.
\end{proof}

\begin{remark}\label{rem:adams-source}
Lemma~\ref{lem:hedberg} and Corollary~\ref{cor:adams} are
\cite[Lemma 5.3 and Corollary 5.1, p.~93--94]{LemarieRieusset2016}, stated
there for a general space of homogeneous type of homogeneous dimension $Q$; we
have specialized to $Q=5$ and written the proof out. Corollary~\ref{cor:adams}
is due to Adams \cite{Adams1975}. Note that only the \emph{Morrey} exponent
$\tau$ enters the condition $a\tau<5$ and the gain
\eqref{eq:adams-exponents}; the integrability exponent $P$ is carried along.
\end{remark}

\begin{lemma}[Parabolic Campanato lemma]\label{lem:campanato}
Let $P\in[1,\infty)$, $\gamma\in(0,1)$ and $h\in L^{P}_{loc}(\R^{3}\times\R)$,
and suppose
\[
  H\;=\;\sup_{z}\sup_{r>0}\ \frac{1}{r^{\gamma}}
    \Bigl(\frac{1}{\abs{\mathfrak{B}_{r}(z)}}
      \iint_{\mathfrak{B}_{r}(z)}
        \bigl|h-\textstyle\stavg{h}{\mathfrak{B}_{r}(z)}\bigr|^{P}\Bigr)^{1/P}
  \;<\;\infty ,
\]
where $\stavg{h}{\mathfrak{B}_{r}(z)}$ is the average of $h$ over
$\mathfrak{B}_{r}(z)$. Then $h$ agrees almost everywhere with a function
$\bar h$ with
$\abs{\bar h(z)-\bar h(z')}\le C_{22}(\gamma,P)\,H\,\dpar(z,z')^{\gamma}$ for all
$z,z'$, and moreover, for every $z$ and every $R>0$,
\begin{equation}\label{eq:campanato-Linfty}
  \norm{\bar h}_{L^{\infty}(\mathfrak{B}_{R}(z))}
  \;\le\;2^{\gamma}C_{22}(\gamma,P)\,H\,R^{\gamma}
   +\Bigl(\frac1{\abs{\mathfrak{B}_{R}(z)}}
      \iint_{\mathfrak{B}_{R}(z)}\abs{h}^{P}\Bigr)^{1/P}.
\end{equation}
Conversely, if $h$ satisfies such a H\"older bound with constant $[h]$ then
$H\le2^{\gamma}[h]$.
\end{lemma}

\begin{proof}
The converse is immediate: for $w\in\mathfrak{B}_{r}(z)$,
\[
\abs{h(w)-\stavg{h}{\mathfrak{B}_{r}(z)}}
 \le\abs{\mathfrak{B}_{r}(z)}^{-1}\iint_{\mathfrak{B}_{r}(z)}\abs{h(w)-h(w')}dw'
 \le[h]\,(2r)^{\gamma}
,
\] since $\dpar(w,w')<2r$ for $w,w'\in\mathfrak{B}_{r}(z)$.

For the direct statement, fix $\Phi\in C_{c}^{\infty}(\R^{3}\times\R)$ supported
in $\mathfrak{B}_{1}(0)$ with $\iint\Phi=1$, and set
$\Phi_{\epsilon}(x,t)=\epsilon^{-5}\Phi(x/\epsilon,t/\epsilon^{2})$, so that
$\Phi_{\epsilon}$ is supported in $\mathfrak{B}_{\epsilon}(0)$,
$\iint\Phi_{\epsilon}=1$,
$\norm{\Phi_{\epsilon}}_{\infty}\le\epsilon^{-5}\norm{\Phi}_{\infty}$,
$\norm{\nabla_{x}\Phi_{\epsilon}}_{\infty}\le\epsilon^{-6}\norm{\nabla\Phi}_{\infty}$
and
$\norm{\partial_{t}\Phi_{\epsilon}}_{\infty}\le\epsilon^{-7}\norm{\partial_{t}\Phi}_{\infty}$.

\emph{(a) Successive scales.} Since $\iint(\Phi_{\epsilon}-\Phi_{\epsilon/2})=0$
and both are supported in $\mathfrak{B}_{\epsilon}(0)$,
\[
  \Phi_{\epsilon}*h(z)-\Phi_{\epsilon/2}*h(z)
  =\iint_{\mathfrak{B}_{\epsilon}(z)}
    \bigl(\Phi_{\epsilon}(z-w)-\Phi_{\epsilon/2}(z-w)\bigr)
    \bigl(h(w)-\stavg{h}{\mathfrak{B}_{\epsilon}(z)}\bigr)dw ,
\]
so, bounding the kernel by
$C\epsilon^{-5}\le C'\abs{\mathfrak{B}_{\epsilon}(z)}^{-1}$ and using
H\"older's inequality and the hypothesis,
\begin{equation}\label{eq:campanato-scales}
  \abs{\Phi_{\epsilon}*h(z)-\Phi_{\epsilon/2}*h(z)}\le C H\epsilon^{\gamma}.
\end{equation}

\emph{(b) Two points at one scale.} Let $z,z'$ with
$\varrho=\dpar(z,z')\le\epsilon$. Since
$\iint(\Phi_{\epsilon}(z-\cdot)-\Phi_{\epsilon}(z'-\cdot))=0$ and both kernels
are supported in $\mathfrak{B}_{\epsilon+\varrho}(z)\subset\mathfrak{B}_{2\epsilon}(z)$,
\[
  \Phi_{\epsilon}*h(z)-\Phi_{\epsilon}*h(z')
  =\iint_{\mathfrak{B}_{2\epsilon}(z)}
    \bigl(\Phi_{\epsilon}(z-w)-\Phi_{\epsilon}(z'-w)\bigr)
    \bigl(h(w)-\stavg{h}{\mathfrak{B}_{2\epsilon}(z)}\bigr)dw .
\]
Writing $z=(x,t)$, $z'=(x',t')$ and using the mean value theorem separately in
$x$ and in $t$,
\[
  \abs{\Phi_{\epsilon}(z-w)-\Phi_{\epsilon}(z'-w)}
  \le\frac{\abs{x-x'}}{\epsilon^{6}}\norm{\nabla\Phi}_{\infty}
    +\frac{\abs{t-t'}}{\epsilon^{7}}\norm{\partial_{t}\Phi}_{\infty}
  \le\frac{C}{\epsilon^{5}}\cdot\frac{\varrho}{\epsilon},
\]
because $\abs{x-x'}\le\varrho$ and $\abs{t-t'}\le\varrho^{2}\le\varrho\epsilon$.
Hence
\begin{equation}\label{eq:campanato-points}
  \abs{\Phi_{\epsilon}*h(z)-\Phi_{\epsilon}*h(z')}
  \le C H\epsilon^{\gamma}\frac{\varrho}{\epsilon}.
\end{equation}

\emph{(c) A continuous representative.} By \eqref{eq:campanato-scales}, for
$0<\epsilon'\le\epsilon$ with $\epsilon'=2^{-m}\epsilon$,
\begin{equation}\label{eq:campanato-cauchy}
  \abs{\Phi_{\epsilon}*h(z)-\Phi_{\epsilon'}*h(z)}
  \le CH\sum_{k=0}^{m-1}(2^{-k}\epsilon)^{\gamma}
  \le\frac{CH}{1-2^{-\gamma}}\,\epsilon^{\gamma},
\end{equation}
uniformly in $z$, so $(\Phi_{2^{-k}}*h)_{k\ge0}$ has a uniform limit
$\bar h$. Repeat part (a) with $\epsilon'\in[\epsilon/2,\epsilon]$
in place of $\epsilon/2$. Both kernels still have integral one,
support in $\mathfrak B_\epsilon(0)$, and supremum bounded by
$C\epsilon^{-5}$. Consequently
\[
 \sup_z|\Phi_\epsilon*h(z)-\Phi_{\epsilon'}*h(z)|
       \le CH\epsilon^\gamma
 \qquad(\epsilon/2\le\epsilon'\le\epsilon).
\]
For any small $s>0$, choose $k$ with $2^{-k-1}\le s\le2^{-k}$
and compare $\Phi_s*h$ to $\Phi_{2^{-k}}*h$. This comparison tends to
zero uniformly as $s\downarrow0$, proving
$\bar h=\lim_{s\downarrow0}\Phi_s*h$. Telescoping from an arbitrary
initial scale $\epsilon$ now gives
$\sup_z|\bar h(z)-\Phi_\epsilon*h(z)|\le C'H\epsilon^\gamma$.
The limit is continuous because each mollification is continuous.
By the Lebesgue differentiation theorem on the homogeneous space
$(\R^{3}\times\R,\dpar,dz)$ (a consequence of
External Input~\ref{ext:maximal}), $\Phi_{\epsilon}*h(z)\to h(z)$ at every
Lebesgue point of $h$, hence $\bar h=h$ almost everywhere.

\emph{(d) The H\"older bound.} Let $z\ne z'$ and $\varrho=\dpar(z,z')$. By
\eqref{eq:campanato-cauchy} with $\epsilon=\varrho$ and $\epsilon'\to0$,
$\abs{\bar h(z)-\Phi_{\varrho}*h(z)}\le C'H\varrho^{\gamma}$ and likewise at
$z'$; by \eqref{eq:campanato-points} with $\epsilon=\varrho$,
$\abs{\Phi_{\varrho}*h(z)-\Phi_{\varrho}*h(z')}\le CH\varrho^{\gamma}$. Adding
the three gives $\abs{\bar h(z)-\bar h(z')}\le C_{22}(\gamma,P)H\varrho^{\gamma}$.

\emph{(e) The supremum bound \eqref{eq:campanato-Linfty}.} For
$w\in\mathfrak{B}_{R}(z)$,
\[
  \bigl|\bar h(w)-\stavg{h}{\mathfrak{B}_{R}(z)}\bigr|
  =\Bigl|\frac1{\abs{\mathfrak{B}_{R}(z)}}\iint_{\mathfrak{B}_{R}(z)}
     \bigl(\bar h(w)-\bar h(w')\bigr)dw'\Bigr|
  \le C_{22}(\gamma,P)H(2R)^{\gamma},
\]
because $\dpar(w,w')<2R$, and
$\bigl|\stavg{h}{\mathfrak{B}_{R}(z)}\bigr|
 \le\bigl(\abs{\mathfrak{B}_{R}(z)}^{-1}\iint_{\mathfrak{B}_{R}(z)}\abs{h}^{P}\bigr)^{1/P}$
by Jensen's inequality.
\end{proof}

\begin{proposition}[Heat potentials of Morrey sources are H\"older]
\label{prop:heat-morrey-hoelder}
Let $\gamma\in(0,1)$ and define
\begin{equation}\label{eq:q0q1}
  \frac1{\theta_{0}}=\frac{2-\gamma}{5},\qquad
  \frac1{\theta_{1}}=\frac{1-\gamma}{5},
  \qquad\text{so}\qquad \theta_{0}=\frac{5}{2-\gamma}>\tfrac52,
  \quad\theta_{1}=\frac{5}{1-\gamma}>5 .
\end{equation}
Let $1\le P\le\theta_{0}<\theta_{1}<\infty$, let $K\in\N$, let
$\varsigma_{1},\dots,\varsigma_{K}$ be smooth on $\R^{3}\setminus\{0\}$ and
homogeneous of degree $1$, and let $\varsigma_{k}(D)$ be the associated Fourier
multipliers, acting in the space variable only. Let
$F\in\mathcal{M}^{P,\theta_{0}}$ and $G_{1},\dots,G_{K}\in\mathcal{M}^{P,\theta_{1}}$,
all supported in a fixed bounded set, and put
\begin{equation}\label{eq:heat-potential}
  h \;=\; W_{+}*\Bigl(F+\sum_{k=1}^{K}\varsigma_{k}(D)G_{k}\Bigr),
  \qquad W_{+}(x,t)=\mathbf{1}_{t>0}\,(4\pi t)^{-3/2}e^{-\abs{x}^{2}/(4t)} ,
\end{equation}
the convolution being taken in $\R^{3}\times\R$ in the distributional,
equivalently almost-everywhere kernel, sense. In particular, its slice
formula is
\[
 \begin{aligned}
 h(x,t)&=\int_{\R}\bigl(W_+(\cdot,t-s)*_x F(\cdot,s)\bigr)(x)\,ds\\
 &\quad+\sum_k\int_{\R}
   \bigl(\varsigma_k(D)(W_+(\cdot,t-s)*_xG_k(\cdot,s))\bigr)(x)\,ds
 \end{aligned}
 \qquad\text{a.e.}
\]
The translated kernel products are absolutely integrable almost everywhere;
this is a conclusion, not an additional source hypothesis. The
distributional pairing remains an identity for every compactly supported
smooth test. The finite constant $C_{23}$ below is chosen from
$\gamma,P$ and the finite symbol family before the sources or their support.
Then $h$ agrees a.e.\ with a
function $\bar h$ that is H\"older of exponent $\gamma$ with respect to
$\dpar$, with
\begin{equation}\label{eq:heat-hoelder}
  [\bar h]_{C^{\gamma,\gamma/2}}\le C_{23}(\gamma,P,\varsigma_{1},\dots,\varsigma_{K})
     \Bigl(\norm{F}_{\mathcal{M}^{P,\theta_{0}}}
          +\sum_{k}\norm{G_{k}}_{\mathcal{M}^{P,\theta_{1}}}\Bigr)=:C_{23}\mathcal{N},
\end{equation}
and, for every $z$ and $R>0$,
\begin{equation}\label{eq:heat-Linfty}
  \norm{\bar h}_{L^{\infty}(\mathfrak{B}_{R}(z))}
  \le 2^{\gamma}C_{23}\mathcal{N}\,R^{\gamma}
   +\Bigl(\frac1{\abs{\mathfrak{B}_{R}(z)}}\iint_{\mathfrak{B}_{R}(z)}
      \abs{h}^{P}\Bigr)^{1/P}.
\end{equation}
\end{proposition}

\begin{proof}
By External Input~\ref{ext:heat-kernel} the kernels satisfy, for
$w=(y,s)$ with $s>0$ and $\varrho=\dpar(w,0)$,
\begin{equation}\label{eq:heat-kernel-bounds}
 \begin{aligned}
  \abs{W_{+}}\le C\varrho^{-3},\quad
  \abs{\nabla_{y}W_{+}}\le C\varrho^{-4},\quad
  \abs{\partial_{s}W_{+}}\le C\varrho^{-5},\\
  \abs{\varsigma(D)W_{+}}\le C\varrho^{-4},\quad
  \abs{\nabla_{y}\varsigma(D)W_{+}}\le C\varrho^{-5},\quad
  \abs{\partial_{s}\varsigma(D)W_{+}}\le C\varrho^{-6}.
 \end{aligned}
\end{equation}
The causal kernels vanish for $s<0$; their values on $s=0$ are set to zero
and do not affect convolution integrals. A general multiplier kernel need
not have zero right trace at $s=0$, even when $y\ne0$. Accordingly no
classical time derivative or cross-zero mean-value estimate is asserted
there.
By linearity it suffices to treat $K=1$ and to sum the resulting bounds, so we
write $\varsigma=\varsigma_{1}$, $G=G_{1}$. Note that $\varsigma(D)W_{+}$ is a
locally integrable function by \eqref{eq:heat-kernel-bounds}
($\iint_{\{\varrho<R\}}\varrho^{-4}=\tfrac{40\pi}{3}R<\infty$), so
$W_{+}*\varsigma(D)G=(\varsigma(D)W_{+})*G$.

Fix $z$ and $r>0$. Split $h=h^{\mathrm{near}}+\sum_{j\ge6}h_{j}$, where
\[
  h^{\mathrm{near}}=W_{+}*\bigl(F^{\mathrm{near}}+\varsigma(D)G^{\mathrm{near}}\bigr),
  \quad
  F^{\mathrm{near}}=\mathbf{1}_{\mathfrak{B}_{2^{6}r}(z)}F,
  \quad
  G^{\mathrm{near}}=\mathbf{1}_{\mathfrak{B}_{2^{6}r}(z)}G,
\]
and, for $j\ge6$,
$\Gamma_{j}=\mathfrak{B}_{2^{j+1}r}(z)\setminus\mathfrak{B}_{2^{j}r}(z)$,
$F_{j}=\mathbf{1}_{\Gamma_{j}}F$, $G_{j}=\mathbf{1}_{\Gamma_{j}}G$,
$h_{j}=W_{+}*(F_{j}+\varsigma(D)G_{j})$. The sum over $j\ge6$ is finite because
$F$ and $G$ have bounded support. We estimate
\begin{equation}\label{eq:campanato-quantity}
  \mathcal{E}_{j}(z,r)
  =\Bigl(\frac{1}{\abs{\mathfrak{B}_{r}(z)}^{2}}
    \iint_{\mathfrak{B}_{r}(z)}\iint_{\mathfrak{B}_{r}(z)}
      \abs{h_{j}(w)-h_{j}(w')}^{P}\,dw\,dw'\Bigr)^{1/P}.
\end{equation}

\emph{Near part, treated as one block.} No cancellation is used. For
$w\in\mathfrak{B}_{r}(z)$ and $v\in\mathfrak{B}_{2^{6}r}(z)$ we have
$\dpar(w,v)<2^{7}r$, so Young's inequality with the truncated kernels gives
\[
  \norm{h^{\mathrm{near}}}_{L^{P}(\mathfrak{B}_{r}(z))}
  \le\norm{\mathbf{1}_{\varrho<2^{7}r}\varrho^{-3}}_{L^{1}}
     \norm{F^{\mathrm{near}}}_{L^{P}}
   +C_{\varsigma}\norm{\mathbf{1}_{\varrho<2^{7}r}\varrho^{-4}}_{L^{1}}
     \norm{G^{\mathrm{near}}}_{L^{P}} .
\]
Since $\abs{\mathfrak{B}_{s}}=\tfrac{8\pi}{3}s^{5}$, we have
$\iint_{\{\varrho<R\}}\varrho^{-3}=\tfrac{40\pi}{3}\int_{0}^{R}s\,ds=\tfrac{20\pi}{3}R^{2}$
and $\iint_{\{\varrho<R\}}\varrho^{-4}=\tfrac{40\pi}{3}R$, so the two kernel
norms are $Cr^{2}$ and $Cr$. By Definition~\ref{def:parabolic-morrey},
$\norm{F^{\mathrm{near}}}_{L^{P}}\le(2^{6}r)^{5(\frac1P-\frac1{\theta_{0}})}
 \norm{F}_{\mathcal{M}^{P,\theta_{0}}}$ and
$\norm{G^{\mathrm{near}}}_{L^{P}}\le(2^{6}r)^{5(\frac1P-\frac1{\theta_{1}})}
 \norm{G}_{\mathcal{M}^{P,\theta_{1}}}$. Using
$\mathcal{E}^{\mathrm{near}}\le2\abs{\mathfrak{B}_{r}}^{-1/P}
 \norm{h^{\mathrm{near}}}_{L^{P}(\mathfrak{B}_{r}(z))}$,
$\abs{\mathfrak{B}_{r}}^{-1/P}=Cr^{-5/P}$, and
$2-5/\theta_{0}=1-5/\theta_{1}=\gamma$ from \eqref{eq:q0q1},
\[
  \mathcal{E}^{\mathrm{near}}(z,r)\le
    C\,r^{\gamma}\bigl(\norm{F}_{\mathcal{M}^{P,\theta_{0}}}
      +\norm{G}_{\mathcal{M}^{P,\theta_{1}}}\bigr).
\]
No series and no constraint beyond $P\le\theta_{0}$ is involved.

\emph{Far shells, $j\ge6$.} Write $L=2^jr$, let
$w,w'\in\mathfrak B_r(z)$ and split $\Gamma_j$ into the regions where
$t_w-s$ and $t_{w'}-s$ are both positive, both nonpositive, or have opposite
signs. Spatial separation along the intervening segments, or time separation
when it is the larger coordinate, is comparable to $L$.
On the first region the separate space and time mean-value estimates give
\[
 \begin{aligned}
 |W_+(w-v)-W_+(w'-v)|&\le C\,\dpar(w,w')L^{-4},\\
 |(\varsigma(D)W_+)(w-v)-(\varsigma(D)W_+)(w'-v)|
       &\le C\,\dpar(w,w')L^{-5}.
 \end{aligned}
\]
The contribution from the second region vanishes. By H\"older and the
Morrey bounds, the first region contributes at most
\[
 C\,r L^{\gamma-1}
   \bigl(\norm F_{\mathcal M^{P,\theta_0}}
                   +\norm G_{\mathcal M^{P,\theta_1}}\bigr).
\]
Here $1-5/\theta_0=\gamma-1$ and $-5/\theta_1=\gamma-1$.

On the crossing region the source time $s$ lies between $t_w$ and $t_{w'}$,
an interval of length at most $2r^2$. Since $j\ge6$, membership in
$\Gamma_j$ then forces spatial separation comparable to $L$.
Cover this spatial shell and time interval by at most $C(L/r)^3$
parabolic balls of radius $Cr$. H\"older on each such ball gives
\[
 \iint_{\text{crossing region}}|G|
 \le C(L/r)^3 r^{5-5/\theta_1}
             \norm G_{\mathcal M^{P,\theta_1}}
 = C L^3r^{1+\gamma}\norm G_{\mathcal M^{P,\theta_1}}.
\]
The kernel size bound $CL^{-4}$ therefore bounds its crossing contribution
by $C(r/L)r^\gamma\norm G_{\mathcal M^{P,\theta_1}}$.
This estimate does not differentiate through time zero.
For the Gaussian part, the zero extension and its time derivative vanish
at $t=0$ at nonzero spatial separation. Alternatively,
$W_+(x,\tau)\le C\tau L^{-5}$ for $|x|\asymp L$, $0<\tau\le2r^2$,
by the exponential formula. The same cover yields
\[
 \iint_{\text{crossing region}}|F|
       \le C L^3r^\gamma\norm F_{\mathcal M^{P,\theta_0}},
\]
so this contribution is at most
$C(r/L)^2r^\gamma\norm F_{\mathcal M^{P,\theta_0}}$.
Consequently
\[
 \mathcal E_j(z,r)\le C r^\gamma
       \bigl(2^{j(\gamma-1)}+2^{-j}+2^{-2j}\bigr)
       \bigl(\norm F_{\mathcal M^{P,\theta_0}}
                      +\norm G_{\mathcal M^{P,\theta_1}}\bigr).
\]
All three shell series converge for $0<\gamma<1$.

For comparison with a pointwise shell-difference estimate, set
$\delta=\dpar(w,w')$. The case $\delta=0$ is immediate. Otherwise repeat
the crossing cover at radius $C\delta$ instead of $Cr$; the time strip
has length at most $\delta^2$. The two crossing contributions are bounded by
$C\delta^{1+\gamma}L^{-1}\norm G_{\mathcal M^{P,\theta_1}}$ and
$C\delta^{2+\gamma}L^{-2}\norm F_{\mathcal M^{P,\theta_0}}$.
Since $\delta\le2r\le L$, each is at most the corresponding source norm
times $C\delta L^{\gamma-1}$. Together with the positive-time region,
this gives
\[
 |h_j(w)-h_j(w')|\le C\delta L^{\gamma-1}
    \bigl(\norm F_{\mathcal M^{P,\theta_0}}
                   +\norm G_{\mathcal M^{P,\theta_1}}\bigr).
\]
%
%
%
%
%

\emph{Conclusion.} By the triangle inequality in
$L^{P}(\mathfrak{B}_{r}\times\mathfrak{B}_{r})$ and Jensen's inequality,
\[
  \Bigl(\frac{1}{\abs{\mathfrak{B}_{r}(z)}}\iint_{\mathfrak{B}_{r}(z)}
     \bigl|h-\stavg{h}{\mathfrak{B}_{r}(z)}\bigr|^{P}\Bigr)^{1/P}
  \le\mathcal{E}^{\mathrm{near}}(z,r)+\sum_{j\ge6}\mathcal{E}_{j}(z,r)
  \le C\,r^{\gamma}\bigl(\norm{F}_{\mathcal{M}^{P,\theta_{0}}}
       +\norm{G}_{\mathcal{M}^{P,\theta_{1}}}\bigr),
\]
uniformly in $z$ and $r$. Lemma~\ref{lem:campanato} now gives
\eqref{eq:heat-hoelder} with $C_{23}=C_{22}(\gamma,P)\cdot C$, and
\eqref{eq:heat-Linfty} follows by averaging the resulting H\"older
bound over $\mathfrak B_R(z)$, as in part (e) of that lemma. Its coefficient
is $2^\gamma C_{23}$, since $C_{23}$ already includes $C_{22}$.
For $K>1$ apply the
above to each pair $(0,G_{k})$ and to $(F,0)$ and add the $K+1$ resulting
seminorm bounds.
\end{proof}

\begin{remark}[Averages and the continuous representative]
There is also a direct dyadic proof of Lemma~\ref{lem:campanato}, avoiding
mollification. Set $a_r(z)=\stavg h{\mathfrak B_r(z)}$.
H\"older's inequality and the volume ratio give
\[
 |a_{r/2}(z)-a_r(z)|\le 2^{5/P}H r^\gamma .
\]
Telescoping at dyadic radii gives a locally uniform limit of the averages,
equal to $h$ almost everywhere by differentiation. To compare two centres
at distance $r$, compare both radius-$r$ averages to the average on a
radius-$3r$ ball containing both balls. This gives a bound $CHr^\gamma$;
the two telescoping tails give the same bound. The resulting single
representative is H\"older everywhere. Averaging its H\"older estimate
gives \eqref{eq:campanato-Linfty}. This is an alternative to the
mollification proof above; its constants may depend on $(\gamma,P)$.
\end{remark}

\subsection{Step 1: the local energy and pressure estimates}\label{ss:step1}

All two-scale estimates in this section are applied with the closure of
the outer cylinder inside the solution domain. Cutoff coefficients are
fixed before the solution. The absolute energy constants and the
$q$-dependent force constants may be enlarged to dominate the finitely
many explicit cutoff and norm-comparison constants, without changing
their stated dependence. The iteration uses one fixed contraction
$\kappa$ chosen in \eqref{eq:kappa-eta}; a bound proved only at that
choice is not an assertion uniform over every contraction.

\begin{lemma}[The backward Gaussian test function]\label{lem:gaussian}
Let $G(x,\vartheta)=(4\pi\vartheta)^{-3/2}e^{-\abs{x}^{2}/(4\vartheta)}$ and, for
$r>0$ and $z_{0}=(x_{0},t_{0})$, set
\begin{equation}\label{eq:psi-r}
  \psi_{r}(x,t)\;=\;r^{2}\,G\bigl(x-x_{0},\,r^{2}-(t-t_{0})\bigr),
  \qquad (x,t)\in\R^{3}\times(-\infty,t_{0}+r^{2}).
\end{equation}
Then $\psi_{r}$ is smooth and positive there, and there is an absolute constant
$C_{24}\ge1$ such that, for all $0<r\le\rho/2$,
\begin{align}
  \partial_{t}\psi_{r}+\Delta\psi_{r}&=0
    &&\text{on }\R^{3}\times(-\infty,t_{0}+r^{2}),
    \label{eq:psi-backward}\\
  \psi_{r}&\ge C_{24}^{-1}r^{-1} &&\text{on }\Cyl{r}{z_{0}},
    \label{eq:psi-lower}\\
  \psi_{r}\le C_{24}r^{-1},\qquad
  \abs{\nabla\psi_{r}}&\le C_{24}r^{-2}
    &&\text{on }\R^{3}\times(-\infty,t_{0}],
    \label{eq:psi-upper}\\
  \psi_{r}\le C_{24}\frac{r^{2}}{\rho^{3}},\qquad
  \abs{\nabla\psi_{r}}&\le C_{24}\frac{r^{2}}{\rho^{4}}
    &&\text{on }\Cyl{\rho}{z_{0}}\setminus\Cyl{\rho/2}{z_{0}} .
    \label{eq:psi-far}
\end{align}
\end{lemma}

\begin{proof}
We may take $z_{0}=(0,0)$. Write $\vartheta=r^{2}-t>0$ for $t<r^{2}$ and
$s=\abs{x}^{2}/(4\vartheta)$. By External Input~\ref{ext:heat},
$\partial_{\vartheta}G=\Delta G$, so
$\partial_{t}\psi_{r}=-r^{2}(\partial_{\vartheta}G)(x,\vartheta)
 =-r^{2}(\Delta G)(x,\vartheta)=-\Delta\psi_{r}$, which is
\eqref{eq:psi-backward}.

\emph{\eqref{eq:psi-lower}.} On $Q_{r}$ we have $\abs{x}<r$ and
$t\in(-r^{2},0]$, so $\vartheta\in[r^{2},2r^{2}]$ and $s\le r^{2}/(4r^{2})=1/4$.
Hence
$\psi_{r}\ge r^{2}(8\pi r^{2})^{-3/2}e^{-1/4}=(8\pi)^{-3/2}e^{-1/4}r^{-1}$.

\emph{\eqref{eq:psi-upper}.} For $t\le0$, $\vartheta\ge r^{2}$, so
$\psi_{r}\le r^{2}(4\pi\vartheta)^{-3/2}\le(4\pi)^{-3/2}r^{-1}$. For the
gradient, $\nabla\psi_{r}=-r^{2}\frac{x}{2\vartheta}G(x,\vartheta)$, and
$\frac{\abs{x}}{2\vartheta}=\vartheta^{-1/2}s^{1/2}$, so
\begin{equation}\label{eq:grad-psi}
  \abs{\nabla\psi_{r}}
  =r^{2}(4\pi)^{-3/2}\vartheta^{-2}\,s^{1/2}e^{-s}
  \le r^{2}(4\pi)^{-3/2}(2e)^{-1/2}\vartheta^{-2},
\end{equation}
using $\sup_{s\ge0}s^{1/2}e^{-s}=(2e)^{-1/2}$. With $\vartheta\ge r^{2}$ this
gives $\abs{\nabla\psi_{r}}\le(4\pi)^{-3/2}(2e)^{-1/2}r^{-2}$.

\emph{\eqref{eq:psi-far}.} Let $(x,t)\in Q_{\rho}\setminus Q_{\rho/2}$, so
$t\in(-\rho^{2},0]$ and either (A) $t\le-\rho^{2}/4$, or (B)
$\abs{x}\ge\rho/2$.

In case (A), $\vartheta=r^{2}-t\ge\rho^{2}/4$, hence
$\psi_{r}\le r^{2}(4\pi\vartheta)^{-3/2}\le r^{2}(\pi\rho^{2})^{-3/2}
 =\pi^{-3/2}r^{2}\rho^{-3}$, and by \eqref{eq:grad-psi}
$\abs{\nabla\psi_{r}}\le(4\pi)^{-3/2}(2e)^{-1/2}r^{2}(4/\rho^{2})^{2}
 =16(4\pi)^{-3/2}(2e)^{-1/2}r^{2}\rho^{-4}$.

In case (B), put $\varsigma=\rho^{2}/(16\vartheta)$, so that
$s=\abs{x}^{2}/(4\vartheta)\ge\varsigma$ and $\vartheta=\rho^{2}/(16\varsigma)$.
Then
\[\begin{aligned}
  &\psi_{r}\le r^{2}(4\pi\vartheta)^{-3/2}e^{-\varsigma}
  =r^{2}(4\pi)^{-3/2}\Bigl(\frac{16\varsigma}{\rho^{2}}\Bigr)^{3/2}e^{-\varsigma}
  \\ &=64(4\pi)^{-3/2}\frac{r^{2}}{\rho^{3}}\,\varsigma^{3/2}e^{-\varsigma}
  \le64(4\pi)^{-3/2}\Bigl(\tfrac32\Bigr)^{3/2}e^{-3/2}\frac{r^{2}}{\rho^{3}},
\end{aligned}\]
using $\sup_{\varsigma\ge0}\varsigma^{3/2}e^{-\varsigma}=(3/2)^{3/2}e^{-3/2}$.
For the gradient, \eqref{eq:grad-psi} and
$s^{1/2}e^{-s}\le e^{-1/2}e^{-s/2}\le e^{-1/2}e^{-\varsigma/2}$ (valid because
$\sup_{s}s^{1/2}e^{-s/2}=e^{-1/2}$) give
\[
  \abs{\nabla\psi_{r}}\le r^{2}(4\pi)^{-3/2}e^{-1/2}
    \Bigl(\frac{16\varsigma}{\rho^{2}}\Bigr)^{2}e^{-\varsigma/2}
  =256(4\pi)^{-3/2}e^{-1/2}\frac{r^{2}}{\rho^{4}}\varsigma^{2}e^{-\varsigma/2}
  \le C\frac{r^{2}}{\rho^{4}} .
\]
Taking $C_{24}$ to be the largest of the finitely many constants produced in
\eqref{eq:psi-upper}--\eqref{eq:psi-far} and of
$\bigl((8\pi)^{-3/2}e^{-1/4}\bigr)^{-1}$, which is what makes
\eqref{eq:psi-lower} read $\psi_{r}\ge C_{24}^{-1}r^{-1}$, completes the proof.
\end{proof}

\begin{lemma}[Caccioppoli inequality; \protect{\cite[(13.30)]{LemarieRieusset2016}}]
\label{lem:caccioppoli}
There is an absolute constant $C_{25}$ and a constant $C_{26}=C_{26}(q)$ such
that for every suitable weak solution $(u,p)$ on $\mathcal{O}$ with force
$f\in L^{q}_{loc}$, $q>5/2$, every $z_{0}$ and $\rho>0$ with
$\Cyl{\rho}{z_{0}}\subset\mathcal{O}$, and every $0<r\le\rho/2$, writing
$\kappa=r/\rho$ and suppressing the base point $z_{0}$,
\begin{equation}\label{eq:caccioppoli}\begin{aligned}
  &\alpha(r)+\beta(r)\;\le\;
    C_{25}\kappa\,\alpha(\rho)
   \\ &\quad+C_{25}\kappa^{-1}\alpha(\rho)^{1/2}\beta(\rho)^{1/2}\gamma(\rho)^{1/2}
   \\ &\quad+C_{25}\kappa^{-1}\delta(\rho)\gamma(\rho)^{1/2}
   +C_{26}\kappa^{-1/2}\gamma(\rho)^{1/2}\lambda(\rho)^{1/2}.
\end{aligned}\end{equation}
\end{lemma}

\begin{proof}
Write $B_{\varrho}=\Ball{\varrho}{x_{0}}$ and $J_{\rho}=(t_{0}-\rho^{2},t_{0})$.
Since $t_{0}\in I$ and $I$ is open, there is $\epsilon_{0}>0$ with
$[t_{0}-\rho^{2}+\tfrac12\epsilon_{0},\,t_{0}+\epsilon_{0}]\subset I$; shrinking
$\epsilon_{0}$ we may and do also require
\begin{equation}\label{eq:eps0-small}
  \epsilon_{0}<\tfrac12 r^{2},
\end{equation}
which is legitimate because $r>0$ is fixed before $\epsilon_{0}$ is chosen and
only smallness of $\epsilon_{0}$ is used. By \eqref{eq:eps0-small} the function
$\psi_{r}$ of Lemma~\ref{lem:gaussian} is smooth on a neighbourhood of
$\overline{\Ball{\rho}{x_{0}}}\times[t_{0}-\rho^{2},t_{0}+\epsilon_{0}]$, since
that set lies in $\R^{3}\times(-\infty,t_{0}+r^{2})$.

\emph{Step 1: the test function.} Let $\eta_{1}\in C_{c}^{\infty}(B_{\rho})$ with
$0\le\eta_{1}\le1$, $\eta_{1}\equiv1$ on $B_{\rho/2}$,
$\supp\eta_{1}\subset B_{3\rho/4}$ and
$\abs{\partial^{a}\eta_{1}}\le C_{10}(\abs{a})\rho^{-\abs{a}}$; this is
Lemma~\ref{lem:cutoff} verbatim, up to renaming the two radii
$13\rho/20\rightsquigarrow\rho/2$ and $3\rho/4\rightsquigarrow3\rho/4$, which
only changes the numerical constants. Let $\eta_{2}\in C^{\infty}(\R)$ with
$0\le\eta_{2}\le1$, $\eta_{2}\equiv0$ on
$(-\infty,t_{0}-\rho^{2}+\tfrac14\epsilon_{0}]$, $\eta_{2}\equiv1$ on
$[t_{0}-\rho^{2}/4,\,t_{0}+\tfrac12\epsilon_{0}]$, $\supp\eta_{2}\subset
(t_{0}-\rho^{2}+\tfrac18\epsilon_{0},\,t_{0}+\epsilon_{0})$ and
$\abs{\eta_{2}'}+\rho^{2}\abs{\eta_{2}''}\le C\rho^{-2}$ \emph{on}
$(-\infty,t_{0}]$, where the transition occupies the interval
$[t_{0}-\rho^{2}+\tfrac14\epsilon_{0},t_{0}-\rho^{2}/4]$ of length
$\ge\rho^{2}/2$; such $\eta_{2}$ exists by mollifying a piecewise affine
function as in Lemma~\ref{lem:cutoff}. Put $\eta=\eta_{1}\eta_{2}$ and
\[
  \phi=\psi_{r}\,\eta\in C_{c}^{\infty}(\mathcal{O}),\qquad \phi\ge0 ,
\]
with $\psi_{r}$ as in Lemma~\ref{lem:gaussian}. Then $\eta\equiv1$ on
$\Cyl{\rho/2}{z_{0}}$, and $\nabla\eta$, $\partial_{t}\eta$, $\Delta\eta$ all
vanish on $\Cyl{\rho/2}{z_{0}}$ and, for $t\le t_{0}$, are supported in
$\Cyl{\rho}{z_{0}}\setminus\Cyl{\rho/2}{z_{0}}$.

\emph{Step 2: the four terms.} Apply Lemma~\ref{lem:lei-pointwise} with this
$\phi$ and $t\in(t_{0}-r^{2},t_{0}]$. Since
$\partial_{t}\psi_{r}+\Delta\psi_{r}=0$,
\begin{equation}\label{eq:phi-heat}
  \partial_{t}\phi+\Delta\phi
  =(\partial_{t}\eta+\Delta\eta)\psi_{r}+2\nabla\eta\cdot\nabla\psi_{r},
\end{equation}
which vanishes on $\Cyl{\rho/2}{z_{0}}$ and, by \eqref{eq:psi-far} and the
bounds on $\eta$, satisfies
$\abs{\partial_{t}\phi+\Delta\phi}\le C\rho^{-2}\cdot C_{24}r^{2}\rho^{-3}
 +2C\rho^{-1}C_{24}r^{2}\rho^{-4}\le C'r^{2}\rho^{-5}$ on
$\Cyl{\rho}{z_{0}}$, $t\le t_{0}$. Also, by \eqref{eq:psi-upper} and
$\rho\ge2r$,
\begin{equation}\label{eq:grad-phi}
  \abs{\nabla\phi}\le\eta\abs{\nabla\psi_{r}}+\psi_{r}\abs{\nabla\eta}
  \le C_{24}r^{-2}+C_{24}r^{-1}\,C\rho^{-1}\le C'r^{-2}
  \qquad\text{on }\Cyl{\rho}{z_{0}} ,
\end{equation}
and $\phi\le C_{24}r^{-1}$. Since $\supp\phi(\cdot,s)\subset\overline{B_{3\rho/4}}=:E$ for every $s$,
Lemma~\ref{lem:divfree-slice} supplies a single null set $N\subset J_{\rho}$ off
which $\int_{B_{\rho}}u(\cdot,s)\cdot\nabla\psi=0$ for \emph{all}
$\psi\in C^{1}_{c}(B_{\rho})$ supported in $E$; applying it to
$\psi=\phi(\cdot,s)$ gives
$\int_{B_{\rho}}\Gamma(s)\,u(\cdot,s)\cdot\nabla\phi(\cdot,s)\,dy=0$ for
$s\notin N$ and any function $\Gamma$ of time alone. We use this with
$\Gamma(s)=\spavg{\abs{u}^{2}}{B_{\rho}}(s)$. Therefore
Lemma~\ref{lem:lei-pointwise} gives, for almost every $t\in(t_{0}-r^{2},t_{0}]$,
\[
  \int\abs{u(\cdot,t)}^{2}\phi(\cdot,t)
  +2\int_{t_{0}-\rho^{2}}^{t}\!\!\int\abs{\nabla u}^{2}\phi
  \;\le\;I_{1}+I_{2}+I_{3}+I_{4},
\]
where
\[\begin{aligned}
  &I_{1}=\iint_{\Cyl{\rho}{z_{0}}}\abs{u}^{2}\abs{\partial_{s}\phi+\Delta\phi},
  \\
  &I_{2}=\iint_{\Cyl{\rho}{z_{0}}}
     \bigl|\abs{u}^{2}-\spavg{\abs{u}^{2}}{B_{\rho}}\bigr|\abs{u}\abs{\nabla\phi},
  \\
  &I_{3}=2\iint_{\Cyl{\rho}{z_{0}}}\abs{p}\abs{u}\abs{\nabla\phi},
  \\
  &I_{4}=2\iint_{\Cyl{\rho}{z_{0}}}\abs{f}\abs{u}\phi .
\end{aligned}\]

\emph{$I_{1}$.} By the bound on $\partial_{s}\phi+\Delta\phi$,
\[
  I_{1}\le C\frac{r^{2}}{\rho^{5}}\iint_{\Cyl{\rho}{z_{0}}}\abs{u}^{2}
  \le C\frac{r^{2}}{\rho^{5}}\rho^{2}\operatorname*{ess\,sup}_{J_{\rho}}\int_{B_{\rho}}\abs{u}^{2}
  = C\frac{r^{2}}{\rho^{3}}\,\rho\,\alpha(\rho)^{2}
  = C\kappa^{2}\alpha(\rho)^{2}.
\]

\emph{$I_{2}$.} By \eqref{eq:grad-phi} and H\"older's inequality with
$\tfrac23+\tfrac13=1$,
\[
  I_{2}\le\frac{C}{r^{2}}
    \bigl\|\abs{u}^{2}-\spavg{\abs{u}^{2}}{B_{\rho}}\bigr\|_{L^{3/2}(\Cyl{\rho}{z_{0}})}
    \norm{u}_{L^{3}(\Cyl{\rho}{z_{0}})} .
\]
For a.e.\ $s$, apply the squared-velocity Sobolev--Poincar\'e estimate
recorded after External Input~\ref{ext:poincare}. The closed outer-cylinder
containment supplies a slightly larger spatial ball on which the slice
has its $H^1$ weak gradient. The estimate gives
$\bigl\|\abs{u}^{2}-\spavg{\abs{u}^{2}}{B_{\rho}}(s)\bigr\|_{L^{3/2}(B_{\rho})}
 \le2C\norm{u(\cdot,s)}_{L^{2}(B_{\rho})}\norm{\nabla u(\cdot,s)}_{L^{2}(B_{\rho})}$.
Taking the $L^{3/2}$ norm in $s$ over $J_{\rho}$ and using H\"older in time as
in \eqref{eq:time-holder},
\[\begin{aligned}
  &\bigl\|\abs{u}^{2}-\spavg{\abs{u}^{2}}{B_{\rho}}\bigr\|_{L^{3/2}(\Cyl{\rho}{z_{0}})}
  \\ &\le C\bigl(\operatorname*{ess\,sup}\mathfrak{u}\bigr)
     \abs{J_{\rho}}^{1/6}\norm{\mathfrak{g}}_{L^{2}(J_{\rho})}
  \\ &\le C\rho^{1/2}\alpha(\rho)\cdot\rho^{1/3}\cdot\rho^{1/2}\beta(\rho)
  = C\rho^{4/3}\alpha(\rho)\beta(\rho),
\end{aligned}\]
with $\mathfrak{u},\mathfrak{g}$ as in \eqref{eq:slice-norms}. Since
$\norm{u}_{L^{3}(\Cyl{\rho}{z_{0}})}=\rho^{2/3}\gamma(\rho)$,
\[
  I_{2}\le\frac{C}{r^{2}}\rho^{4/3}\alpha(\rho)\beta(\rho)\rho^{2/3}\gamma(\rho)
  =C\frac{\rho^{2}}{r^{2}}\alpha(\rho)\beta(\rho)\gamma(\rho)
  =C\kappa^{-2}\alpha(\rho)\beta(\rho)\gamma(\rho).
\]

\emph{$I_{3}$.} By \eqref{eq:grad-phi} and H\"older,
$I_{3}\le Cr^{-2}\norm{p}_{L^{3/2}(\Cyl{\rho}{z_{0}})}\norm{u}_{L^{3}(\Cyl{\rho}{z_{0}})}
 =Cr^{-2}\rho^{4/3}\delta(\rho)^{2}\rho^{2/3}\gamma(\rho)
 =C\kappa^{-2}\delta(\rho)^{2}\gamma(\rho)$.

\emph{$I_{4}$.} Since $\phi\le C_{24}r^{-1}$ and $q>5/2$, H\"older with
$\tfrac1q+\tfrac1{q'}=1$ and then H\"older on $\Cyl{\rho}{z_{0}}$ (using
$q'<\tfrac53<3$) give
\[\begin{aligned}
  &I_{4}\le\frac{C}{r}\norm{f}_{L^{q}(\Cyl{\rho}{z_{0}})}\norm{u}_{L^{q'}(\Cyl{\rho}{z_{0}})}
  \\ &\le\frac{C}{r}\,\rho^{-\sigma}\lambda(\rho)\cdot
     \abs{\Cyl{\rho}{z_{0}}}^{\frac1{q'}-\frac13}\rho^{2/3}\gamma(\rho)
  \\ &\le\frac{C'}{r}\rho^{\frac5q-3}\lambda(\rho)\rho^{\frac{10}{3}-\frac5q}\rho^{2/3}\gamma(\rho)
  =C'\kappa^{-1}\gamma(\rho)\lambda(\rho),
\end{aligned}\]
where we used $\abs{\Cyl{\rho}{z_{0}}}^{\frac1{q'}-\frac13}
 =C\rho^{5(\frac23-\frac1q)}=C\rho^{\frac{10}3-\frac5q}$ and
$\tfrac5q-3-\sigma$ cancels.

\emph{Step 3: the left side.} On $\Cyl{r}{z_{0}}$ we have $\eta=1$ and
$\psi_{r}\ge C_{24}^{-1}r^{-1}$, so
$\int\abs{u(\cdot,t)}^{2}\phi(\cdot,t)\ge C_{24}^{-1}r^{-1}\int_{B_{r}}\abs{u(\cdot,t)}^{2}$
for $t\in(t_{0}-r^{2},t_{0}]$, and
$2\int_{t_{0}-\rho^{2}}^{t_{0}}\!\int\abs{\nabla u}^{2}\phi
 \ge2C_{24}^{-1}r^{-1}\iint_{\Cyl{r}{z_{0}}}\abs{\nabla u}^{2}
 =2C_{24}^{-1}\beta(r)^{2}$. Taking the essential supremum over the full-measure set of good times
$t\in(t_{0}-r^{2},t_{0})$ bounds the terminal energy term. For dissipation,
choose good times $t_n\uparrow t_0$, discard the nonnegative terminal
energy, and use monotone convergence on
$\mathbf1_{(t_0-\rho^2,t_n)}\abs{\nabla u}^2\phi$.
This gives the integral up to $t_0$ without assuming that $t_0$ is a good
time. Since $I_1,\ldots,I_4$ do not depend on the terminal time,
\[
  \alpha(r)^{2}+\beta(r)^{2}
  \le \tfrac32\,C_{24}\bigl(I_{1}+I_{2}+I_{3}+I_{4}\bigr)
  \le C\bigl[\kappa^{2}\alpha(\rho)^{2}+\kappa^{-2}\alpha\beta\gamma
    +\kappa^{-2}\delta^{2}\gamma+\kappa^{-1}\gamma\lambda\bigr](\rho).
\]
Taking square roots and using
$\alpha(r)+\beta(r)\le\sqrt2(\alpha(r)^{2}+\beta(r)^{2})^{1/2}$ and
$(s_{1}+s_{2}+s_{3}+s_{4})^{1/2}\le\sum s_{i}^{1/2}$ gives
\eqref{eq:caccioppoli}.
\end{proof}

\begin{remark}[Comparison with the sources]\label{rem:caccioppoli-sources}
Inequality \eqref{eq:caccioppoli} is \cite[(2.18)]{Kukavica2008} with the force
term in the $L^{q}$ form of \cite[Theorem 2]{Kukavica2008}, and it is
\cite[(13.30)]{LemarieRieusset2016} after the substitutions
$U_{\rho}=\rho\alpha(\rho)^{2}$, $V_{\rho}=\rho\beta(\rho)^{2}$,
$W_{\rho}=\rho^{2}\gamma(\rho)^{3}$, $P_{\rho}=\rho^{2}\delta(\rho)^{3}$. We
have used Kukavica's test function $\psi_{r}\eta$ rather than Lemarié-Rieusset's
$r^{3}\omega\,\theta\,H(4r^{2}+t-s,x-y)$; the two differ by the normalizing
power of $r$ (his $\psi$ is $\approx r$ times ours, matching his un-normalized
$U_{r},V_{r}$) and by the auxiliary cut-off $\theta$ in time, which is not
needed here because Lemma~\ref{lem:lei-pointwise} already provides the
pointwise-in-time inequality. The four terms and their exponents agree.
\end{remark}

\subsection{Step 2: the scale iteration}\label{ss:step2}

\begin{lemma}[Combined decay inequality]\label{lem:theta-decay}
For $\kappa\in(0,\tfrac12]$ and $\Cyl{\rho}{z}\subset\mathcal{O}$ set
\begin{equation}\label{eq:theta}
  \theta(z,\rho)\;=\;\alpha(z,\rho)+\beta(z,\rho)+\kappa^{-4}\delta(z,\rho)^{2}.
\end{equation}
There are an absolute constant $C_{27}$ and a constant $C_{28}=C_{28}(q)$ such
that for every such $z,\rho,\kappa$, suppressing $z$,
\begin{equation}\label{eq:theta-decay-1}
  \theta(\kappa\rho)\;\le\;C_{27}\kappa^{2/3}\theta(\rho)
   +C_{27}\kappa^{-5}\bigl(\beta(\rho)^{1/2}+\beta(\rho)\bigr)\theta(\rho)
   +C_{28}\kappa^{-1/2}\theta(\rho)^{1/2}\lambda(\rho)^{1/2}
   +C_{28}\kappa^{-3}\lambda(\rho),
\end{equation}
and, if in addition $\theta(\rho)\le1$,
\begin{equation}\label{eq:theta-decay-2}
  \theta(\kappa\rho)\;\le\;C_{27}\kappa^{2/3}\theta(\rho)
   +2C_{27}\kappa^{-5}\theta(\rho)^{1/2}\theta(\rho)
   +C_{28}\kappa^{-1/2}\theta(\rho)^{1/2}\lambda(\rho)^{1/2}
   +C_{28}\kappa^{-3}\lambda(\rho).
\end{equation}
\end{lemma}

\begin{proof}
Write $r=\kappa\rho$ and suppress $\rho$ on the right. All of $\alpha,\beta$ and
$\kappa^{-4}\delta^{2}$ are $\le\theta$ by definition.

\emph{Pressure part.} By Theorem~\ref{thm:pressure-decay} and
$(s_{1}+s_{2}+s_{3})^{2}\le3(s_{1}^{2}+s_{2}^{2}+s_{3}^{2})$,
\[
  \kappa^{-4}\delta(r)^{2}
  \le3C_{14}^{2}\kappa^{-5}\alpha\beta
   +3C_{14}^{2}\kappa^{2/3}\bigl(\kappa^{-4}\delta^{2}\bigr)
   +3C_{15}^{2}\kappa^{-3}\lambda
  \le3C_{14}^{2}\kappa^{-5}\beta\,\theta+3C_{14}^{2}\kappa^{2/3}\theta
   +3C_{15}^{2}\kappa^{-3}\lambda .
\]

\emph{Energy part.} By Corollary~\ref{cor:gagliardo},
$\gamma\le C_{9}(\alpha^{1/2}\beta^{1/2}+\alpha)\le C_{9}(\theta^{1/2}\beta^{1/2}+\theta)$
and hence
$\gamma^{1/2}\le C_{9}^{1/2}(\theta^{1/4}\beta^{1/4}+\theta^{1/2})$. Insert this
into \eqref{eq:caccioppoli}, term by term, using $\beta\le\theta$ and
$\delta\le\kappa^{2}\theta^{1/2}$ (from $\kappa^{-4}\delta^{2}\le\theta$):
\[
\begin{aligned}
  \kappa\alpha&\le\kappa\theta;\\
  \kappa^{-1}\alpha^{1/2}\beta^{1/2}\gamma^{1/2}
    &\le C_{9}^{1/2}\kappa^{-1}\theta^{1/2}\beta^{1/2}
      \bigl(\theta^{1/4}\beta^{1/4}+\theta^{1/2}\bigr)
    \le2C_{9}^{1/2}\kappa^{-1}\beta^{1/2}\theta;\\
  \kappa^{-1}\delta\gamma^{1/2}
    &\le\kappa\theta^{1/2}\,C_{9}^{1/2}\bigl(\theta^{1/4}\beta^{1/4}+\theta^{1/2}\bigr)
    \le2C_{9}^{1/2}\kappa\theta;\\
  \kappa^{-1/2}\gamma^{1/2}\lambda^{1/2}
    &\le C_{9}^{1/2}\kappa^{-1/2}\bigl(\theta^{1/4}\beta^{1/4}+\theta^{1/2}\bigr)\lambda^{1/2}
    \le2C_{9}^{1/2}\kappa^{-1/2}\theta^{1/2}\lambda^{1/2},
\end{aligned}
\]
where in the second line we used $\beta^{1/4}\theta^{1/4}\le\theta^{1/2}$ and in
the third and fourth $\beta^{1/4}\theta^{1/4}\le\theta^{1/2}$ again. Hence
\[
  \alpha(r)+\beta(r)\le
    C_{25}\bigl(1+2C_{9}^{1/2}\bigr)\kappa\theta
   +2C_{25}C_{9}^{1/2}\kappa^{-1}\beta^{1/2}\theta
   +2C_{26}C_{9}^{1/2}\kappa^{-1/2}\theta^{1/2}\lambda^{1/2}.
\]
Adding the two parts and using $\kappa\le\kappa^{2/3}$ and
$\kappa^{-1}\le\kappa^{-5}$, valid for $\kappa\le\tfrac12$, gives
\eqref{eq:theta-decay-1} with
$C_{27}=3C_{14}^{2}+C_{25}(1+2C_{9}^{1/2})+2C_{25}C_{9}^{1/2}$ and
$C_{28}=3C_{15}^{2}+2C_{26}C_{9}^{1/2}$ (sums, not maxima, so that nothing is
lost when the two groups are added). If $\theta(\rho)\le1$ then
$\beta\le\theta\le\theta^{1/2}$, so
$\beta^{1/2}+\beta\le2\theta^{1/2}$, which gives \eqref{eq:theta-decay-2}.
\end{proof}

\begin{convention}[The constants $\kappa$, $\eta$, $\varepsilon_{*}$]\label{conv:kappa}
Recall $\varepsilon=2/5$. Define, once and for all,
\begin{equation}\label{eq:kappa-eta}
  \kappa=\min\Bigl\{\tfrac12,\ \bigl(8C_{27}\bigr)^{-\frac{1}{2/3-\varepsilon}}\Bigr\},
  \qquad
  \eta=\min\Bigl\{1,\ \Bigl(\frac{\kappa^{5+\varepsilon}}{16\,C_{27}}\Bigr)^{2}\Bigr\},
  \qquad
  \varepsilon_{*}=\min\Bigl\{1,\ \Bigl(\frac{\kappa^{5+\varepsilon}}{32\,C_{27}}\Bigr)^{2}\Bigr\},
\end{equation}
so that
\begin{equation}\label{eq:kappa-props}
  C_{27}\kappa^{2/3-\varepsilon}\le\tfrac18,
  \qquad
  2C_{27}\kappa^{-5-\varepsilon}\eta^{1/2}\le\tfrac18,
  \qquad
  2C_{27}\kappa^{-5-\varepsilon}\bigl(\varepsilon_{*}^{1/2}+\varepsilon_{*}\bigr)\le\tfrac18 .
\end{equation}
Since $\varepsilon=2/5$ is fixed and $C_{27}$ is absolute, $\kappa$, $\eta$ and
$\varepsilon_{*}$ are \emph{absolute constants}. Finally set
\begin{equation}\label{eq:C29}
  C_{29}=C_{29}(q)=2C_{28}^{2}\kappa^{-1-2\varepsilon}+C_{28}\kappa^{-3-\varepsilon},
  \qquad
  \Lambda_{0}=\Lambda_{0}(q)=\frac{\eta}{2C_{29}} .
\end{equation}
\end{convention}

\begin{proposition}[The iteration]\label{prop:iteration}
Let $z\in\mathcal{O}$ and $r_{5}>0$ with $\Cyl{r_{5}}{z}\subset\mathcal{O}$, and
suppose
\begin{equation}\label{eq:iteration-hyp}
  \theta(z,r_{5})\le\eta
  \qquad\text{and}\qquad
  \lambda(z,r_{5})\le\Lambda_{0}.
\end{equation}
Then $\theta(z,\kappa^{n}r_{5})\le\eta\,\kappa^{n\varepsilon}$ for every
$n\in\N_{0}$, and consequently
\begin{equation}\label{eq:iteration-concl}
  \max\bigl\{\alpha(z,r),\ \beta(z,r),\ \delta(z,r)^{2}\bigr\}
  \;\le\;\theta(z,r)\;\le\;M\,r^{\varepsilon},
  \qquad 0<r\le r_{5},\qquad
  M:=\kappa^{-4/3-\varepsilon}\,\eta\,r_{5}^{-\varepsilon}.
\end{equation}
\end{proposition}

\begin{proof}
Put $\rho_{n}=\kappa^{n}r_{5}$, $\theta_{n}=\theta(z,\rho_{n})$,
$\lambda_{n}=\lambda(z,\rho_{n})$, $T_{n}=\theta_{n}\kappa^{-n\varepsilon}$ and
$L_{n}=\lambda_{n}\kappa^{-n\varepsilon}$. By Lemma~\ref{lem:monotonicity},
$\lambda_{n}\le\kappa^{n\sigma}\lambda_{0}$, and $\sigma>1>\varepsilon$ gives
$L_{n}\le\kappa^{n(\sigma-\varepsilon)}\lambda_{0}\le\Lambda_{0}$ for all $n$.

We prove $T_{n}\le\eta$ by induction; $T_{0}=\theta(z,r_{5})\le\eta$ by
hypothesis. Assume $T_{n}\le\eta$; then $\theta_{n}\le\eta\le1$, so
\eqref{eq:theta-decay-2} applies with $\rho=\rho_{n}$. Dividing it by
$\kappa^{(n+1)\varepsilon}$,
\[
  T_{n+1}\le C_{27}\kappa^{2/3-\varepsilon}T_{n}
   +2C_{27}\kappa^{-5-\varepsilon}\theta_{n}^{1/2}T_{n}
   +C_{28}\kappa^{-1/2-\varepsilon}T_{n}^{1/2}L_{n}^{1/2}
   +C_{28}\kappa^{-3-\varepsilon}L_{n}.
\]
By \eqref{eq:kappa-props} the first term is $\le\tfrac18T_{n}$ and, since
$\theta_{n}^{1/2}\le\eta^{1/2}$, the second is $\le\tfrac18T_{n}$. For the third,
Young's inequality $ab\le\tfrac18a^{2}+2b^{2}$ with $a=T_{n}^{1/2}$ and
$b=C_{28}\kappa^{-1/2-\varepsilon}L_{n}^{1/2}$ gives
$\le\tfrac18T_{n}+2C_{28}^{2}\kappa^{-1-2\varepsilon}L_{n}$. Hence
\[
  T_{n+1}\le\tfrac38T_{n}+C_{29}L_{n}
   \le\tfrac38\eta+C_{29}\Lambda_{0}=\tfrac38\eta+\tfrac12\eta=\tfrac78\eta\le\eta,
\]
completing the induction.

For \eqref{eq:iteration-concl}, let $0<r\le r_{5}$ and choose $n\in\N_{0}$ with
$\kappa^{n+1}r_{5}<r\le\kappa^{n}r_{5}$. By Lemma~\ref{lem:monotonicity} with
$\rho_{1}=r$, $\rho_{2}=\rho_{n}$ and $\rho_{n}/r<\kappa^{-1}$,
\[
  \alpha(z,r)\le\kappa^{-1/2}\alpha(z,\rho_{n}),\quad
  \beta(z,r)\le\kappa^{-1/2}\beta(z,\rho_{n}),\quad
  \delta(z,r)^{2}\le\kappa^{-4/3}\delta(z,\rho_{n})^{2},
\]
so $\theta(z,r)\le\kappa^{-4/3}\theta(z,\rho_{n})\le\kappa^{-4/3}\eta\kappa^{n\varepsilon}$.
Finally $\kappa^{n}<r/(\kappa r_{5})$ gives
$\kappa^{n\varepsilon}<\kappa^{-\varepsilon}r_{5}^{-\varepsilon}r^{\varepsilon}$,
whence $\theta(z,r)\le\kappa^{-4/3-\varepsilon}\eta r_{5}^{-\varepsilon}r^{\varepsilon}$.
\end{proof}

\begin{lemma}[Upper semicontinuity of $\alpha$ in the base point]
\label{lem:alpha-usc}
Let $0<r<R$ and $t_{1}<t_{2}$ be such that
$\overline{\Ball{R}{x_{0}}}\times[t_{1},t_{2}]\subset\mathcal{O}$ and
$t_{2}\in I$. Then
\begin{equation}\label{eq:alpha-usc}
  \lim_{h\to0^{+}}\ \operatorname*{ess\,sup}_{s\in(t_{1},t_{2}+h)}
     \int_{\Ball{r}{x_{0}}}\abs{u(y,s)}^{2}dy
  \;\le\;\operatorname*{ess\,sup}_{s\in(t_{1},t_{2})}
     \int_{\Ball{R}{x_{0}}}\abs{u(y,s)}^{2}dy .
\end{equation}
\end{lemma}

\begin{proof}
Write $S$ for the right side, which is finite by (S1). Choose
$\epsilon_{0}>0$ with $[t_{1},t_{2}+\epsilon_{0}]\subset I$ and
$\overline{\Ball{R}{x_{0}}}\times[t_{1},t_{2}+\epsilon_{0}]\subset\mathcal{O}$;
this is possible because $\mathcal{O}$ is open and the left-hand set is compact.
Let $\psi_{1}\in C_{c}^{\infty}(\Ball{R}{x_{0}})$ with $0\le\psi_{1}\le1$ and
$\psi_{1}\equiv1$ on $\Ball{r}{x_{0}}$. Define
\[
  \Xi(a,b)=\int_{a}^{b}\!\!\int_{\Ball{R}{x_{0}}}
   \Bigl[\abs{u}^{2}\abs{\Delta\psi_{1}}
    +\bigl(\abs{u}^{2}+2\abs{p}\bigr)\abs{u}\abs{\nabla\psi_{1}}
    +2\abs{f}\abs{u}\Bigr]dy\,ds ,
\]
which is finite for $[a,b]\subset[t_{1},t_{2}+\epsilon_{0}]$ by the
integrability checks in Step 2 of the proof of Lemma~\ref{lem:lei-pointwise};
by absolute continuity of the integral, $\Xi(a,b)\to0$ as $b-a\to0$.

Fix $0<h<\min\{\epsilon_{0}/2,t_2-t_1\}$ and let $T\in(t_{2},t_{2}+h)$. Let
$T_{0}\in(t_{2}-h,t_{2})$ and $h'\in(0,t_{2}-T_{0})$, and let
$\chi=\chi_{T_{0},h'}$ be a smooth nondecreasing function with $\chi\equiv0$ on
$(-\infty,T_{0}]$, $\chi\equiv1$ on $[T_{0}+h',\infty)$ and
$0\le\chi'\le2/h'$ supported in $(T_{0},T_{0}+h')$; such $\chi$ is obtained by
mollifying the piecewise affine function as in \eqref{eq:time-cutoff}. Multiply
$\chi$ by a smooth time cut-off equal to $1$ on
$(-\infty,t_{2}+\epsilon_{0}/2]$ and vanishing near $t_{2}+\epsilon_{0}$, which
does not affect the computation below because we integrate only up to
$T<t_{2}+\epsilon_{0}/2$. Then
$\phi(y,s)=\psi_{1}(y)\chi(s)\in C_{c}^{\infty}(\mathcal{O})$ is nonnegative and
vanishes for $s$ near $t_{1}$, so Lemma~\ref{lem:lei-pointwise} applies for almost every
$T$ (for each fixed cutoff). Since $\chi(T)=1$, $\partial_{s}\phi=\psi_{1}\chi'$,
$\Delta\phi=\chi\Delta\psi_{1}$, $\nabla\phi=\chi\nabla\psi_{1}$ and
$0\le\chi\le1$, discarding the (nonnegative) dissipation term gives
\[
  \int_{\Ball{r}{x_{0}}}\abs{u(y,T)}^{2}dy
  \;\le\;\int\abs{u(y,T)}^{2}\psi_{1}\,dy
  \;\le\;\int_{T_{0}}^{T_{0}+h'}\!\!\int_{\Ball{R}{x_{0}}}\abs{u}^{2}\chi'\,dy\,ds
   \;+\;\Xi(T_{0},T).
\]
The first term on the right is at most
$\frac{2}{h'}\int_{T_{0}}^{T_{0}+h'}\bigl(\int_{\Ball{R}{x_{0}}}\abs{u(\cdot,s)}^{2}\bigr)ds
 \le2S$ provided $(T_{0},T_{0}+h')\subset(t_{1},t_{2})$. To remove the factor
$2$, choose instead $T_{0}$ to be a Lebesgue point in $(t_{2}-h,t_{2})$ of the
function $s\mapsto\int_{\Ball{R}{x_{0}}}\abs{u(\cdot,s)}^{2}$ (almost every
point is one) and let $h'\to0^{+}$ through a fixed sequence, intersecting the countably
many full-measure sets of terminal times: the first term then converges to
$\int_{\Ball{R}{x_{0}}}\abs{u(\cdot,T_{0})}^{2}\le S$ if $T_{0}$ is moreover a
point where the essential supremum bound holds, which is the case for a.e.\
$T_{0}\in(t_{1},t_{2})$. Hence
\[
  \int_{\Ball{r}{x_{0}}}\abs{u(y,T)}^{2}dy\le S+\Xi(t_{2}-h,\,t_{2}+h)
  \qquad\text{for a.e. }T\in(t_{2},t_{2}+h).
\]
The same bound holds trivially for a.e.\ $T\in(t_{1},t_{2})$ since
$\Ball{r}{x_{0}}\subset\Ball{R}{x_{0}}$. Taking the essential supremum over
$T\in(t_{1},t_{2}+h)$ and letting $h\to0^{+}$, so that
$\Xi(t_{2}-h,t_{2}+h)\to0$, gives \eqref{eq:alpha-usc}.
\end{proof}

\begin{corollary}\label{cor:theta-usc}
Let $z_{0}\in\mathcal{O}$ and $0<\varrho<R$ with
$\overline{\Ball{R}{x_{0}}}\times[t_{0}-R^{2},t_{0}]\subset\mathcal{O}$ and
$t_{0}\in I$. Then
\begin{equation}\label{eq:theta-usc}
  \limsup_{z\to z_{0}}\ \alpha(z,\varrho)^{2}\;\le\;\frac R\varrho\,\alpha(z_{0},R)^{2},
  \qquad
  \lim_{z\to z_{0}}\beta(z,\varrho)=\beta(z_{0},\varrho),
  \qquad
  \lim_{z\to z_{0}}\delta(z,\varrho)=\delta(z_{0},\varrho).
\end{equation}
\end{corollary}

\begin{proof}
Let $z=(x,t)\to z_{0}$. For $z$ close enough to $z_{0}$ we have
$\Ball{\varrho}{x}\subset\Ball{R}{x_{0}}$ and
$(t-\varrho^{2},t)\subset(t_{0}-R^{2},t_{0}+h)$ for any prescribed $h>0$. Hence
\[
  \alpha(z,\varrho)^{2}
  =\frac1\varrho\operatorname*{ess\,sup}_{(t-\varrho^{2},t)}\int_{\Ball{\varrho}{x}}\abs{u}^{2}
  \le\frac1\varrho\operatorname*{ess\,sup}_{(t_{0}-R^{2},t_{0}+h)}\int_{\Ball{R}{x_{0}}}\abs{u}^{2},
\]
Since $\varrho<R$ strictly, fix an intermediate radius $r$ with
$\varrho<r<R$. For $z$ with $\abs{x-x_{0}}<r-\varrho$ we have
$\Ball{\varrho}{x}\subset\Ball{r}{x_{0}}$, so
\[
  \alpha(z,\varrho)^{2}
  \le\frac1\varrho\operatorname*{ess\,sup}_{s\in(t_{0}-R^{2},t_{0}+h)}
     \int_{\Ball{r}{x_{0}}}\abs{u(y,s)}^{2}dy .
\]
Letting $z\to z_{0}$ and then $h\to0^{+}$ and applying
Lemma~\ref{lem:alpha-usc} with $t_{1}=t_{0}-R^{2}$, $t_{2}=t_{0}$, inner radius
$r$ and outer radius $R$ (so that $r<R$ as the lemma requires) gives
$\limsup_{z\to z_{0}}\alpha(z,\varrho)^{2}
 \le\varrho^{-1}\operatorname*{ess\,sup}_{(t_{0}-R^{2},t_{0})}\int_{\Ball{R}{x_{0}}}\abs{u}^{2}
 =(R/\varrho)\alpha(z_{0},R)^{2}$.
For $\beta$ and $\delta$, the maps
$z\mapsto\iint_{\Cyl{\varrho}{z}}\abs{\nabla u}^{2}$ and
$z\mapsto\iint_{\Cyl{\varrho}{z}}\abs{p}^{3/2}$ are continuous, because
$\abs{\Cyl{\varrho}{z}\,\triangle\,\Cyl{\varrho}{z_{0}}}\to0$ as $z\to z_{0}$ and
$\abs{\nabla u}^{2},\abs{p}^{3/2}\in L^{1}$ near $z_{0}$, so absolute continuity
of the integral applies.
\end{proof}

\begin{proposition}[Step 2: Morrey decay on a neighbourhood]\label{prop:morrey-decay}
Let $q>5/2$, $\varepsilon=2/5$, and let $\kappa,\eta,\varepsilon_{*},\Lambda_{0}$
be as in Convention~\ref{conv:kappa}. Let $(u,p)$ be a suitable weak solution on
$\mathcal{O}$ with force $f\in L^{q}_{loc}$, and let $z_{0}\in\mathcal{O}$
satisfy
\begin{equation}\label{eq:morrey-decay-hyp}
  \limsup_{r\to0^{+}}\beta(z_{0},r)\;<\;\varepsilon_{*}.
\end{equation}
Then there are $r_{2}>0$ with
$\mathfrak{B}_{2r_{2}}(z_{0})\subset\mathcal{O}$ and $M\ge1$ such that
\begin{equation}\label{eq:morrey-decay}
  \max\bigl\{\alpha(z,r),\beta(z,r),\delta(z,r)^{2}\bigr\}\le M r^{\varepsilon}
  \qquad\text{for all }z\in\mathfrak{B}_{r_{2}}(z_{0})\text{ and }0<r<r_{2}.
\end{equation}
\end{proposition}

\begin{proof}
\emph{Step (a): descent at $z_{0}$.} Choose $r_{4}>0$ with
$\overline{\mathfrak{B}_{4r_{4}}(z_{0})}\subset\mathcal{O}$ and
$\beta(z_{0},\rho)\le\varepsilon_{*}$ for all $\rho\le r_{4}$; this is possible
by \eqref{eq:morrey-decay-hyp}. Put $M_{f}=\norm{f}_{L^{q}(\mathfrak{B}_{4r_{4}}(z_{0}))}$,
so that $\lambda(z,\rho)\le M_{f}\rho^{\sigma}$ whenever
$\Cyl{\rho}{z}\subset\mathfrak{B}_{4r_{4}}(z_{0})$. Put $T_{n}=\theta(z_{0},\kappa^{n}r_{4})\kappa^{-n\varepsilon}$ and
$L_{n}=\lambda(z_{0},\kappa^{n}r_{4})\kappa^{-n\varepsilon}$, so that
$L_{n}\le M_{f}r_{4}^{\sigma}\kappa^{n(\sigma-\varepsilon)}$. Apply
\eqref{eq:theta-decay-1} (the hypothesis
$\theta(z_{0},\cdot)\le\eta$ of \eqref{eq:theta-decay-2} is not yet known) at $z_{0}$ with
$\rho=\kappa^{n}r_{4}$ and divide by $\kappa^{(n+1)\varepsilon}$:
\[\begin{aligned}
  &T_{n+1}\le C_{27}\kappa^{2/3-\varepsilon}T_{n}
   \\ &\quad+C_{27}\kappa^{-5-\varepsilon}
     \bigl(\beta(z_{0},\kappa^{n}r_{4})^{1/2}+\beta(z_{0},\kappa^{n}r_{4})\bigr)T_{n}
   \\ &\quad+C_{28}\kappa^{-1/2-\varepsilon}T_{n}^{1/2}L_{n}^{1/2}
   +C_{28}\kappa^{-3-\varepsilon}L_{n}.
\end{aligned}\]
Since $\kappa^{n}r_{4}\le r_{4}$ we have $\beta(z_{0},\kappa^{n}r_{4})\le\varepsilon_{*}$
for every $n$, so the second term is at most
\[
C_{27}\kappa^{-5-\varepsilon}(\varepsilon_{*}^{1/2}+\varepsilon_{*})T_{n}
 \le\tfrac1{16}T_{n}
\]
 by the third inequality of \eqref{eq:kappa-props}; the
first is $\le\tfrac18T_{n}$ by the first inequality there; and the third is
$\le\tfrac18T_{n}+2C_{28}^{2}\kappa^{-1-2\varepsilon}L_{n}$ by Young's
inequality $ab\le\tfrac18a^{2}+2b^{2}$. Hence
$T_{n+1}\le\tfrac5{16}T_{n}+C_{29}L_{n}\le\tfrac38T_{n}+C_{29}L_{n}$, and
\[
  T_{n}\le\Bigl(\tfrac38\Bigr)^{n}T_{0}
   +C_{29}M_{f}r_{4}^{\sigma}\sum_{k=0}^{n-1}\Bigl(\tfrac38\Bigr)^{n-1-k}
     \kappa^{k(\sigma-\varepsilon)}
  \le\Bigl(\tfrac38\Bigr)^{n}T_{0}+C\,M_{f}r_{4}^{\sigma},
  \qquad C:=\tfrac85\,C_{29},
\]
because $\sum_{k=0}^{n-1}(3/8)^{n-1-k}\kappa^{k(\sigma-\varepsilon)}
 \le\sum_{m\ge0}(3/8)^{m}=\tfrac85$ (each factor $\kappa^{k(\sigma-\varepsilon)}\le1$). Shrinking
$r_{4}$ so that $C M_{f}r_{4}^{\sigma}\le\eta/8$ and choosing $n_{0}$ with
$(3/8)^{n_{0}}T_{0}\le\eta/8$ (possible since $T_{0}=\theta(z_{0},r_{4})<\infty$
by (S1)), we obtain
\begin{equation}\label{eq:descent}
  \theta(z_{0},R_{0})\le\tfrac14\eta\,\kappa^{n_{0}\varepsilon}\le\tfrac14\eta,
  \qquad R_{0}:=\kappa^{n_{0}}r_{4}.
\end{equation}

\emph{Step (b): transfer to nearby centres.} Let $\varrho\in(0,R_{0})$ with
$(R_{0}/\varrho)^{4/3+\varepsilon}\le2$. By Lemma~\ref{lem:monotonicity},
$\beta(z_{0},\varrho)\le(R_{0}/\varrho)^{1/2}\beta(z_{0},R_{0})$ and
$\delta(z_{0},\varrho)^{2}\le(R_{0}/\varrho)^{4/3}\delta(z_{0},R_{0})^{2}$, and
by Corollary~\ref{cor:theta-usc} applied with $R=R_{0}$,
\[
  \limsup_{z\to z_{0}}\theta(z,\varrho)
  \le\Bigl(\frac{R_{0}}{\varrho}\Bigr)^{4/3}\theta(z_{0},R_{0})
  \le2\cdot\tfrac14\eta=\tfrac12\eta .
\]
Hence there is $r_{2}\in(0,\varrho/4)$ with
$\mathfrak{B}_{2r_{2}}(z_{0})\subset\mathcal{O}$ and
\begin{equation}\label{eq:transfer}
  \theta(z,\varrho)\le\eta
  \qquad\text{for every }z\in\mathfrak{B}_{r_{2}}(z_{0}).
\end{equation}
Shrinking $r_{4}$ further at the start if necessary, we may also assume
$\lambda(z,\varrho)\le M_{f}\varrho^{\sigma}\le\Lambda_{0}$ for all such $z$.

\emph{Step (c): iterate at every centre.} For each
$z\in\mathfrak{B}_{r_{2}}(z_{0})$, Proposition~\ref{prop:iteration} with
$r_{5}=\varrho$ gives $\theta(z,r)\le M r^{\varepsilon}$ for $0<r\le\varrho$,
with coefficient $\kappa^{-4/3-\varepsilon}\eta\varrho^{-\varepsilon}$
independent of $z$. Set
$M=\max\{1,\kappa^{-4/3-\varepsilon}\eta\varrho^{-\varepsilon}\}$;
the estimate remains valid and $M\ge1$ as asserted.
Since $r_{2}<\varrho$, this is \eqref{eq:morrey-decay}.
\end{proof}

\begin{lemma}[Step 2 in Morrey form]\label{lem:step2-morrey-balls}
Let $z_{0},r_{2},M$ be as in Proposition~\ref{prop:morrey-decay} and put
$Q_{2}=\mathfrak{B}_{r_{2}/4}(z_{0})$. Then, with $\tau_{2},\tau_{3},\tau_{p}$
as in \eqref{eq:standing},
\begin{equation}\label{eq:step2-morrey}
  \mathbf{1}_{Q_{2}}u\in\mathcal{M}^{3,\tau_{2}},
  \qquad
  \mathbf{1}_{Q_{2}}\nabla u\in\mathcal{M}^{2,\tau_{3}},
  \qquad
  \mathbf{1}_{Q_{2}}p\in\mathcal{M}^{3/2,\tau_{p}},
\end{equation}
with norms bounded by $C(M,r_{2})$. Explicitly $\tau_{2}=25/3>5$ and
$\tau_{3}=\tau_{p}=25/8>5/2$, and $\tfrac1{\tau_{3}}=\tfrac1{\tau_{2}}+\tfrac15$.
\end{lemma}

\begin{proof}
By Lemma~\ref{lem:morrey-cylinders} it suffices to bound the integrals over
one-sided cylinders $\Cyl{\varrho}{z}$, $z\in\R^{3}\times\R$, $\varrho>0$. If
$\Cyl{\varrho}{z}\cap Q_{2}=\emptyset$ there is nothing to prove.

\emph{Small radii, $\varrho\le r_{2}/8$.} Pick
$z'=(x',t')\in\Cyl{\varrho}{z}\cap Q_{2}$, so that
$\Cyl{\varrho}{z}\subset\Cyl{2\varrho}{(x',t'+\varrho^{2})}$ by the triangle
inequality for $\dpar$ and Lemma~\ref{lem:parabolic-metric}. The centre
$z''=(x',t'+\varrho^{2})$ satisfies
$\dpar(z'',z_{0})\le\dpar(z',z_{0})+\varrho<r_{2}/4+r_{2}/8<r_{2}$, so
$z''\in\mathfrak{B}_{r_{2}}(z_{0})$, and $\varrho'=2\varrho<r_{2}$. By Proposition~\ref{prop:morrey-decay} and
Corollary~\ref{cor:gagliardo},
\[
  \iint_{\Cyl{\varrho}{z}}\abs{u}^{3}
  \le\iint_{\Cyl{\varrho'}{z''}}\abs{u}^{3}
  =\varrho'^{2}\gamma(z'',\varrho')^{3}
  \le\varrho'^{2}\bigl(2C_{9}M\varrho'^{\varepsilon}\bigr)^{3}
  =C M^{3}\varrho'^{2+3\varepsilon},
\]
and $2+3\varepsilon=16/5=5(1-3/\tau_{2})$ with $\tau_{2}=25/3$, which is the
Morrey condition for $\mathcal{M}^{3,\tau_{2}}$. Similarly
\[
  \iint_{\Cyl{\varrho}{z}}\abs{\nabla u}^{2}\le\varrho'\beta(z'',\varrho')^{2}
  \le M^{2}\varrho'^{1+2\varepsilon},
  \qquad
  \iint_{\Cyl{\varrho}{z}}\abs{p}^{3/2}\le\varrho'^{2}\delta(z'',\varrho')^{3}
  \le M^{3/2}\varrho'^{2+\frac32\varepsilon},
\]
and $1+2\varepsilon=9/5=5(1-2/\tau_{3})$ with $\tau_{3}=25/8$, while
$2+\tfrac32\varepsilon=13/5=5(1-\tfrac{3/2}{\tau_{p}})$ with $\tau_{p}=25/8$.
Finally $\tfrac1{\tau_{3}}=\tfrac8{25}=\tfrac3{25}+\tfrac5{25}
 =\tfrac1{\tau_{2}}+\tfrac15$.

\emph{Large radii, $\varrho>r_{2}/8$.} Cover the bounded set $Q_{2}$ by finitely
many cylinders of radius $r_{2}/8$ and apply the small-radius case to each: this
gives $M_{0}:=\iint_{Q_{2}}\abs{u}^{3}<\infty$, and likewise
$M_{0}':=\iint_{Q_{2}}\abs{\nabla u}^{2}<\infty$ and
$M_{0}'':=\iint_{Q_{2}}\abs{p}^{3/2}<\infty$ (all three are finite already by
(S1)). For $\varrho>r_{2}/8$,
\[
  \iint_{\Cyl{\varrho}{z}}\abs{\mathbf{1}_{Q_{2}}u}^{3}\le M_{0}
  =M_{0}\Bigl(\frac{r_{2}}8\Bigr)^{-16/5}\Bigl(\frac{r_{2}}8\Bigr)^{16/5}
  \le M_{0}\Bigl(\frac{r_{2}}8\Bigr)^{-16/5}\varrho^{16/5},
\]
which is the required Morrey bound with the constant enlarged by
$M_{0}(r_{2}/8)^{-16/5}$; identically for $\nabla u$ with the exponent $9/5$ and
for $p$ with the exponent $13/5$.
\end{proof}

\begin{remark}[Comparison with \protect{\cite[Lemma 13.4]{LemarieRieusset2016}}]
\label{rem:LR-pressure}
Lemarié-Rieusset states the conclusion of Step 2 as
$\mathbf{1}_{Q_{2}}u\in\mathcal{M}^{3,\tau_{2}}$ and
$\mathbf{1}_{Q_{2}}p\in\mathcal{M}^{q_{0},\tau_{2}/2}$ under the restriction
$1<\tau_{2}/5<\min\{q_{0},2\}$, equivalently
$5<\tau_{2}<5\min\{q_{0},2\}$. This is the condition printed in his
Lemma 13.4. In particular $\tau_2<5q_0$, so both terms in
$\lambda=\min\{10/\tau_{2},10q_{0}/\tau_{2}-2\}$ are positive,
as required by his iteration on p.~437. With $q_{0}=3/2$
that restriction forces $\tau_{2}<15/2$, i.e.\ $\varepsilon<1/3$. We instead
iterate Kukavica's functional $\theta=\alpha+\beta+\kappa^{-4}\delta^{2}$, in
which the pressure is normalized by $\delta^{2}$ rather than by $P_{r}$; this
allows $\varepsilon=2/5$ and hence $\tau_{2}=25/3$ for the velocity, at the
price of the weaker pressure exponent $\tau_{p}=25/8$ instead of
$\tau_{2}/2=25/6$. The weaker exponent is sufficient: the pressure Morrey norm
is used below only through fixed-scale $L^{3/2}$ integrability in
Lemma~\ref{lem:pressure-gradient-morrey}, never through its Morrey exponent. The gradient bound
$\mathbf{1}_{Q_{2}}\nabla u\in\mathcal{M}^{2,\tau_{3}}$,
$\tfrac1{\tau_{3}}=\tfrac1{\tau_{2}}+\tfrac15$, is the one extracted in
\cite[p.~439]{LemarieRieusset2016} from the same iteration. We use this
gradient information explicitly in the pressure-gradient estimate and the
velocity bootstrap below.
\end{remark}

\subsection{Step 3: the local representation and the Morrey bootstrap}\label{ss:step3}

Throughout \S\ref{ss:step3}--\S\ref{ss:step4}, $z_{0},r_{2},M$ are as in
Proposition~\ref{prop:morrey-decay}, $Q_{2}=\mathfrak{B}_{r_{2}/4}(z_{0})$, and
$\varphi\in C_{c}^{\infty}(\R^{3}\times\R)$ satisfies
\begin{equation}\label{eq:phi-zeta}
  \varphi\equiv1\text{ on }Q_{3}:=\mathfrak{B}_{r_{3}}(z_{0}),
  \qquad
  \supp\varphi\subset Q_{2}.
\end{equation}
with $0<r_{3}<r_{2}/4$. We write $N(y)=-\tfrac1{4\pi\abs{y}}$ as in
Proposition~\ref{prop:pressure-decomposition}, $W_{+}$ for the forward heat
kernel of \eqref{eq:heat-potential} (unit viscosity, as throughout), and
$I_{1},I_{2}$ for the parabolic Riesz potentials of
Definition~\ref{def:riesz-potential}.

We keep the pressure gradient in the undifferentiated source of the
localized heat equation. The pressure estimate below uses a fixed spatial
localization and one measurable weak gradient; it does not require a
commutator expansion.

\begin{lemma}[The localized equation]\label{lem:local-equation}
Set $v=\varphi u$. Then, in $\mathcal{D}'(\R^{3}\times\R)$,
\begin{equation}\label{eq:local-equation}
  \partial_{t}v-\Delta v
  \;=\;\mathbf{g}+\sum_{i=1}^{3}\partial_{i}\mathbf{h}_{i}
   \;-\;\varphi\nabla p,
  \qquad
  \mathbf{g}=(\partial_{t}\varphi)u+(\Delta\varphi)u-\varphi\,u\cdot\nabla u+\varphi f,
  \qquad
  \mathbf{h}_{i}=-2(\partial_{i}\varphi)u ,
\end{equation}
and the compact-support adjoint identity proved below gives
\begin{equation}\label{eq:duhamel}
  v=W_{+}*\Bigl(\mathbf{g}+\sum_{i}\partial_{i}\mathbf{h}_{i}-\varphi\nabla p\Bigr).
\end{equation}
\end{lemma}

\begin{proof}
By (S3), $\partial_{t}u-\Delta u=-u\cdot\nabla u-\nabla p+f$ in
$\mathcal{D}'(\mathcal{O})$, where $u\cdot\nabla u=\dv(u\otimes u)$ makes sense
as a distribution because $u\otimes u\in L^{5/3}_{loc}$. Multiplying by
$\varphi$ and using the Leibniz rules
$\partial_{t}(\varphi u)=(\partial_{t}\varphi)u+\varphi\partial_{t}u$ and
$\Delta(\varphi u)=(\Delta\varphi)u+2\nabla\varphi\cdot\nabla u+\varphi\Delta u$,
together with
$2\nabla\varphi\cdot\nabla u=\sum_{i}\partial_{i}\bigl(2(\partial_{i}\varphi)u\bigr)
 -2(\Delta\varphi)u$, gives \eqref{eq:local-equation}. For \eqref{eq:duhamel}, write
$\mathcal F$ for the compactly supported distribution on the right-hand
side of \eqref{eq:local-equation}. Given a compactly supported smooth
vector test $\psi$, let
\[
 \Psi(y,s)=\int_s^\infty\int_{\R^3}
                 W_+(x-y,t-s)\psi(x,t)\,dx\,dt .
\]
Then $-\partial_s\Psi-\Delta_y\Psi=\psi$.
Choose a smooth compact cutoff equal to one on a neighbourhood of the
supports of $v$ and $\mathcal F$. Testing the localized equation with
this cutoff times $\Psi$ gives
\[
 \langle v,\psi\rangle
   =\langle\mathcal F,\Psi\rangle
   =\langle W_+*\mathcal F,\psi\rangle .
\]
The last equality is the defining adjoint convolution identity; for the
locally integrable source terms it also follows from Fubini. The source
support makes the cutoff immaterial in the middle pairing. Thus
\eqref{eq:duhamel} holds distributionally, and almost everywhere whenever
the source potential is represented by a locally integrable function.
This proof requires neither an algebra of arbitrary causal distributions
nor unrestricted uniqueness for the heat equation.
\end{proof}

\begin{lemma}[Localized pressure-gradient Morrey bound]\label{lem:pressure-gradient-morrey}
Let $\mathfrak B_{2R}(z)\subset\mathcal O$, $R>0$, and
$25/3\le\tau\le25$. If
$\mathbf1_{\mathfrak B_R(z)}u\in\mathcal M^{3,\tau}$ and
$\mathbf1_{\mathfrak B_R(z)}\nabla u\in\mathcal M^{2,25/8}$, then the
distributional spatial gradient of $p$ has a jointly measurable representative
$Dp$ on $\mathfrak B_{R/2}(z)$, and
\begin{equation}\label{eq:pressure-gradient-morrey}
 \mathbf1_{\mathfrak B_{R/2}(z)}Dp\in\mathcal M^{6/5,\varkappa},
 \qquad \varkappa=\min\left\{q,\left(\frac1\tau+\frac8{25}\right)^{-1}\right\}.
\end{equation}
In particular, for every smooth test $\psi$ compactly supported there,
$\iint p\,\partial_i\psi=-\iint D_i p\,\psi$. The bounds are finite
functions of the incoming Morrey norms, the local energy, pressure and force
norms on a fixed surrounding cylinder, $q,\tau$, and the localization radii.
No bound is required on a shrinking-cell norm of the unnormalized pressure.
\end{lemma}

\begin{proof}
Write $z=(x,t)$ and fix the cylinder
$Q=B_R(x)\times(t-3R^2/4,t+R^2/4]$, which lies in
$\mathfrak B_R(z)$ and whose closure lies in $\mathfrak B_{2R}(z)$.
Let $\eta$ be the spatial cutoff of Lemma~\ref{lem:cutoff}, and let
$c(s)=\spavg{u(\cdot,s)}{B_R(x)}$.
In the pressure decomposition of Proposition~\ref{prop:pressure-decomposition}
use the singly centred tensor $U_{ij}=-u_i(u_j-c_j)$. Its differentiated,
localized source is
\[
 V_i=\sum_j\left[\eta(u_j-c_j)\partial_j u_i
                    +(\partial_j\eta)u_i(u_j-c_j)\right].
\]
All sources in this proof are extended by zero outside the fixed time
window of $Q$. The weak product rule, incompressibility and the completed
second Riesz transforms give, on the inner cylinder,
\begin{equation}\label{eq:pressure-gradient-decomposition}
 D_i p=-\sum_j\partial_i\partial_j\Delta^{-1}(\mathbf1_Q V_j)
       +\sum_j\partial_i\partial_j\Delta^{-1}(\mathbf1_Q\eta f_j)+H_i,
 \qquad H=\nabla(p_{\rm har}+p_8).
\end{equation}
Here $p_{\rm har}=p_2+\cdots+p_6$ is the smooth harmonic remainder on
the inner ball; it agrees almost everywhere there with $p-p_1-p_7-p_8$.
The gradient in $H$ is taken of this smooth representative and the smooth
annular potential $p_8$, not of an arbitrary representative of $p$.
The operators are their $L^{6/5}$ extensions, not pointwise
differentiation of a possibly undefined singular integral. Spatial
mollification and almost-everywhere limits select jointly measurable
representatives; the slice weak-derivative identities, followed by Fubini,
give the asserted space--time pairing. Uniqueness of locally integrable
weak derivatives identifies these representatives on overlapping cylinders.

We give the estimates on the three terms. The fixed spatial mean $c$ is
essentially bounded in time by the local energy bound. Thus its restriction to $Q$
belongs to every Morrey class used below. H\"older's inequality gives
$u\nabla u\in\mathcal M^{6/5,(1/\tau+8/25)^{-1}}$.
The remaining products in $V$ have at least this regularity: for example,
$u(u-c)$ has integrability $3/2$ and Morrey exponent $\tau/2$ before
lowering the two exponents. The bounded cutoff derivatives do not change
membership. Since $q>5/2$, the force term also belongs to
$\mathcal M^{6/5,\varkappa}$. These facts use only the fixed support and
H\"older's inequality, with its volume factors.

For completeness, the spatial Riesz operators preserve this parabolic
Morrey bound for the spatially supported sources just used. Split each
source into $B_{2r}(y)$ and the spatial annuli of radii $2^k r$, on the
time interval of a testing cylinder. The near term is bounded by the
$L^{6/5}$ Calder\'on--Zygmund estimate. For the annuli use the kernel
bound $C|x-y|^{-3}$, H\"older in space and time, and the source Morrey
bound on the larger cylinders. After normalization, the series is bounded
by a constant times
\[
 \sum_{k\ge0}2^{k(5/3-5/\varkappa)}
 \le \sum_{k\ge0}2^{-2k/15}<\infty,
 \qquad \varkappa\le25/9.
\]
The constants here depend only on $\varkappa$ and the fixed-exponent
Calder\'on--Zygmund constant, not the number of annuli or the source.
The Newtonian kernel derivatives used in these estimates include, for
$y\ne0$,
\begin{equation}\label{eq:third-derivative-N}
 4\pi\,\partial_i\partial_j\partial_l N(y)
 =-3\delta_{jl}\frac{y_i}{|y|^5}
  -3\delta_{ij}\frac{y_l}{|y|^5}
  -3\delta_{il}\frac{y_j}{|y|^5}
  +15\frac{y_i y_j y_l}{|y|^7},
\end{equation}
obtained by differentiating $N(y)=-(4\pi|y|)^{-1}$.

For the remainder, harmonic interior estimates on the fixed ball and
the separated support of the density of $p_8$ give, for almost every $s$,
\[
 \sup_{B_{R/2}(x)}|H(\cdot,s)|\le M_R(s),\qquad
 \int_{t-3R^2/4}^{t+R^2/4} M_R(s)^{3/2}\,ds<\infty .
\]
More explicitly one may take, with absolute constants,
\[
 M_R(s)=C R^{-3}\left(\alpha((x,t+R^2/4),R)^2
              +\|p(\cdot,s)\|_{L^{3/2}(B_R(x))}\right)
           +C R^{-2}\|f(\cdot,s)\|_{L^{3/2}(B_R(x))}.
\]
The pressure term is integrated at this fixed scale; finiteness follows
from (S1) and $q>5/2$. H\"older in time shows that on any testing cylinder
of radius $r$ the $6/5$-power integral of the restricted remainder is at
most $C r^{17/5}\|M_R\|_{L^{3/2}}^{6/5}$.
For small $r$ this is stronger than the required power
$r^{5-6/\varkappa}$, since $\varkappa\le25/9$; for large $r$ use the finite
integral on its bounded support. Adding the three bounds proves
\eqref{eq:pressure-gradient-morrey}.
\end{proof}

We record three elementary facts about parabolic Morrey spaces.

\begin{lemma}[H\"older, embeddings, and Minkowski]\label{lem:morrey-minkowski}
Let all functions be measurable on $\R^{3}\times\R$.
\begin{enumerate}
\item[(i)] \emph{(H\"older.)} If $g_{1}\in\mathcal{M}^{P_{1},\tau_{1}}$,
  $g_{2}\in\mathcal{M}^{P_{2},\tau_{2}}$,
  $\tfrac1P=\tfrac1{P_{1}}+\tfrac1{P_{2}}$ and
  $\tfrac1\tau=\tfrac1{\tau_{1}}+\tfrac1{\tau_{2}}$, then
  \begin{equation}\label{eq:morrey-holder}
    \norm{g_{1}g_{2}}_{\mathcal{M}^{P,\tau}}
    \le\norm{g_{1}}_{\mathcal{M}^{P_{1},\tau_{1}}}
       \norm{g_{2}}_{\mathcal{M}^{P_{2},\tau_{2}}} .
  \end{equation}
\item[(ii)] \emph{(Lowering the integrability exponent.)} If
  $1\le P'\le P\le\tau$ then
\[
\norm{g}_{\mathcal{M}^{P',\tau}}\le(\tfrac{8\pi}{3})^{\frac1{P'}-\frac1P}
   \norm{g}_{\mathcal{M}^{P,\tau}}
,
\] with no support hypothesis.
\item[(iii)] \emph{(Lowering the Morrey exponent.)} If $P\le\varkappa\le\tau$ and
  $\supp g\subset E$ with $\dpar\text{-}\diam E\le D$, then
  \begin{equation}\label{eq:morrey-lower}
    \norm{g}_{\mathcal{M}^{P,\varkappa}}
    \le\max\{1,D^{5(\frac1\varkappa-\frac1\tau)}\}\,
       \norm{g}_{\mathcal{M}^{P,\tau}} .
  \end{equation}
\item[(iv)] \emph{(Minkowski.)} If $k\in L^{1}(\R^{3})$ and
  $g\in\mathcal{M}^{P,\tau}$ then
  $\bigl\|\int_{\R^{3}}k(y)\,g(\cdot-y,\cdot)\,dy\bigr\|_{\mathcal{M}^{P,\tau}}
   \le\norm{k}_{L^{1}}\norm{g}_{\mathcal{M}^{P,\tau}}$; the same holds for a
  convolution in all four space-time variables with $k\in L^{1}(\R^{3}\times\R)$.
\end{enumerate}
\end{lemma}

\begin{proof}
(i) H\"older's inequality on each ball $\mathfrak{B}_{r}(z)$ with exponents
$P_{1}/P$ and $P_{2}/P$, and
$5(\tfrac1P-\tfrac1\tau)=5(\tfrac1{P_{1}}-\tfrac1{\tau_{1}})
 +5(\tfrac1{P_{2}}-\tfrac1{\tau_{2}})$.

(ii) H\"older on $\mathfrak{B}_{r}(z)$ against
$\abs{\mathfrak{B}_{r}}^{\frac1{P'}-\frac1P}
 =(\tfrac{8\pi}{3})^{\frac1{P'}-\frac1P}r^{5(\frac1{P'}-\frac1P)}$, and
$5(\tfrac1{P'}-\tfrac1P)+5(\tfrac1P-\tfrac1\tau)=5(\tfrac1{P'}-\tfrac1\tau)$.

(iii) If $D=0$, then $E$ contains at most one point and $g=0$ almost
everywhere, so the assertion is immediate. Assume $D>0$.
Let $r>0$ and $z$. If $r\le D$ then
$r^{-5(\frac1P-\frac1\varkappa)}=r^{5(\frac1\varkappa-\frac1\tau)}
 r^{-5(\frac1P-\frac1\tau)}\le D^{5(\frac1\varkappa-\frac1\tau)}
 r^{-5(\frac1P-\frac1\tau)}$, giving the bound. If $r>D$ then
$E$ is contained in the closed ball of radius $D$ centred at any point
of $E$. Apply the Morrey bound on the open ball of radius $D+\delta$
and let $\delta\downarrow0$ (the empty-set case is immediate). This gives
$\bigl(\iint_{\mathfrak{B}_{r}(z)}\abs{g}^{P}\bigr)^{1/P}
 \le D^{5(\frac1P-\frac1\tau)}\norm{g}_{\mathcal{M}^{P,\tau}}
 \le r^{5(\frac1P-\frac1\varkappa)}D^{5(\frac1\varkappa-\frac1\tau)}
    \norm{g}_{\mathcal{M}^{P,\tau}}$, again as claimed.

(iv) Minkowski's integral inequality in $L^{P}(\mathfrak{B}_{r}(z))$ together
with the translation invariance of $\norm{\cdot}_{\mathcal{M}^{P,\tau}}$:
\[
\begin{aligned}
&\bigl(\iint_{\mathfrak{B}_{r}(z)}\abs{\int k(y)g(\cdot-y)dy}^{P}\bigr)^{1/P}
\\ &\le\int\abs{k(y)}\bigl(\iint_{\mathfrak{B}_{r}(z-(y,0))}\abs{g}^{P}\bigr)^{1/P}dy
\\ &\le\norm{k}_{L^{1}}r^{5(\frac1P-\frac1\tau)}\norm{g}_{\mathcal{M}^{P,\tau}}
\end{aligned}
.
\]
\end{proof}

\begin{proposition}[Step 3: one bootstrap round]\label{prop:bootstrap}
Assume $q>5/2$ and the conclusions \eqref{eq:step2-morrey} of Step 2.
Let $\tau>5$ be such that
$\mathbf{1}_{Q_{2}}u\in\mathcal{M}^{3,\tau}$ and
\begin{equation}\label{eq:bootstrap-cond}
  \frac1{\tau_{3}}+\frac1\tau>\frac25 .
\end{equation}
Since $\tau_3=25/8$, this last condition is equivalent, for $\tau>0$,
to $\tau<25/2$. It is a hypothesis, not a consequence of $\tau>5$ alone.
Then $\mathbf{1}_{Q_{3}}u\in\mathcal{M}^{3,\varsigma}$ where
\begin{equation}\label{eq:bootstrap-gain}
  \frac1\varsigma=\frac1\tau-\varpi,
  \qquad \varpi:=\frac15-\frac1{\tau_{2}}=\frac{\varepsilon}{5}=\frac2{25}.
\end{equation}
\end{proposition}

\begin{proof}
First suppose $25/3\le\tau<25/2$. We work on a small symmetric
parabolic ball $\mathfrak B_R(z)$ with
$\mathfrak B_{2R}(z)\subset Q_2$, and choose a smooth cutoff $\varphi$
equal to one on $\mathfrak B_{R/4}(z)$, supported in
$\mathfrak B_{R/2}(z)$. Its derivatives are bounded by constants depending
only on the fixed radius. Set
\[
 \frac1\kappa=\frac1\tau+\frac8{25},\qquad
 \frac1\theta=\frac1\tau+\frac3{25},\qquad
 \frac1\varsigma=\frac1\tau-\frac2{25}.
\]
Then $3/2<\kappa<5/2$, $3<\theta<5$ and $\varsigma>5$.
Since $\kappa<5/2<q$, Lemma~\ref{lem:pressure-gradient-morrey} gives
$\mathbf1_{\mathfrak B_{R/2}(z)}Dp\in\mathcal M^{6/5,\kappa}$.

\emph{The undifferentiated source.}
With the notation of Lemma~\ref{lem:local-equation}, put
\[
 F=(\partial_t\varphi+\Delta\varphi)u
           -\varphi u\cdot\nabla u+\varphi f-\varphi Dp,
 \qquad G_i=-2(\partial_i\varphi)u .
\]
By Morrey H\"older,
$\varphi u\cdot\nabla u\in\mathcal M^{6/5,\kappa}$ because
$5/6=1/3+1/2$ and $1/\kappa=1/\tau+1/\tau_3$.
The cutoff-linear terms have this same membership by lowering exponents
on their bounded support. The force term does too, since
$f\in L^q_{\rm loc}$, $6/5<\kappa<q$.
Together with the pressure-gradient estimate, this proves
$F\in\mathcal M^{6/5,\kappa}$.

\emph{The differentiated sources.}
The fields $G_i$ initially belong to $\mathcal M^{3,\tau}$.
Since $3<\theta<\tau$ and they have bounded support,
Lemma~\ref{lem:morrey-minkowski}(iii) yields
$G_i\in\mathcal M^{3,\theta}$.

\emph{The two potential estimates.}
The localized equation and its causal representation give
\[
 v=W_+*F+\sum_i(\partial_i W_+)*G_i,\qquad
 |v|\le C\left(I_2|F|+\sum_i I_1|G_i|\right)
 \quad\hbox{almost everywhere}.
\]
The second inequality uses only the heat kernel and its spatial
derivatives, with the maximum gauge of Definition~\ref{def:riesz-potential}.
Its constant is absolute and chosen before the sources. The same proof
may first use the sum gauge $d_\Sigma=|x-y|+|t-s|^{1/2}$: for $a=1,2$,
\[
 I_a^\Sigma g\le I_a g\le 2^{5-a}I_a^\Sigma g .
\]
These absolute comparison factors transfer both the pointwise bound and
the Morrey estimates to the maximum gauge; the two potentials are not
identified literally.
Corollary~\ref{cor:adams} supplies both the almost-everywhere finiteness
needed for these convolutions and the bounds
\[
 I_2|F|\in\mathcal M^{P_F,\varsigma},
 \quad P_F=\frac{6/5}{1-2\kappa/5},\qquad
 I_1|G_i|\in\mathcal M^{P_G,\varsigma},
 \quad P_G=\frac{3}{1-\theta/5},
\]
because
\[
 \frac1\varsigma=\frac1\kappa-\frac25
                =\frac1\theta-\frac15=\frac1\tau-\frac2{25}.
\]
Both $P_F$ and $P_G$ are at least $3$ and at most $\varsigma$.
Lower their integrability exponents to $3$ using
Lemma~\ref{lem:morrey-minkowski}(ii), and add the finitely many terms.
Since $v=u$ on $\mathfrak B_{R/4}(z)$, this proves the desired local
velocity bound. All constants are finite functions of the incoming
bounds and the cutoff radius; no uniform bound independent of the
solution's incoming norms is asserted by this proposition.

Choose $R$ from the positive gap $r_2/4-r_3$ and cover
$\overline{Q_3}$ by finitely many such inner balls with their doubled
balls in $Q_2$. Restriction and the triangle inequality for Morrey norms
give the conclusion on $Q_3$.

Finally, if $5<\tau<25/3$, use the Step 2 velocity exponent $25/3$
instead of $\tau$. The case just proved gives exponent $25$.
Here $(1/\tau-2/25)^{-1}<25$, so lowering the Morrey exponent on
the bounded set $Q_3$ gives exactly \eqref{eq:bootstrap-gain}.
\end{proof}

\begin{corollary}[The bootstrap terminates after one round]\label{cor:one-round}
With $\varepsilon=2/5$, Proposition~\ref{prop:bootstrap} applies once with
$\tau=\tau_{2}=25/3$ and yields $\mathbf{1}_{Q_{3}}u\in\mathcal{M}^{3,\varsigma}$
with $\varsigma=25$; and then
\begin{equation}\label{eq:one-round}
  \frac1\varsigma+\frac1{\tau_{2}}=\frac1{25}+\frac3{25}=\frac4{25}<\frac15 .
\end{equation}
A second round is not available, since
$\tfrac1{\tau_{3}}+\tfrac1\varsigma=\tfrac8{25}+\tfrac1{25}=\tfrac9{25}<\tfrac25$
violates \eqref{eq:bootstrap-cond}.
\end{corollary}

\begin{proof}
Condition \eqref{eq:bootstrap-cond} for $\tau=\tau_{2}$ reads
$\tfrac8{25}+\tfrac3{25}=\tfrac{11}{25}>\tfrac{10}{25}$, which holds. The gain
is $\tfrac1\varsigma=\tfrac3{25}-\tfrac2{25}=\tfrac1{25}$. The two displayed
computations are immediate.
\end{proof}

\subsection{Step 4: H\"older continuity}\label{ss:step4}

\begin{theorem}[Endgame]\label{thm:endgame}
Let $q>5/2$ and let $z_{0},r_{2},M$ be as in
Proposition~\ref{prop:morrey-decay} with $\varepsilon=2/5$. Then $u$ agrees
almost everywhere on $Q_{3}=\mathfrak{B}_{r_{3}}(z_{0})$, for every
$r_{3}<r_{2}/4$, with a function that is H\"older continuous with respect to
$\dpar$ of exponent
\begin{equation}\label{eq:gamma-value}
  \gamma_{0}\;=\;\min\Bigl\{\,2-\frac5q,\ \frac15\,\Bigr\}\;>\;0,
\end{equation}
with norm bounded in terms of $q$, $M$, $r_{2}$, $r_{3}$ and the norms
$\norm{u}_{L^{10/3}}$, $\norm{p}_{L^{3/2}}$, $\norm{f}_{L^{q}}$ on
$\mathfrak{B}_{r_{2}}(z_{0})$. In particular every point of $Q_{3}$ is a regular
point of $u$ in the sense of Definition~\ref{def:regular}.
\end{theorem}

\begin{proof}
Choose an intermediate symmetric ball $Q_2'$ with
$\overline{Q_3}\subset Q_2'\Subset Q_2$.
Corollary~\ref{cor:one-round}, applied with $Q_2'$ as its inner set,
gives $\mathbf1_{Q_2'}u\in\mathcal M^{3,25}$; Step 2 also gives
$\mathbf1_{Q_2'}u\in\mathcal M^{3,25/3}$.
Choose a new cutoff $\varphi\in C_c^\infty(Q_2')$ with
$\varphi=1$ near $\overline{Q_3}$. Thus every localized source below is supported where
the improved velocity exponent is available.
Apply Lemma~\ref{lem:pressure-gradient-morrey} on finitely many balls
whose doubled balls lie in $Q_2'$ and whose inner balls cover
$\supp\varphi$.
Weak-derivative uniqueness identifies the pressure gradients on overlaps.
Since $1/25+8/25=9/25$, the resulting pressure gradient belongs locally to
$\mathcal M^{6/5,\min\{q,25/9\}}$.
Use the same sources $F=\mathbf g-\varphi Dp$ and $G_i=\mathbf h_i$
as in the proof of Proposition~\ref{prop:bootstrap}. Morrey H\"older
and bounded-support lowering now give
\[
 F\in\mathcal M^{6/5,\theta_0},\qquad
 G_i\in\mathcal M^{3,\theta_1},\qquad
 \theta_0=\min\{q,25/9\},\quad \theta_1=25/3.
\]
The pressure is retained in $F$, not converted into a differentiated
quadratic source. The local equation therefore has the representation
$v=W_+*(F+\sum_i\partial_iG_i)$.
Lower the integrability of the $G_i$ to $6/5$, and, on the fixed bounded
supports, lower the Morrey exponents to the values required by
Proposition~\ref{prop:heat-morrey-hoelder}. The conditions are
\[
 \gamma\le 2-\frac5{\theta_0}
       =\min\{2-5/q,1/5\},\qquad
 \gamma\le 1-\frac5{\theta_1}=\frac25 .
\]
They hold for $\gamma=\gamma_0>0$ in \eqref{eq:gamma-value}.
The resulting single representative agrees almost everywhere with $u$
on $Q_3$. On overlapping local balls the continuous representatives
coincide, since they agree almost everywhere; a finite cover gives the
claimed H\"older representative and norm bound.

For the quantitative bound, fix the cutoff family from the radii before
the solution and retain the source-norm bounds in the preceding estimates.
The fixed-scale harmonic remainder is controlled by the local energy and
pressure bounds; the local energy inequalities bound these quantities
in terms of the Step 2 bound and the listed surrounding norms.
The Calder\'on--Zygmund constants occur only at fixed exponents, and the
finite cover and cutoff derivatives depend only on the radii.
Together with the heat-potential estimate this gives exactly the dependence
asserted in the theorem.
\end{proof}

\begin{lemma}[One-sided pressure-gradient estimate]\label{lem:pressure-gradient-past}
Let $(u,p)$ be a suitable weak solution on $\mathcal O$ with force $f$
and exponent $q>5/2$. Write $E_R=B_R(0)\times(-R^2,0]$. Suppose
$\overline{E_R}\subset\mathcal O$, $0<r<R$, and $25/3\le\tau\le25$.
If
\[
 \mathbf1_{E_R}u\in\mathcal M^{3,\tau},\qquad
 \mathbf1_{E_R}\nabla u\in\mathcal M^{2,25/8},
\]
then the weak spatial gradient of $p$ has a jointly measurable,
locally integrable representative $Dp$ near $\overline{E_r}$, and
\[
 \mathbf1_{E_r}Dp\in\mathcal M^{6/5,\varkappa},\qquad
 \varkappa=\min\left\{q,\left(\frac1\tau+\frac8{25}\right)^{-1}\right\}.
\]
Its Morrey norm is bounded in terms of the two displayed norms, $q,\tau$,
the radii, and the following \emph{past-time} quantities only:
\[
 A_R=\operatorname*{ess\,sup}_{-R^2<s<0}
           \int_{B_R}|u(y,s)|^2\,dy,\qquad
 \|p\|_{L^{3/2}(E_R)},\qquad \|f\|_{L^q(E_R)}.
\]
No Morrey estimate or quantitative norm at positive times is assumed.
\end{lemma}

\begin{proof}
Choose $r<a<b<R$ and a fixed spatial cutoff equal to one on $B_b$
and supported in $B_R$. At almost every time $s\in(-a^2,0)$, apply
the spatial pressure decomposition with
$c(s)=\spavg u{B_R}(s)$ and the singly centred tensor
$-u_i(u_j-c_j)$. The differentiated source is
\[
 V_i=\sum_j\bigl[\eta(u_j-c_j)\partial_j u_i
                 +(\partial_j\eta)u_i(u_j-c_j)\bigr].
\]
Set all sources to zero outside the time interval $(-a^2,0]$.
The spatial operators do not change this time support. On $E_r$, the
same slice weak-derivative identity used in
Lemma~\ref{lem:pressure-gradient-morrey} is
\[
 D_i p=-\sum_j\partial_i\partial_j\Delta^{-1}
              (\mathbf1_{(-a^2,0]}V_j)
       +\sum_j\partial_i\partial_j\Delta^{-1}
              (\mathbf1_{(-a^2,0]}\eta f_j)+H_i.
\]
All cutoffs in this identity are spatial; no time derivative of the
indicator is taken. The mean satisfies
$|c(s)|\le |B_R|^{-1/2}A_R^{1/2}$ for a.e.\ past time.
Morrey H\"older gives
$\mathbf1_{(-a^2,0]}\eta u\nabla u
 \in\mathcal M^{6/5,(1/\tau+8/25)^{-1}}$.
The terms containing $c$ or a cutoff derivative have at least this
membership after bounded-support lowering: $u(u-c)$ first has
integrability $3/2$ and Morrey exponent $\tau/2$.
The restricted force belongs to $\mathcal M^{6/5,\varkappa}$ by
H\"older and bounded-support lowering from $L^q$.

For clarity, the Riesz estimate remains valid across the truncation
time. On a testing cylinder of radius $\ell$, split a source spatially
into the ball of radius $2\ell$ and dyadic annuli of radii $L=2^k\ell$,
while keeping its time interval unchanged. At exponent $P=6/5$, the
near term follows from the spatial $L^P$ estimate integrated in time.
For an annulus the kernel bound and spatial H\"older give
\[
 \|Tg_k\|_{L^P(B_\ell\times J_\ell)}
 \le C\ell^{3/P}L^{-3/P}
                 \|g\|_{L^P(B_{CL}\times J_\ell)}
 \le C\ell^{3/P}L^{2/P-5/\varkappa}
                 \|g\|_{\mathcal M^{P,\varkappa}}.
\]
Here the last norm is estimated on a parabolic ball of radius $CL$
containing the spatial annulus and the original time interval.
After division by $\ell^{5/P-5/\varkappa}$, the annular factor is
$C2^{k(5/3-5/\varkappa)}$. Since $\varkappa\le25/9$, it is summable,
with exponent at most $-2/15$. The argument applies directly to the
zero-extended past sources; it never asks for their future values.

Harmonic interior estimates and separation from the annular force
source give a majorant on $B_r$ of the form
\[
 \sup_{B_r}|H(\cdot,s)|\le m(s),\qquad
 m(s)\le C_{r,R}\bigl(A_R+
       \|p(\cdot,s)\|_{L^{3/2}(B_R)}+
       \|f(\cdot,s)\|_{L^{3/2}(B_R)}\bigr).
\]
Thus $m\in L^{3/2}(-a^2,0)$ with a bound depending only on the stated
past quantities. For a testing cylinder of radius $\ell$, H\"older in
time bounds the $6/5$-power integral of the restricted remainder by
$C\ell^{17/5}\|m\|_{L^{3/2}}^{6/5}$. This is sufficient at small radii,
since $17/5>5-6/\varkappa$; bounded support handles large radii.
This proves the asserted Morrey estimate.

To obtain a representative on an open neighbourhood as well, use the
same spatial construction on a short time collar of $[-a^2,0]$ lying
inside the solution domain. Local suitability gives
$u\nabla u\in L^{5/4}_{\rm loc}$ there by
$u\in L^{10/3}_{\rm loc}$ and $\nabla u\in L^2_{\rm loc}$.
The spatial $L^{5/4}$ Riesz estimate and the harmonic remainder therefore
give a locally integrable gradient on that collar. Spatial mollification
and a.e.\ limits give joint measurability; the slice identities and
Fubini give the spacetime weak-derivative identity. Weak-derivative
uniqueness identifies it with the past representative already estimated.
The collar is used only for this qualitative identity, not for a
quantitative Morrey bound.
\end{proof}

\begin{proposition}[One-sided bootstrap and endgame]\label{prop:past-bootstrap}
Let $(u,p)$ be a suitable weak solution on $\mathcal O$ with force $f$.
Let $E_R=B_R(0)\times(-R^2,0]$ with
$\overline{E_R}\subset\mathcal O$, and assume the finite past bounds
\[\begin{gathered}
 \mathbf1_{E_R}u\in\mathcal M^{3,25/3},\qquad
 \mathbf1_{E_R}\nabla u\in\mathcal M^{2,25/8},\\
 A_R<\infty,\quad p\in L^{3/2}(E_R),\quad f\in L^q(E_R),\quad q>5/2.
\end{gathered}\]
For every $0<r<R$ one has
$\mathbf1_{E_r}u\in\mathcal M^{3,25}$.
For every $0<\rho<r$, there is a globally defined parabolically
H\"older function $\bar v$ of exponent
$\gamma_0=\min\{2-5/q,1/5\}>0$, agreeing a.e.\ with $u$ on $E_\rho$.
Its H\"older seminorm and its supremum on $\overline{E_\rho}$ are
bounded in terms of the displayed past bounds, $q$, and $R,r,\rho$.
No conclusion of regularity on the top face in the open-neighbourhood
sense is asserted.
\end{proposition}

\begin{proof}
First choose $r<a<b<R$. By Lemma~\ref{lem:pressure-gradient-past},
$\mathbf1_{E_b}Dp\in\mathcal M^{6/5,25/11}$, because
$1/(25/3)+8/25=11/25$ and $25/11<q$.
Choose a smooth product cutoff $\varphi$ equal to one near
$\overline{E_r}$, whose past support lies in $E_a$ and whose full
support lies in the open solution domain and inside the neighbourhood
where $Dp$ was constructed. Explicitly, the spatial cutoff is supported
in $B_a$ and is one on $B_r$; the temporal cutoff vanishes before
$-a^2$ and is one on $[-r^2,0]$, remaining one for a short future interval
before vanishing inside the available time collar. All derivatives at
times $t\le0$ are bounded in terms of the fixed radii. The future cutoff
may depend on the collar, but its derivatives occur only at $t>0$.

Define the full localized sources
\[\begin{aligned}
 &F=(\partial_t\varphi+\Delta\varphi)u
       -\varphi u\cdot\nabla u+\varphi f-\varphi Dp,
 \qquad G_i=-2(\partial_i\varphi)u,\\
 &F^-=\mathbf1_{\{t\le0\}}F,\qquad
 G_i^-=\mathbf1_{\{t\le0\}}G_i.
\end{aligned}\]
The local equation and its adjoint-test proof give the full distributional
representation for $\varphi u$. Splitting the source in that
representation, rather than differentiating $\mathbf1_{\{t\le0\}}\varphi u$,
introduces no terminal-time distribution. The kernels $W_+$ and
$\partial_iW_+$ vanish at negative time differences. Hence
\[
 \varphi u=W_+*F^-+\sum_i(\partial_iW_+)*G_i^-
 \quad\hbox{a.e. on }\{t\le0\}.
\]
All future sources make zero contribution on this set; the single time
slice $t=0$ has spacetime measure zero.

Morrey H\"older and bounded-support lowering, applied only to the past
sources, give
\[
 F^-\in\mathcal M^{6/5,25/11},\qquad
 G_i^-\in\mathcal M^{3,25/6}.
\]
For the latter use $25/6<25/3$; for the former use the pressure estimate
above and $u\nabla u$ with reciprocal Morrey exponent $11/25$.
Adams' estimate with potential orders two and one gives respectively
\[
 I_2|F^-|\in\mathcal M^{66/5,25},\qquad
 I_1|G_i^-|\in\mathcal M^{18,25}.
\]
Lower the integrability exponents to three. The heat-kernel pointwise
bound and the a.e.\ identity on $E_r$ now prove
$\mathbf1_{E_r}u\in\mathcal M^{3,25}$, with constants depending only
on the past bounds and the chosen radii.

For the final step choose new intermediate radii
$\rho<a'<b'<r$ and a new cutoff equal to one near $\overline{E_\rho}$
whose past support lies in $E_{a'}$. Its sources are therefore supported
where the improved velocity bound is available. Apply
Lemma~\ref{lem:pressure-gradient-past} with outer radius $r$ and
inner radius $b'$ at $\tau=25$. The resulting past pressure-gradient
exponent is $\theta_0=\min\{q,25/9\}$. The same source calculation gives
\[
 F^-\in\mathcal M^{6/5,\theta_0},\qquad
 G_i^-\in\mathcal M^{3,25/3}.
\]
Here the nonlinear term uses the improved exponent 25, not 25/3;
all other past bounds on $E_r$ follow by restriction from $E_R$.
Lower the integrability of $G_i^-$ to $6/5$ and their Morrey exponents,
on their bounded supports, to those prescribed by
Proposition~\ref{prop:heat-morrey-hoelder} for $\gamma=\gamma_0$.
The requisite inequalities are
\[
 \gamma_0\le2-5/\theta_0=\min\{2-5/q,1/5\},\qquad
 \gamma_0\le1-5/(25/3)=2/5.
\]
That proposition gives a global H\"older representative of the
potential of these past sources. By causality it agrees a.e.\ with u
on $E_\rho$. Finally, for $z\in\overline{E_\rho}$ average its
H\"older inequality over $y\in E_\rho$ to obtain
\[
 |\bar v(z)|\le |E_\rho|^{-1/3}\|u\|_{L^3(E_\rho)}
       +[\bar v]_{C^{\gamma_0,\gamma_0/2}}
            \operatorname{diam}_{\dpar}(\overline{E_\rho})^{\gamma_0}.
\]
This controls the supremum using past data only and proves the proposition.
\end{proof}

\begin{remark}[The exponent delivered by this route]\label{rem:exponent}
With the fixed choice $\varepsilon=2/5$, the single bootstrap round gives
velocity exponent $25$, source exponents
$\theta_0=\min\{q,25/9\}$ and $\theta_1=25/3$, and hence
$\gamma_0=\min\{2-5/q,1/5\}$. No optimization over a variable
$\varepsilon$ is asserted here: the pressure-gradient estimates above
were used only in the stated range $25/3\le\tau\le25$.
The ceiling $2/3$ appearing in the blow-up route
\cite[Lemma 2.4]{LadyzhenskayaSeregin1999},
\cite[Remark 6.6]{Tsai2018} is not attained by this argument; only the
positivity of $\gamma_0$ is used in Theorems~\ref{thm:A}--\ref{thm:C}.
\end{remark}

\subsection{Proof of Theorems~\ref{thm:A} and \ref{thm:B}}\label{ss:assembly}

\begin{proof}[Proof of Theorem~\ref{thm:B}]
Let $\varepsilon_{1}=\varepsilon_{*}$, the absolute constant of
Convention~\ref{conv:kappa}; note that the hypothesis \eqref{eq:thmB-hyp} of
Theorem~\ref{thm:B} is stated for $\beta^{2}$ against $\varepsilon_{1}^{2}$,
whereas Proposition~\ref{prop:morrey-decay} is stated for $\beta$ against
$\varepsilon_{*}$, so the two thresholds are the same number and no squaring
discrepancy arises. If
$\limsup_{r\to0^{+}}\beta(z_{0},r)^{2}<\varepsilon_{1}^{2}$ then
$\limsup_{r\to0^{+}}\beta(z_{0},r)<\varepsilon_{*}$, so
Proposition~\ref{prop:morrey-decay} provides $r_{2}$ and $M$, and
Lemma~\ref{lem:step2-morrey-balls}, Proposition~\ref{prop:bootstrap},
Corollary~\ref{cor:one-round} and Theorem~\ref{thm:endgame} show that $u$
agrees a.e.\ on the open set $\mathfrak{B}_{r_{2}/8}(z_{0})$ with a function
that is H\"older continuous of exponent $\gamma_{0}$ with respect to $\dpar$.
Since $z_{0}\in\mathfrak{B}_{r_{2}/8}(z_{0})$, Definition~\ref{def:regular} is
satisfied and $z_{0}$ is a regular point.
\end{proof}

\begin{lemma}[Caccioppoli inequality, $\gamma$-form]\label{lem:caccioppoli-gamma}
Under the hypotheses of Lemma~\ref{lem:caccioppoli}, and with the same notation,
\begin{equation}\label{eq:caccioppoli-gamma}
  \alpha(r)+\beta(r)\;\le\;
    C_{25}'\Bigl[\kappa\,\gamma(\rho)+\kappa^{-1}\gamma(\rho)^{3/2}
      +\kappa^{-1}\delta(\rho)\gamma(\rho)^{1/2}\Bigr]
   +C_{26}'\kappa^{-1/2}\gamma(\rho)^{1/2}\lambda(\rho)^{1/2},
\end{equation}
with $C_{25}'$ absolute and $C_{26}'=C_{26}'(q)$.
\end{lemma}

\begin{proof}
Run the proof of Lemma~\ref{lem:caccioppoli} verbatim, changing only the bounds
of $I_{1}$ and $I_{2}$.

\emph{$I_{1}$.} By H\"older on $\Cyl{\rho}{z_{0}}$ with $1=\tfrac23+\tfrac13$,
$\iint_{\Cyl{\rho}{z_{0}}}\abs{u}^{2}
 \le\abs{\Cyl{\rho}{z_{0}}}^{1/3}\norm{u}^{2}_{L^{3}(\Cyl{\rho}{z_{0}})}
 =C\rho^{5/3}\rho^{4/3}\gamma(\rho)^{2}=C\rho^{3}\gamma(\rho)^{2}$, so
$I_{1}\le C r^{2}\rho^{-5}\cdot\rho^{3}\gamma(\rho)^{2}=C\kappa^{2}\gamma(\rho)^{2}$.

\emph{$I_{2}$.} Instead of the Poincar\'e step, use the spatial integral bound
\[\int_{B_{\rho}}\bigl|\abs{u}^{2}-\spavg{\abs{u}^{2}}{B_{\rho}}\bigr|\abs{u}\le2\int_{B_{\rho}}\abs{u}^{3},\]
valid at almost every time by H\"older's inequality on the mean term. Thus
\[
I_{2}\le Cr^{-2}\iint_{\Cyl{\rho}{z_{0}}}\abs{u}^{3}
 =Cr^{-2}\rho^{2}\gamma(\rho)^{3}=C\kappa^{-2}\gamma(\rho)^{3}
.
\]

The bounds on $I_{3}$ and $I_{4}$ are unchanged. Adding and taking square roots
as in Step 3 of that proof gives \eqref{eq:caccioppoli-gamma}.
\end{proof}

\begin{lemma}[Smallness at one scale propagates to all centres]\label{lem:thmA-start}
There is $\varepsilon_{0}=\varepsilon_{0}(q)>0$ such that: if $(u,p)$ is a
suitable weak solution on an open set containing $\overline{Q_{1}}$ with force
$f\in L^{q}(Q_{1})$ and \eqref{eq:thmA-hyp} holds, then, with
$r_{5}:=\kappa/4$,
\[
  \theta(z,r_{5})\le\eta
  \qquad\text{and}\qquad
  \lambda(z,r_{5})\le\Lambda_{0}
  \qquad\text{for every }z\in Q_{3/4}.
\]
\end{lemma}

\begin{proof}
Let $z=(x,t)\in Q_{3/4}$, so $\abs{x}<3/4$ and $t\in(-\tfrac9{16},0]$. Then
$\Cyl{1/4}{z}\subset Q_{1}$: indeed $\Ball{1/4}{x}\subset\Ball{1}{0}$ and
$(t-\tfrac1{16},t]\subset(-\tfrac9{16}-\tfrac1{16},0]\subset(-1,0]$. By
\eqref{eq:thmA-hyp},
\[
  \gamma(z,\tfrac14)^{3}=16\!\iint_{\Cyl{1/4}{z}}\abs{u}^{3}\le16\varepsilon_{0},
  \quad
  \delta(z,\tfrac14)^{3}\le16\varepsilon_{0},
  \quad
  \lambda(z,\tfrac14)=\bigl(\tfrac14\bigr)^{\sigma}\norm{f}_{L^{q}(\Cyl{1/4}{z})}
    \le\varepsilon_{0}^{1/q}.
\]
Apply Lemma~\ref{lem:caccioppoli-gamma} with $\rho=1/4$ and $r=r_{5}=\kappa/4$,
and Proposition~\ref{prop:lin34}(ii-b). For a.e.\ time $s$, write
$c(s)=\spavg u{B_{1/4}(x)}(s)$. The inequalities
$|a-b|^3\le4(|a|^3+|b|^3)$ and Jensen's inequality give
\[
 \int_{B_{1/4}(x)}|u(y,s)-c(s)|^3\,dy
 \le4\int_{B_{1/4}(x)}|u(y,s)|^3\,dy+4|B_{1/4}|\,|c(s)|^3
 \le8\int_{B_{1/4}(x)}|u(y,s)|^3\,dy.
\]
Integrating in time and multiplying by $(1/4)^{-2}$ proves
$\widehat C(z,\tfrac14)\le8\gamma(z,\tfrac14)^3$ for the spatial mean
used here. Taking cube roots in the pressure estimate gives
\[
  \delta(z,r_{5})\le C_{32}(q)^{1/3}\bigl[2\kappa^{-2/3}\gamma(z,\tfrac14)
    +\kappa^{1/3}\delta(z,\tfrac14)+\kappa^{1/2}\lambda(z,\tfrac14)^{1/2}\bigr].
\]
Every term on the right of both displays tends to $0$ as
$\varepsilon_{0}\to0$, with $\kappa$ fixed. Choose $\varepsilon_{0}(q)$ so small
that $\theta(z,r_{5})=\alpha(z,r_{5})+\beta(z,r_{5})+\kappa^{-4}\delta(z,r_{5})^{2}
\le\eta$; and $\lambda(z,r_{5})\le\lambda(z,\tfrac14)\le\varepsilon_{0}^{1/q}
\le\Lambda_{0}$ after shrinking $\varepsilon_{0}$ once more.
\end{proof}

\begin{proof}[Proof of Theorem~\ref{thm:A}]
Let $\varepsilon_{0}(q)$ be as in Lemma~\ref{lem:thmA-start} and
$r_{5}=\kappa/4$. The set of admissible centres is the \emph{one-sided}
cylinder $Q_{3/4}$, so the centre-shift device of
Lemma~\ref{lem:step2-morrey-balls}, which moves a centre forward in time by
$\varrho^{2}$, is not available as it stands; and $\overline{Q_{1/2}}$ contains
its top face, so no cut-off supported in $\{t\le0\}$ can equal $1$ there. Three
modifications repair this.

\emph{Step 1: the scale iteration at every centre of $Q_{3/4}$.} By
Lemma~\ref{lem:thmA-start}, $\theta(z,r_{5})\le\eta$ and
$\lambda(z,r_{5})\le\Lambda_{0}$ for every $z\in Q_{3/4}$, with $r_{5}$
\emph{fixed}. Proposition~\ref{prop:iteration} therefore gives
\begin{equation}\label{eq:thmA-morrey}
  \max\bigl\{\alpha(z,r),\beta(z,r),\delta(z,r)^{2}\bigr\}\le M\,r^{\varepsilon}
  \qquad\text{for all }z\in Q_{3/4},\ 0<r\le r_{5},
\end{equation}
with $M=\kappa^{-4/3-\varepsilon}\eta\,r_{5}^{-\varepsilon}$ an \emph{absolute}
constant.

\emph{Step 2: a one-sided Morrey transfer.} Put
$Q_{2}^{\sharp}=\Ball{5/8}{0}\times(-\tfrac{25}{64},0]$, a one-sided cylinder
with $\overline{Q_{1/2}}\cap\{t\le0\}\subset Q_{2}^{\sharp}\subset Q_{3/4}$. We
claim that \eqref{eq:step2-morrey} holds with $Q_{2}$ replaced by
$Q_{2}^{\sharp}$. By Lemma~\ref{lem:morrey-cylinders} it suffices to bound the
integrals over one-sided cylinders $\Cyl{\varrho}{w}$. Since
$Q_{2}^{\sharp}\subset\{t\le0\}$, only $\Cyl{\varrho}{w}\cap\{t\le0\}$
contributes. If that set meets $Q_{2}^{\sharp}$, pick
$w'\in\Cyl{\varrho}{w}\cap Q_{2}^{\sharp}$ and set
\[
  w''=\bigl(x_{w'},\ \min\{t_{w}+\varrho^{2},\,0\}\bigr),
\]
the \emph{truncated} centre shift. Then $w''\in Q_{3/4}$, because
$\abs{x_{w'}}<5/8<3/4$ and $t_{w''}\le0$ while
$t_{w''}\ge t_{w'}>-\tfrac{25}{64}>-\tfrac9{16}$. Moreover
$\Cyl{\varrho}{w}\cap\{t\le0\}\subset\Cyl{2\varrho}{w''}$: for the spatial part,
$\abs{x-x_{w''}}\le\abs{x-x_{w}}+\abs{x_{w}-x_{w'}}<2\varrho$; for the upper
time end, $s\le\min\{t_{w},0\}\le t_{w''}$; and for the lower end,
$s>t_{w}-\varrho^{2}\ge t_{w''}-3\varrho^{2}>t_{w''}-4\varrho^{2}$. The
computation of Lemma~\ref{lem:step2-morrey-balls} now applies with $w''$ in
place of $z''$ and \eqref{eq:thmA-morrey} in place of
\eqref{eq:morrey-decay}, for $\varrho\le r_{5}/2$; the large-radius case is
unchanged. This gives
$\mathbf{1}_{Q_{2}^{\sharp}}u\in\mathcal{M}^{3,\tau_{2}}$,
$\mathbf{1}_{Q_{2}^{\sharp}}\nabla u\in\mathcal{M}^{2,\tau_{3}}$,
$\mathbf{1}_{Q_{2}^{\sharp}}p\in\mathcal{M}^{3/2,\tau_{p}}$, with norms bounded
in terms of $M$ and $r_{5}$, hence by absolute constants.

\emph{Step 3: the one-sided bootstrap and endgame.}
We apply Proposition~\ref{prop:past-bootstrap} with
$R=5/8$, $r=9/16$ and $\rho=1/2$.
The velocity and gradient Morrey hypotheses on $E_R=Q_2^\sharp$
are precisely those proved in Step 2. Its pressure Morrey bound also
controls $\|p\|_{L^{3/2}(E_R)}$, and
$\|f\|_{L^q(E_R)}\le\varepsilon_0^{1/q}$ by the smallness assumption.
For the remaining past energy bound, choose a fixed
$0<\ell<\min\{r_5,3/4-R\}/4$ and cover $B_R$ by finitely many
balls $B_\ell(x_j)$ with $x_j\in B_R$. Cover $(-R^2,0)$ by finitely
many time intervals $(t_k-\ell^2,t_k)$ with $-R^2<t_k\le0$.
Their cylinders have centres in $Q_{3/4}$, so Step 1 gives
\[
 \operatorname*{ess\,sup}_{t_k-\ell^2<s<t_k}
    \int_{B_\ell(x_j)}|u(y,s)|^2\,dy
   \le \ell\,\alpha((x_j,t_k),\ell)^2
   \le M^2\ell^{1+2\varepsilon}.
\]
Summing over the finite spatial cover and taking the maximum over the
time intervals bounds $A_R$ by a constant depending only on $M,r_5$
and the fixed radii. Thus every quantitative hypothesis of
Proposition~\ref{prop:past-bootstrap} is controlled by the past data.

The first conclusion of that proposition upgrades the velocity to
$\mathbf1_{E_{9/16}}u\in\mathcal M^{3,25}$ before its final
pressure-gradient estimate and H\"older step are applied.
It supplies a single function $\bar v$, of exponent
$\gamma_0=\min\{2-5/q,1/5\}$, agreeing a.e.\ with $u$ on $Q_{1/2}$
and satisfying $[\bar v]_{C^{\gamma_0,\gamma_0/2}}\le C(q)$.
The causal source splitting in that proposition requires no positive-time
Morrey bound. To control its value as well as its oscillation, put
$E=Q_{1/2}$ and average over this same past cylinder:
\[
 \sup_{z\in\overline E}|\bar v(z)|
 \le |E|^{-1/3}\|u\|_{L^3(E)}
       +[\bar v]_{C^{\gamma_0,\gamma_0/2}}
           \operatorname{diam}_{\dpar}(\overline E)^{\gamma_0}
 \le C\varepsilon_0^{1/3}+C(q).
\]
Taking $w=\bar v|_{\overline{Q_{1/2}}}$ gives the required bound
$\|w\|_{C^{\gamma_0,\gamma_0/2}(\overline{Q_{1/2}})}\le C_4(q)$.
Every point of the open interior has an open neighbourhood within
$Q_{1/2}$ and is therefore regular. The argument makes no assertion of
open-neighbourhood regularity on its top face.
\end{proof}
\section{From the gradient criterion to the partial regularity theorem}
\label{sec:covering}

This section proves Theorem~\ref{thm:C} assuming Theorem~\ref{thm:B}.

\begin{lemma}[$\Phaus^{1}$ finite implies Lebesgue null]\label{lem:P1-implies-null}
Let $E\subset\R^{3}\times\R$ with $\Phaus^{1}(E)<\infty$. Then $E$ has Lebesgue
measure zero in $\R^{4}$.
\end{lemma}

\begin{proof}
Let $M=\Phaus^{1}(E)$ and let $\delta\in(0,1]$. Since
$\Phaus^{1}_{\delta}(E)\le\Phaus^{1}(E)=M$, there are $z_{j}$ and
$r_{j}\in(0,\delta]$ with $E\subset\bigcup_{j}\mathfrak{B}_{r_{j}}(z_{j})$ and
$\sum_{j}r_{j}\le M+1$. By \eqref{eq:parabolic-ball},
$\abs{\mathfrak{B}_{r}(z)}=\abs{\Ball{r}{x}}\cdot2r^{2}=\tfrac{8\pi}{3}r^{5}$,
so
\[
  \abs{E}\le\sum_{j}\abs{\mathfrak{B}_{r_{j}}(z_{j})}
  =\frac{8\pi}{3}\sum_{j}r_{j}^{5}
  \le\frac{8\pi}{3}\,\delta^{4}\sum_{j}r_{j}
  \le\frac{8\pi}{3}\,\delta^{4}(M+1).
\]
Letting $\delta\to0^{+}$ gives $|E|=0$.
Equivalently, the same cover gives zero five-dimensional parabolic
Hausdorff measure, which dominates Lebesgue measure up to an absolute
normalization. This comparison gives an alternative proof of the lemma.
\end{proof}

\begin{proof}[Proof of Theorem~\ref{thm:C}]
Let $\varepsilon_{1}=\varepsilon_{1}(q)$ be the constant of
Theorem~\ref{thm:B} and set $\varepsilon_{\mathrm{sing}}=\varepsilon_{1}^{2}>0$.

\emph{Step 1: the defect inequality at singular points.} Let $z\in S$. By
Theorem~\ref{thm:B} (in contrapositive form), \eqref{eq:thmB-hyp} must fail, so
\begin{equation}\label{eq:defect}
  \limsup_{r\to0^{+}}\;\frac1r\iint_{\Cyl{r}{z}}\abs{\nabla u}^{2}\,dz'
  \;\ge\;\varepsilon_{\mathrm{sing}}.
\end{equation}
Consequently, for every $\tau>0$ there exists $r\in(0,\tau)$ with
$\Cyl{r}{z}\subset\mathcal{O}$ and
\begin{equation}\label{eq:defect-r}
  \iint_{\Cyl{r}{z}}\abs{\nabla u}^{2}\,dz'\;>\;\tfrac12\varepsilon_{\mathrm{sing}}\,r .
\end{equation}

\emph{Step 2: reduction to relatively compact sets.}
Equivalently, one can cover the domain by countably many relatively
compact local boxes, run the following argument on each closed smaller
box, and use countable subadditivity. This uses only second countability
and local finiteness of the Dirichlet energy. We give the compact-exhaustion
version for comparison.
$\mathcal{O}$ is open in $\R^{4}$, so
there are compact sets $K_{1}\subset K_{2}\subset\cdots$ with
$\bigcup_{n}K_{n}=\mathcal{O}$ and $K_{n}\subset\operatorname{int}K_{n+1}$; for
instance $K_{n}=\{z\in\mathcal{O}:\abs{z}\le n,\ \dist(z,\R^{4}\setminus\mathcal{O})\ge1/n\}$
when $\mathcal{O}\ne\R^{4}$, and $K_{n}=\{\abs{z}\le n\}$ otherwise. Since
$\Phaus^{1}$ is an outer measure
(External Input~\ref{ext:hausdorff}) and $S=\bigcup_{n}(S\cap K_{n})$, it
suffices to prove $\Phaus^{1}(S\cap K_{n})=0$ for each $n$. Fix $n$, write
$K=K_{n}$ and $V_{0}=\operatorname{int}K_{n+1}$, so that $K\subset V_{0}$,
$\overline{V_{0}}\subset K_{n+2}\subset\mathcal{O}$ is compact, and by (S1)
\begin{equation}\label{eq:finite-dirichlet}
  \mathcal{E}_{0}:=\iint_{V_{0}}\abs{\nabla u}^{2}\,dz'<\infty .
\end{equation}
By Lemma~\ref{lem:S-closed}, $S$ is relatively closed in $\mathcal{O}$, so
$S\cap K$ is compact.

\emph{Step 3: the covering.} Let $V$ be any open set with
$S\cap K\subset V\subset V_{0}$, and let $\delta>0$. For each $z\in S\cap K$
choose, by Step 1, a radius $r_{z}\in(0,\delta)$ such that
$\mathfrak{B}_{r_{z}}(z)\subset V$ and \eqref{eq:defect-r} holds with $r=r_{z}$;
this is possible because $V$ is open and $z\in V$, so
$\mathfrak{B}_{r}(z)\subset V$ for all small $r$. The family
$\mathcal{F}=\{\mathfrak{B}_{r_{z}}(z):z\in S\cap K\}$ consists of balls of the
metric space $(\R^{4},\dpar)$ with radii bounded by $\delta$. That space is
separable: $\Q^{4}$ is a countable $\dpar$-dense subset, because
$\dpar\le\abs{x-x'}+\abs{t-t'}^{1/2}$ and both terms can be made small by a
rational approximation. By
External Input~\ref{ext:vitali} there is a countable subfamily
$\{\mathfrak{B}_{r_{j}}(z_{j})\}_{j\in J}$, $J\subset\N$, $r_{j}=r_{z_{j}}$, which is pairwise
disjoint and such that every member of $\mathcal{F}$ is contained in
$\mathfrak{B}_{5r_{j}}(z_{j})$ for some $j$. Since each $z\in S\cap K$ lies in
$\mathfrak{B}_{r_{z}}(z)\in\mathcal{F}$,
\begin{equation}\label{eq:covering}
  S\cap K\subset\bigcup_{j\in J}\mathfrak{B}_{5r_{j}}(z_{j}),
  \qquad 5r_{j}<5\delta .
\end{equation}

\emph{Step 4: the measure estimate.} By Lemma~\ref{lem:parabolic-metric},
$\Cyl{r_{j}}{z_{j}}\subset\mathfrak{B}_{r_{j}}(z_{j})\subset V$, and the
cylinders $\Cyl{r_{j}}{z_{j}}$ are pairwise disjoint because the balls
$\mathfrak{B}_{r_{j}}(z_{j})$ are. Hence, by \eqref{eq:defect-r},
\begin{equation}\label{eq:sum-radii}
  \sum_{j}r_{j}
  \;\le\;\frac{2}{\varepsilon_{\mathrm{sing}}}\sum_{j}\iint_{\Cyl{r_{j}}{z_{j}}}\abs{\nabla u}^{2}
  \;\le\;\frac{2}{\varepsilon_{\mathrm{sing}}}\iint_{V}\abs{\nabla u}^{2} ,
\end{equation}
and therefore, by \eqref{eq:covering} and Definition~\ref{def:parabolic-hausdorff},
\begin{equation}\label{eq:P1-delta}
  \Phaus^{1}_{5\delta}(S\cap K)\;\le\;\sum_{j}5r_{j}
  \;\le\;\frac{10}{\varepsilon_{\mathrm{sing}}}\iint_{V}\abs{\nabla u}^{2} .
\end{equation}

\emph{Step 5: $S\cap K$ is Lebesgue null.} Apply \eqref{eq:P1-delta} with
$V=V_{0}$: for every $\delta>0$,
$\Phaus^{1}_{5\delta}(S\cap K)\le10\varepsilon_{\mathrm{sing}}^{-1}\mathcal{E}_{0}$, whence
$\Phaus^{1}(S\cap K)=\lim_{\delta\to0}\Phaus^{1}_{5\delta}(S\cap K)
 \le10\varepsilon_{\mathrm{sing}}^{-1}\mathcal{E}_{0}<\infty$ by
\eqref{eq:finite-dirichlet}. By Lemma~\ref{lem:P1-implies-null},
$\abs{S\cap K}=0$.

\emph{Step 6: conclusion.} Let $\varepsilon>0$. Since
$\abs{\nabla u}^{2}\in L^{1}(V_{0})$ by \eqref{eq:finite-dirichlet}, absolute
continuity of the integral gives $\tau>0$ such that
$\iint_{W}\abs{\nabla u}^{2}<\varepsilon$ for every measurable $W\subset V_{0}$
with $\abs{W}<\tau$. By Step 5 and outer regularity of Lebesgue measure there is
an open $V'\supset S\cap K$ with $\abs{V'}<\tau$; replacing $V'$ by
$V=V'\cap V_{0}$, which is open and still contains $S\cap K$, we get
$\abs{V}<\tau$ and hence $\iint_{V}\abs{\nabla u}^{2}<\varepsilon$. Now
\eqref{eq:P1-delta} gives $\Phaus^{1}_{5\delta}(S\cap K)\le10\varepsilon_{\mathrm{sing}}^{-1}\varepsilon$
for every $\delta>0$, so $\Phaus^{1}(S\cap K)\le10\varepsilon_{\mathrm{sing}}^{-1}\varepsilon$.
As $\varepsilon>0$ was arbitrary, $\Phaus^{1}(S\cap K)=0$. By Step 2,
$\Phaus^{1}(S)\le\sum_{n}\Phaus^{1}(S\cap K_{n})=0$.
\end{proof}

\begin{remark}[Why parabolic balls rather than cylinders in the covering]
\label{rem:balls-vs-cylinders}
The Vitali lemma with the $5r$ enlargement is a statement about balls of a
metric space, and the parabolic cylinders $\Cyl{r}{z}$ are not balls of
$\dpar$. Two intersecting backward cylinders need not be contained in
one another's enlargements about the original top centres: a point later
than one top time remains outside every cylinder with that top time.
We therefore run the covering with the genuine $\dpar$-balls
$\mathfrak{B}_{r}(z)$, and recover the statement about cylinders through
Lemma~\ref{lem:P1-cylinders}. The defect inequality \eqref{eq:defect-r} is still
stated on cylinders, which is where the local energy inequality lives, and the
inclusion $\Cyl{r}{z}\subset\mathfrak{B}_{r}(z)$ of
Lemma~\ref{lem:parabolic-metric} is what links the two. An alternative used in
\cite[\S6.2, pp.~96--97]{Tsai2018} uses the shifted cylinders
$Q^*(z,r)=B_r(x)\times(t-7r^2/8,t+r^2/8)$, which contain both $z$ and
$\Cyl{r/2}{z}$.
\end{remark}

\appendix

\section{External inputs}\label{sec:external}

The analytic inputs and comparison results are collected below with their
precise scope. Several inputs are proved internally in the formal development;
their designation here does not mean that they are additional axioms.
The Campanato comparison and Aubin--Lions compactness entries are not used
by the direct A/B/C proof route. Constants named in this
appendix are absolute unless a dependence is displayed.

\begin{externalinput}[Calder\'on--Zygmund bounds for the Riesz transforms]\label{ext:CZ}
For $i\in\{1,2,3\}$ let $R_{i}$ denote the $i$-th Riesz transform on $\R^{3}$,
that is, for the Fourier convention with kernel $e^{-\mathrm{i}x\cdot\xi}$,
the multiplier operator with symbol $-\mathrm{i}\xi_{i}/\abs{\xi}$;
equivalently $R_{i}=-\partial_{i}(-\Delta)^{-1/2}$, so that
$R_{i}R_{j}=\partial_{i}\partial_{j}(-\Delta)^{-1}$. For every $s\in(1,\infty)$
there is a constant $C_{11}=C_{11}(s)$ such that
\[
  \norm{R_{i}R_{j}g}_{L^{s}(\R^{3})}\le C_{11}(s)\norm{g}_{L^{s}(\R^{3})}
  \qquad\text{for all }g\in L^{s}(\R^{3})\text{ and all }i,j .
\]
For the pressure itself we write $C_{11}=C_{11}(3/2)$; the weak-gradient
route also uses the fixed exponent $6/5$. These fixed-exponent constants
are absolute. The operators are the bounded $L^s$ extensions of the
second Riesz transforms from a dense class, not absolutely convergent
global integrals of the Hessian kernel. Literal kernel formulas are used
only off the source support, or with their proper singular-integral
interpretation. Interpolation is first performed on $L^2\cap L^s$ and
then passed to the completed $L^s$ operator.
Reference: \cite[Ch.~II--III]{Stein1970}.
\end{externalinput}

\begin{externalinput}[Hardy--Littlewood--Sobolev for the Riesz potential of order $1$]
\label{ext:riesz}
Let $N(y)=-\tfrac1{4\pi\abs{y}}$, so that $\abs{\nabla N(y)}=\tfrac1{4\pi\abs{y}^{2}}$.
There is an absolute finite constant $C'_{11}$ such that, for every
$g\in L^{5/2}(\R^3)$ and every $j\in\{1,2,3\}$, the literal convolution
$\partial_jN*g$ is absolutely convergent for almost every $x$ and
\[
  \norm{\partial_jN*g}_{L^{15}(\R^3)}
      \le C'_{11}\norm{g}_{L^{5/2}(\R^3)} .
\]
This is the specialization $1/15=1/(5/2)-1/3$ needed here; the full
Hardy--Littlewood--Sobolev exponent family is not part of this input.
No pointwise measurability or compact-support assumption is added to the
$L^{5/2}$ datum.
Reference: \cite[Ch.~V, Thm.~1]{Stein1970}.
\end{externalinput}

\begin{externalinput}[Newtonian representation of a compactly supported solution]
\label{ext:newtonian}
Let $N(y)=-\tfrac1{4\pi\abs{y}}$, so that $\Delta N=\delta_{0}$ in
$\mathcal{D}'(\R^{3})$. Let $v\in L^{3/2}(\R^{3})$ have compact support and
suppose that in $\mathcal{D}'(\R^{3})$
\[
  \Delta v=\partial_{i}\partial_{j}G_{ij}+\partial_{j}H_{j}+K
\]
where $G\in L^{3/2}(\R^{3};\R^{3\times3})$, $H\in L^{m}(\R^{3};\R^{3})$ for some
$m\in(1,3)$, and $K\in L^{m'}(\R^{3})$ for some $m'\in(1,3/2)$, all with compact
support. Then
\[
  v=-R_{i}R_{j}G_{ij}+\partial_{j}N*H_{j}+N*K
  \qquad\text{almost everywhere in }\R^{3}.
\]
Reference: \cite[Ch.~2 and Ch.~4]{GilbargTrudinger2001} for the Newtonian
potential and \cite[Ch.~III, V]{Stein1970} for the identification
$\partial_{i}\partial_{j}N*G_{ij}=-R_{i}R_{j}G_{ij}$; the uniqueness statement
is Liouville's theorem applied to the difference, which is harmonic on $\R^{3}$
and tends to $0$ in the $L^{3/2}$-average sense at infinity.
\end{externalinput}

\begin{externalinput}[Sobolev embeddings]\label{ext:sobolev-ball}\leavevmode
\begin{enumerate}
\item[(i)] There is an absolute constant such that
  $\norm{g}_{L^{6}(\R^{3})}\le C\norm{\nabla g}_{L^{2}(\R^{3})}$ for all
  $g\in H^{1}(\R^{3})$; consequently
  $\norm{g}_{L^{6}(\R^{3})}\le C\norm{g}_{H^{1}(\R^{3})}$.
\item[(ii)] There is an absolute constant $C$ such that for every Euclidean
  ball $B_r(x)$, $r>0$, and every $g\in H^1(B_r(x))$,
  \[
    \norm{g}_{L^6(B_r(x))}
      \le C\bigl(\norm{\nabla g}_{L^2(B_r(x))}
                       +r^{-1}\norm{g}_{L^2(B_r(x))}\bigr).
  \]
  In particular the unscaled $H^1$ estimate has an absolute constant on $B_1$.
  No general Lipschitz-domain extension theorem is required here.
\item[(iii)] There is an absolute constant with
  $\norm{g}_{L^{\infty}(\R^{3})}\le C\norm{g}_{H^{2}(\R^{3})}$ for all
  $g\in H^{2}(\R^{3})$.
\item[(iv)] Let $J$ be an interval and $U=B_r(x)$, $r>0$. If
  $g\in L^\infty(J;L^2(U))\cap L^2(J;H^1(U))$, then
  $g\in L^{10/3}(U\times J)$ and, with one absolute constant $C$,
  \[
  \norm{g}_{L^{10/3}(U\times J)}^{10/3}
   \le C\bigl(\norm{g}_{L^\infty(J;L^2(U))}^{4/3}
                  \norm{\nabla g}_{L^2(U\times J)}^2
       +r^{-2}\norm{g}_{L^\infty(J;L^2(U))}^{10/3}|J|\bigr).
  \]
  For unbounded $J$, the membership conclusion is asserted separately;
  an infinite right side is not a proof of membership. The displayed
  bound is \eqref{eq:interp-q103} integrated in time with the slice
  $L^2$ mass bounded by the essential supremum. No embedding on an
  arbitrary bounded open set, or sharper public actual-mass estimate,
  is asserted here.
\end{enumerate}
References: \cite[Ch.~7]{GilbargTrudinger2001}, \cite[Ch.~V]{Stein1970}.
\end{externalinput}

\begin{externalinput}[Poincar\'e inequalities on the unit ball]\label{ext:poincare}\leavevmode
\begin{enumerate}
\item[(i)] There is an absolute constant with
  $\norm{g-\spavg{g}{B_{1}}}_{L^{1}(B_{1})}\le C\norm{\nabla g}_{L^{1}(B_{1})}$
  for all $g\in W^{1,1}(B_{1})$.
\item[(ii)] There is an absolute constant with
  $\norm{g-\spavg{g}{B_{1}}}_{L^{6}(B_{1})}\le C\norm{\nabla g}_{L^{2}(B_{1})}$
  for all $g\in H^{1}(B_{1})$.
\item[(iii)] \emph{Additional classical statement.} There is an absolute constant with
  $\norm{g-\spavg g{B_1}}_{L^{3/2}(B_1)}
          \le C\norm{\nabla g}_{L^1(B_1)}$
  for all $g\in W^{1,1}(B_1)$. Rescaling gives the same coefficient on
  any Euclidean ball, since both sides have scaling degree $2$.
  This arbitrary $W^{1,1}$ version is recorded as classical paper material,
  not as an available formal export; the formal scalar estimate presently
  treats $C^1$ data.
\end{enumerate}
Reference: \cite[Ch.~7]{GilbargTrudinger2001}.
\end{externalinput}

The formalized Caccioppoli application uses a separate squared-velocity
estimate: for $0<r<R$, $B_R(x_0)\subset U$ with $U$ open, and
$u\in H^1(U;\R^3)$ with its weak gradient,
\[
 \bigl\||u|^2-\spavg{|u|^2}{B_r(x_0)}\bigr\|_{L^{3/2}(B_r(x_0))}
 \le C\norm u_{L^2(B_r(x_0))}\norm{\nabla u}_{L^2(B_r(x_0))}.
\]
The constant is absolute. Approximation on nested balls proves this
particular weak-slice estimate; it does not by itself export clause (iii)
for arbitrary $W^{1,1}$ functions.

\begin{externalinput}[Interior estimates for harmonic functions]\label{ext:harmonic}
Let $h$ be harmonic on an open $U\subset\R^{3}$.
\begin{enumerate}
\item[(i)] \emph{(Mean value property.)} If $\overline{\Ball{s}{x}}\subset U$
  then $h(x)=\abs{\Ball{s}{x}}^{-1}\int_{\Ball{s}{x}}h$.
\item[(ii)] \emph{(Gradient estimate.)} If $\overline{\Ball{s}{x}}\subset U$
  then $\abs{\nabla h(x)}\le\tfrac3s\sup_{\Ball{s}{x}}\abs{h}$.
\end{enumerate}
Reference: \cite[Ch.~2]{GilbargTrudinger2001}.
\end{externalinput}

\begin{externalinput}[Vitali covering lemma, $5r$ form]\label{ext:vitali}
Let $(X,d)$ be a \emph{separable} metric space and let $\mathcal{F}$ be a family
of open balls $\mathfrak{B}_{r_{a}}(z_{a})$, $a\in\mathcal{I}$, with
$\sup_{a}r_{a}<\infty$. Then there is a countable set
$\mathcal{J}\subset\mathcal{I}$ such that the balls
$\{\mathfrak{B}_{r_{j}}(z_{j})\}_{j\in\mathcal{J}}$ are pairwise disjoint and
\[
  \bigcup_{a\in\mathcal{I}}\mathfrak{B}_{r_{a}}(z_{a})
  \subset\bigcup_{j\in\mathcal{J}}\mathfrak{B}_{5r_{j}}(z_{j}) .
\]
In particular every $\mathfrak{B}_{r_{a}}(z_{a})$ is contained in some
$\mathfrak{B}_{5r_{j}}(z_{j})$. Separability is needed for the subfamily to be
\emph{countable}: without it the standard Zorn's-lemma argument produces only a
maximal disjoint subfamily, which may be uncountable.
The open-ball version follows directly from the usual $5r$ argument. Choose
$R>\sup_a r_a$ and successively choose a maximal disjoint family among
balls of radii in $(2^{-n-1}R,2^{-n}R]$ that avoid the earlier choices.
Every unchosen ball meets a chosen ball of radius at least half its own,
and is therefore contained in that ball's fivefold enlargement. The chosen
balls are countable because each contains a distinct point of a fixed
countable dense set. This also proves the individual containment assertion.
For the covering argument in its standard form see \cite[Thm.~2.1]{Mattila1995}.
\end{externalinput}

\begin{externalinput}[Hausdorff measures of a metric space]\label{ext:hausdorff}
Let $(X,d)$ be a metric space and $\alpha\ge0$. The set functions
$\Phaus^{\alpha}_{\delta}$ and $\Phaus^{\alpha}$ of
Definition~\ref{def:parabolic-hausdorff} are outer measures on $X$; in
particular $\Phaus^{\alpha}$ is countably subadditive and monotone, and
$\Phaus^{\alpha}(E)=\lim_{\delta\to0^{+}}\Phaus^{\alpha}_{\delta}(E)
 =\sup_{\delta>0}\Phaus^{\alpha}_{\delta}(E)$, the family
$\delta\mapsto\Phaus^{\alpha}_{\delta}(E)$ being nonincreasing in $\delta$.
Reference: \cite[Ch.~4]{Mattila1995}.
\end{externalinput}

\begin{externalinput}[Parabolic Campanato characterization of H\"older continuity]
\label{ext:campanato}
Let $\alpha\in(0,1)$, $R>0$, $z_{0}\in\R^{3}\times\R$, and
$w\in L^{3}(\Cyl{2R}{z_{0}};\R^{m})$. Suppose
\[
  \Lambda^{3}\;:=\;\sup\Bigl\{\frac1{r^{5+3\alpha}}
    \iint_{\Cyl{r}{z}}\bigl|w-\stavg{w}{\Cyl{r}{z}}\bigr|^{3}dz'
    \ :\ z\in \Cyl{R}{z_{0}},\ 0<r\le R\Bigr\}\;<\;\infty .
\]
Then $w$ agrees almost everywhere on $\Cyl{R}{z_{0}}$ with a function
$\bar w\in C^{\alpha,\alpha/2}(\overline{\Cyl{R}{z_{0}}};\R^{m})$, and there is a
constant $C(\alpha)$ with
\[
  [\bar w]_{C^{\alpha,\alpha/2}(\Cyl{R}{z_{0}})}\le C(\alpha)\Lambda,
  \qquad
  \norm{\bar w}_{L^{\infty}(\Cyl{R}{z_{0}})}
   \le C(\alpha)\Bigl(\Lambda R^{\alpha}
     +R^{-5/3}\norm{w}_{L^{3}(\Cyl{2R}{z_{0}})}\Bigr).
\]
The mean $\stavg{w}{\Cyl{r}{z}}$ is the genuine space-time average
\eqref{eq:spacetime-average}. Note that the normalization is
$r^{-5-3\alpha}$, so that in terms of the excess \eqref{eq:excess} the
hypothesis reads $\widetilde C(z,r)^{1/3}\le\Lambda\,r^{1+\alpha}$.
Reference: \cite[\S4.2]{LadyzhenskayaSolonnikovUraltseva1968},
\cite[\S III.1]{Giaquinta1983} for the elliptic model case, and
\cite[after (6.24)]{Tsai2018} for the form used here. This input is recorded for
comparison with the blow-up route only: Section~\ref{sec:core} uses instead the
self-contained Lemma~\ref{lem:campanato}, and the exponent it delivers is
$\gamma_{0}=\min\{2-5/q,1/5\}$ (Remark~\ref{rem:exponent}).
\end{externalinput}

\begin{externalinput}[Heat kernel]\label{ext:heat}
Let $G(x,\tau)=(4\pi\tau)^{-3/2}\exp\bigl(-\abs{x}^{2}/(4\tau)\bigr)$ for
$x\in\R^{3}$, $\tau>0$. Then $G\in C^{\infty}(\R^{3}\times(0,\infty))$, $G>0$,
$\int_{\R^{3}}G(x,\tau)\,dx=1$ for every $\tau>0$, and
$\partial_{\tau}G=\Delta G$ on $\R^{3}\times(0,\infty)$. Consequently, for
$r>0$, the function $\psi_{r}(x,t)=r^{2}G(x,r^{2}-t)$ is smooth and positive on
$\R^{3}\times(-\infty,r^{2})$ and satisfies
$\partial_{t}\psi_{r}+\Delta\psi_{r}=0$ there.
Reference: \cite[Ch.~2]{Tsai2018}.
\end{externalinput}

\begin{externalinput}[Heat kernel: size and derivative bounds]\label{ext:heat-kernel}
Let
\[
W_{+}(x,t)=\mathbf{1}_{t>0}(4\pi t)^{-3/2}e^{-\abs{x}^{2}/(4t)}
.
\] Then
$W_{+}\in L^{1}_{loc}(\R^{3}\times\R)$, $\supp W_{+}\subset\{t\ge0\}$,
$(\partial_{t}-\Delta)W_{+}=\delta_{(0,0)}$ in $\mathcal{D}'(\R^{3}\times\R)$,
and convolution with a compactly supported distribution is well-defined.
Causal uniqueness, if used, must be restricted to tempered distributions:
a tempered $d$ supported in $\{t\ge c\}$ and satisfying
$(\partial_t-\Delta)d=0$ is zero. No such assertion holds for arbitrary
distributions, and half-space support does not define a universal
convolution algebra. The proof below uses the adjoint-test identity
in Lemma~\ref{lem:local-equation}, not this uniqueness assertion.
There is an absolute constant $C$
such that, with
$\varrho=\max\{|x|,|t|^{1/2}\}$, for $t>0$,
\[
  \abs{W_{+}}\le C\varrho^{-3},\quad
  \abs{\nabla_{x}W_{+}}\le C\varrho^{-4},\quad
  \abs{\partial_{t}W_{+}}\le C\varrho^{-5},
\]
and, for every Fourier multiplier $\varsigma(D)$ with $\varsigma$ smooth on
$\R^{3}\setminus\{0\}$ and homogeneous of degree $1$,
\[
  \abs{\varsigma(D)W_{+}}\le C_{\varsigma}\varrho^{-4},\quad
  \abs{\nabla_{x}\varsigma(D)W_{+}}\le C_{\varsigma}\varrho^{-5},\quad
  \abs{\partial_{t}\varsigma(D)W_{+}}\le C_{\varsigma}\varrho^{-6}.
\]
The causal kernels vanish for $t<0$. Their size estimates hold almost
everywhere on spacetime, but no time derivative at $t=0$ is asserted for
the general multiplier kernels.
Reference: \cite[\S IV.11]{LadyzhenskayaSolonnikovUraltseva1968},
\cite[Ch.~5]{LemarieRieusset2016}.
\end{externalinput}

\begin{externalinput}[Aubin--Lions compactness]\label{ext:aubin-lions}
Let $X\subset\subset Y\subset Z$ be Banach spaces with $X$ reflexive, the
embedding $X\hookrightarrow Y$ compact and $Y\hookrightarrow Z$ continuous. Let
$J$ be a bounded interval, $1<s_{0}\le\infty$, $1<s_{1}<\infty$, and let
$(g_{k})$ be a sequence with
$\sup_{k}\bigl(\norm{g_{k}}_{L^{s_{0}}(J;X)}+\norm{\partial_{t}g_{k}}_{L^{s_{1}}(J;Z)}\bigr)<\infty$.
Then $(g_{k})$ is relatively compact in $L^{s_{0}}(J;Y)$ (and in $C(\overline{J};Y)$
if $s_{0}=\infty$ and $s_{1}>1$). For comparison with a compactness proof one may take
$X=H^{1}(B)$, $Y=L^{2}(B)$, $Z=Z_B$ of \eqref{eq:Z-space},
$s_0=2$, $s_1=3/2$. This input is not consumed in the direct proof here.
Reference: \cite[Corollary 4]{Simon1987}, including the time-continuity endpoint.
\end{externalinput}

\begin{externalinput}[Pettis measurability theorem]\label{ext:pettis}
Let $X$ be a separable Banach space, $J\subset\R$ an interval, and
$w:J\to X$ such that $t\mapsto\langle w(t),\varphi\rangle$ is Lebesgue
measurable for every $\varphi\in X^{*}$. Then $w$ is strongly (Bochner)
measurable, i.e.\ an a.e.\ limit of countably valued measurable functions.
In particular, if $X$ is a separable
Hilbert space or the dual of a separable Hilbert space, weak measurability
implies strong measurability.
Reference: \cite[Ch.~II]{DiestelUhl1977}.
\end{externalinput}

\begin{externalinput}[Vector-valued du Bois-Reymond lemma]\label{ext:dubois-reymond}
Let $X$ be a Banach space, $J\subset\R$ an open interval and
$w\in L^{1}_{loc}(J;X)$ with
$\int_{J}w(t)\,\theta'(t)\,dt=0$ in $X$ for every $\theta\in C_{c}^{\infty}(J)$.
Then there is $\xi\in X$ with $w(t)=\xi$ for almost every $t\in J$.
Reference: \cite[Ch.~III, Lemma 1.1]{Temam1977}.
\end{externalinput}

\begin{externalinput}[Mollification and density]\label{ext:mollify}
Let $\zeta\in C_{c}^{\infty}(B_{1})$ with $\zeta\ge0$ and
$\int_{\R^{3}}\zeta=1$, and $\zeta_{\varepsilon}(y)=\varepsilon^{-3}\zeta(y/\varepsilon)$.
For $g\in L^{s}(\R^{3})$, $1\le s<\infty$, one has
$g*\zeta_{\varepsilon}\in C^{\infty}(\R^{3})$,
$\norm{g*\zeta_{\varepsilon}}_{L^{s}}\le\norm{g}_{L^{s}}$ and
$g*\zeta_{\varepsilon}\to g$ in $L^{s}(\R^{3})$; the analogous statements hold in
one time variable. Moreover $C_{c}^{\infty}(U)$ is dense in $H^{2}_{0}(U)$ by
definition of the latter, and $C_{c}^{\infty}(U)$ is dense in $L^{s}(U)$ for
$1\le s<\infty$ and $U\subset\R^{3}$ open.
Reference: \cite[Ch.~7]{GilbargTrudinger2001}.
\end{externalinput}

\bibliographystyle{amsplain}
\bibliography{refs}

\end{document}
